% Solutions aqueuses — PCT 3ème — Cours signé OnBuch+
% Programme officiel MINESEC. Compiler avec : tectonic N06-solutions-aqueuses.tex
\newcommand{\DOCMATIERE}{PCT}
\newcommand{\DOCNIVEAU}{3ème}
\newcommand{\DOCMODULE}{Module 1 : Chimie}
\newcommand{\DOCLECON}{N06}
\newcommand{\DOCTITRE}{Solutions aqueuses}
% OnBuch+ — préambule « cours signés » ENRICHI (compilé avec tectonic / XeLaTeX)
% Système visuel « Pop » d'OnBuch+ : fond crème (jamais blanc), bordures encre
% épaisses, ombres « dures » décalées, titres Archivo Black en capitales.
% Version MATHÉMATIQUES : graphes (pgfplots/tikz), boîtes pédagogiques
% (définition, propriété, méthode, attention, le savais-tu, exercice résolu,
% exercice, TP, bilan), page de garde et signature OnBuch+ ancrée en bas.
%
% Macros attendues dans main.tex AVANT \input{preamble} :
%   \DOCMATIERE  \DOCNIVEAU  \DOCMODULE  \DOCLECON (numéro)  \DOCTITRE
\documentclass[11pt]{article}

\usepackage{fontspec}
\usepackage[french]{babel}
\usepackage[a4paper,margin=2cm,top=3.4cm,bottom=2.8cm,headheight=1.4cm,footskip=1.3cm]{geometry}
\usepackage{amsmath,amssymb,amsfonts}
\usepackage{mathtools} % \xmapsto, \coloneqq... souvent utilisés par les modèles
\usepackage{cancel} % simplifications barrées (bilans de forces)
% \abs{z} (module, valeur absolue) et \norm{u} (norme), utilisés par les modèles
\providecommand{\abs}{}\let\abs\relax\DeclarePairedDelimiter{\abs}{\lvert}{\rvert}
\providecommand{\norm}{}\let\norm\relax\DeclarePairedDelimiter{\norm}{\lVert}{\rVert}
\usepackage[table]{xcolor} % [table] : fournit aussi \rowcolors (alternance de lignes)
\usepackage{pagecolor}
\usepackage{tikz}
\usetikzlibrary{arrows.meta,positioning,calc,shapes.geometric,shapes.misc,shapes.arrows,decorations.pathmorphing,decorations.pathreplacing,decorations.markings,patterns,shadows,mindmap,trees,fit,backgrounds,matrix,intersections,angles,quotes}
\usepackage{pgfplots}
\usepackage[version=4]{mhchem} % équations nucléaires \ce{^{235}_{92}U}
\usepackage[european,siunitx]{circuitikz} % schémas électriques
\pgfplotsset{compat=1.18}
\usepgfplotslibrary{fillbetween,groupplots} % aires entre courbes (calcul intégral)
\usepackage{enumitem}
\usepackage{fancyhdr}
% [most] charge toutes les bibliothèques tcolorbox (listings, minted, raster,
% documentation...) et gonfle inutilement la mémoire TeX sur un document long
% (~300 boîtes) : on ne charge que ce qui est réellement utilisé.
\usepackage[skins,breakable]{tcolorbox}
\usepackage{graphicx}
\usepackage{colortbl}
\usepackage{booktabs}
\usepackage{tabularx}
\usepackage{multirow}
\usepackage{diagbox} % case d'en-tête diagonale des tableaux à double entrée
% Colonnes extensibles centrée / gauche / droite (C, L, R), souvent utilisées par les modèles
\newcolumntype{C}{>{\centering\arraybackslash}X}
\newcolumntype{L}{>{\raggedright\arraybackslash}X}
\newcolumntype{R}{>{\raggedleft\arraybackslash}X}
\usepackage{array}
\usepackage{multicol}
\usepackage{siunitx}
% parse-numbers=false : accepte aussi \SI{2,5 \times 10^{-3}}{...}
\sisetup{locale=FR,output-decimal-marker={,},group-minimum-digits=4,parse-numbers=false}
\DeclareSIUnit\ohm{\ensuremath{\symup{\Omega}}} % U+2126 absent de la police
\DeclareSIUnit\division{div} % réglages d'oscilloscope (ms/div, V/div)
\DeclareSIUnit\tr{tr} % tours (vitesse de rotation, ex. \SI{1500}{\tr\per\minute})
\DeclareSIUnit\molar{M} % concentration molaire (ex. \SI{3}{\molar}, \SI{1}{\micro\molar})
\DeclareSIUnit\an{an} % année (le modèle confond parfois avec \year, un compteur TeX) — ex. \SI{5}{\kilo\meter\per\an}
\DeclareSIUnit\FCFA{FCFA} % franc CFA (ex. \SI{500}{\FCFA\per\kilo\gram})
\DeclareSIUnit\degreeBrix{\textdegree Brix} % teneur en sucre d'un sirop/jus (réfractométrie)
\DeclareSIUnit\month{mois} % le modèle utilise parfois \month, un compteur TeX, comme unité de durée
\DeclareSIUnit\microliter{\micro\liter} % le modèle écrit parfois \microliter en un seul token
\DeclareSIUnit\million{millions} % dénombrement (ex. \SI{15}{\million\per\mL} spermatozoïdes)
\usepackage{qrcode}
% \degree : commande fréquemment utilisée par les modèles pour un angle ou une
% température en dehors de \SI{}{} (non fournie par les paquets chargés).
\providecommand{\degree}{^{\circ}}
% \qed (fin de démonstration), utilisé par les modèles sans amsthm
\providecommand{\qed}{\hfill$\square$}
% \barre : double barre des valeurs interdites dans un tableau de variations (tkz-tab)
\providecommand{\barre}{\ensuremath{\|}}
% environnement proof (amsthm non chargé) utilisé par les modèles
\ifdefined\proof\else\newenvironment{proof}[1][Démonstration]{\par\noindent\textit{#1.}\ }{\qed\par}\fi
\usepackage{lastpage}
\usepackage{hyperref}
\hypersetup{hidelinks}

% ============================================================
% Palette « Pop » OnBuch+ — mêmes hex que la classe P (plus_ui.dart)
% ============================================================
\definecolor{poporange}{HTML}{F59321}
\definecolor{popdark}{HTML}{E07A0C}
\definecolor{popcream}{HTML}{FFF6E8}
\definecolor{popink}{HTML}{1C1714}
\definecolor{poppurple}{HTML}{7A5AE0}
\definecolor{popgreen}{HTML}{2E9E6A}
\definecolor{poppink}{HTML}{E0457B}
\definecolor{popblue}{HTML}{2F7DE1}
\definecolor{popgold}{HTML}{C98A12}
\definecolor{popmuted}{HTML}{8A7F76}
\definecolor{popline}{HTML}{EADFD2}
\definecolor{popgray}{HTML}{8A7F76}
\colorlet{popgrey}{popgray}
% Alias tolérants (le modèle nomme parfois les couleurs différemment)
\colorlet{popwhite}{white}
\colorlet{popblack}{popink}
\colorlet{popred}{poppink}
\colorlet{popyellow}{popgold}
% Teintes claires (fonds de tableaux/encadrés) — le modèle nomme parfois
% les couleurs avec un suffixe L (ex. popblueL) sans les avoir définies.
\colorlet{poporangeL}{poporange!20}
\colorlet{popdarkL}{popdark!20}
\colorlet{poppurpleL}{poppurple!20}
\colorlet{popgreenL}{popgreen!20}
\colorlet{poppinkL}{poppink!20}
\colorlet{popblueL}{popblue!20}
\colorlet{popgoldL}{popgold!20}
\colorlet{popredL}{popred!20}
\colorlet{popgrayL}{popgray!20}
% Teintes claires pour fonds de tableaux / zones
\colorlet{popblueL}{popblue!10!white}
\colorlet{poporangeL}{poporange!14!white}
\colorlet{popgreenL}{popgreen!12!white}
\colorlet{poppurpleL}{poppurple!12!white}
\colorlet{poppinkL}{poppink!12!white}
\colorlet{popgoldL}{popgold!14!white}

\pagecolor{popcream}

% ============================================================
% Typographie
% ============================================================
\defaultfontfeatures{Ligatures=TeX}
\setmainfont{PlusJakartaSans-Medium.ttf}[Path=../../../fonts/,BoldFont=PlusJakartaSans-Bold.ttf]
% Maths en sans-serif (Fira Math) avec chiffres et lettres droites de Plus
% Jakarta, pour que les formules se fondent dans le texte.
\usepackage[math-style=ISO,bold-style=ISO]{unicode-math}
\setmathfont{FiraMath-Regular.otf}[Path=../../../fonts/]
\setmathfont{PlusJakartaSans-Medium.ttf}[Path=../../../fonts/,range={up/num,up/{Latin,latin},\mathup}]
\usepackage{icomma} % virgule décimale sans espace en mode math
% Caractères mathématiques Unicode absents des polices de texte (Plus Jakarta,
% Archivo Black) : rendus par leur équivalent mathématique, en texte comme en
% formule — sinon ils s'affichent en carré vide (ex. « Arithmétique dans ℤ »).
\usepackage{newunicodechar}
\newunicodechar{ℝ}{\ensuremath{\mathbb{R}}}
\newunicodechar{ℤ}{\ensuremath{\mathbb{Z}}}
\newunicodechar{ℂ}{\ensuremath{\mathbb{C}}}
\newunicodechar{ℕ}{\ensuremath{\mathbb{N}}}
\newunicodechar{ℚ}{\ensuremath{\mathbb{Q}}}
\newunicodechar{𝔻}{\ensuremath{\mathbb{D}}}
\newunicodechar{α}{\ensuremath{\alpha}}
\newunicodechar{β}{\ensuremath{\beta}}
\newunicodechar{γ}{\ensuremath{\gamma}}
\newunicodechar{δ}{\ensuremath{\delta}}
\newunicodechar{ε}{\ensuremath{\varepsilon}}
\newunicodechar{θ}{\ensuremath{\theta}}
\newunicodechar{λ}{\ensuremath{\lambda}}
\newunicodechar{μ}{\ensuremath{\mu}}
\newunicodechar{σ}{\ensuremath{\sigma}}
\newunicodechar{τ}{\ensuremath{\tau}}
\newunicodechar{φ}{\ensuremath{\varphi}}
\newunicodechar{ψ}{\ensuremath{\psi}}
\newunicodechar{ω}{\ensuremath{\omega}}
\newunicodechar{Σ}{\ensuremath{\Sigma}}
% majuscules produites par \MakeUppercase dans les titres (α-aminés -> Α)
\newunicodechar{Α}{\ensuremath{\alpha}}
\newunicodechar{Β}{\ensuremath{\beta}}
\newunicodechar{Γ}{\ensuremath{\Gamma}}
\newunicodechar{Δ}{\ensuremath{\Delta}}
\newunicodechar{Ε}{\ensuremath{\varepsilon}}
\newunicodechar{Θ}{\ensuremath{\Theta}}
\newunicodechar{Λ}{\ensuremath{\Lambda}}
\newunicodechar{Μ}{\ensuremath{\mu}}
\newunicodechar{Τ}{\ensuremath{\tau}}
\newunicodechar{Φ}{\ensuremath{\Phi}}
\newunicodechar{Ψ}{\ensuremath{\Psi}}
\newunicodechar{Ω}{\ensuremath{\Omega}}
\newunicodechar{↦}{\ensuremath{\mapsto}}
\newunicodechar{∈}{\ensuremath{\in}}
\newunicodechar{∉}{\ensuremath{\notin}}
\newunicodechar{∧}{\ensuremath{\wedge}}
\newunicodechar{⇔}{\ensuremath{\Leftrightarrow}}
\newunicodechar{⟺}{\ensuremath{\Longleftrightarrow}}
\newunicodechar{⇒}{\ensuremath{\Rightarrow}}
\newunicodechar{⟹}{\ensuremath{\Longrightarrow}}
\newunicodechar{∘}{\ensuremath{\circ}}
\newunicodechar{∪}{\ensuremath{\cup}}
\newunicodechar{∩}{\ensuremath{\cap}}
\newunicodechar{≡}{\ensuremath{\equiv}}
\newunicodechar{⊂}{\ensuremath{\subset}}
\newunicodechar{⊥}{\ensuremath{\perp}}
\newunicodechar{∃}{\ensuremath{\exists}}
\newunicodechar{∀}{\ensuremath{\forall}}
\newunicodechar{∛}{\ensuremath{\sqrt[3]{\,}}}
\newunicodechar{∜}{\ensuremath{\sqrt[4]{\,}}}
\newunicodechar{⃗}{\ensuremath{{}^{\to}}}
\newunicodechar{ⁿ}{\ifmmode^{n}\else\textsuperscript{\ensuremath{n}}\fi}
\newunicodechar{ˣ}{\ifmmode^{x}\else\textsuperscript{\ensuremath{x}}\fi}
\newunicodechar{ᵇ}{\ifmmode^{b}\else\textsuperscript{\ensuremath{b}}\fi}
\newunicodechar{ʳ}{\ifmmode^{r}\else\textsuperscript{\ensuremath{r}}\fi}
\newunicodechar{ᵘ}{\ifmmode^{u}\else\textsuperscript{\ensuremath{u}}\fi}
\newunicodechar{ʸ}{\ifmmode^{y}\else\textsuperscript{\ensuremath{y}}\fi}
\newunicodechar{ᵃ}{\ifmmode^{a}\else\textsuperscript{\ensuremath{a}}\fi}
\newunicodechar{ᵉ}{\ifmmode^{e}\else\textsuperscript{\ensuremath{e}}\fi}
\newunicodechar{ᵏ}{\ifmmode^{k}\else\textsuperscript{\ensuremath{k}}\fi}
\newunicodechar{ᵖ}{\ifmmode^{p}\else\textsuperscript{\ensuremath{p}}\fi}
\newunicodechar{ᴺ}{\ifmmode^{N}\else\textsuperscript{\ensuremath{N}}\fi}
\newunicodechar{ᶜ}{\ifmmode^{c}\else\textsuperscript{\ensuremath{c}}\fi}
\newunicodechar{ʰ}{\ifmmode^{h}\else\textsuperscript{\ensuremath{h}}\fi}
\newunicodechar{ᵠ}{\ifmmode^{q}\else\textsuperscript{\ensuremath{q}}\fi}
\newunicodechar{ₙ}{\ifmmode_{n}\else\textsubscript{\ensuremath{n}}\fi}
\newunicodechar{ₐ}{\ifmmode_{a}\else\textsubscript{\ensuremath{a}}\fi}
\newunicodechar{ᵢ}{\ifmmode_{i}\else\textsubscript{\ensuremath{i}}\fi}
\newunicodechar{ᵣ}{\ifmmode_{r}\else\textsubscript{\ensuremath{r}}\fi}
\newunicodechar{ₖ}{\ifmmode_{k}\else\textsubscript{\ensuremath{k}}\fi}
\newunicodechar{ᵦ}{\ifmmode_{\beta}\else\textsubscript{\ensuremath{\beta}}\fi}
\newunicodechar{ₓ}{\ifmmode_{x}\else\textsubscript{\ensuremath{x}}\fi}
\newfontfamily\popdisplay{ArchivoBlack-Regular.ttf}[Path=../../../fonts/]
\newfontfamily\poplogo{SpaceGrotesk-Bold.ttf}[Path=../../../fonts/]
\newfontfamily\popsemi{PlusJakartaSans-SemiBold.ttf}[Path=../../../fonts/]
\newfontfamily\popxbold{PlusJakartaSans-ExtraBold.ttf}[Path=../../../fonts/]
\newfontfamily\popxboldls{PlusJakartaSans-ExtraBold.ttf}[Path=../../../fonts/,LetterSpace=3.0]

\setlist[enumerate]{leftmargin=1.7em,itemsep=5pt,topsep=4pt}
\setlist[itemize]{leftmargin=1.7em,itemsep=4pt,topsep=4pt,label={\color{poporange}\textbullet}}
\setlength{\parindent}{0pt}
\setlength{\parskip}{5pt}
\renewcommand{\arraystretch}{1.25}

% pgfplots : style Pop par défaut
\pgfplotsset{
  every axis/.append style={axis line style={popink,line width=1pt},tick style={popink},
    grid style={popline,line width=0.5pt},label style={font=\small},tick label style={font=\footnotesize}},
  popcurve/.style={poporange,line width=1.8pt,no markers,smooth},
}
% Même style disponible dans un \draw TikZ simple (hors pgfplots)
\tikzset{popcurve/.style={poporange,line width=1.8pt}}
\tikzset{popgrid/.style={popline,very thin}}
\colorlet{popcurve}{poporange} % le nom du style est parfois utilisé comme couleur

% ============================================================
% Pill / squiggle
% ============================================================
\newcommand{\poppill}[3]{%
  \tikz[baseline=-0.55ex]\node[draw=popink,line width=1.2pt,rounded corners=2.6pt,%
    fill=#2,inner xsep=6.5pt,inner ysep=3pt]{%
    \popxboldls\fontsize{7}{7}\selectfont\color{#3}\MakeUppercase{#1}};%
}
\newcommand{\popsquiggle}[1]{%
  \tikz[baseline=-0.4ex]{
    \draw[#1,line width=2.6pt,line cap=round]
      plot [smooth] coordinates {(0,0.09) (0.34,0.01) (0.68,0.21) (1.02,0.01) (1.36,0.10)};
  }%
}
% Mot-clé surligné (marqueur orange sous le texte)
\newcommand{\cle}[1]{{\popxbold\color{popdark}#1}}

% ============================================================
% En-tête / pied
% ============================================================
\pagestyle{fancy}
\fancyhf{}
\renewcommand{\headrulewidth}{0pt}
\renewcommand{\footrulewidth}{0pt}
\lhead{\raisebox{-0.15ex}{\poplogo\fontsize{13}{13}\selectfont\color{popink}OnBuch\color{poporange}+}}
\rhead{\poppill{\DOCMATIERE\ \textbullet\ \DOCNIVEAU\ \textbullet\ Leçon \DOCLECON}{white}{popink}}
\lfoot{\popxbold\fontsize{7.6}{7.6}\selectfont\color{popmuted}\MakeUppercase{Cours signé OnBuch+}\ \textbullet\ {\popsemi\DOCTITRE}}
\rfoot{\tikz[baseline=-0.6ex]\node[rounded corners=8.5pt,fill=popink,minimum height=17pt,minimum width=17pt,inner xsep=5pt,inner sep=0]{\popxbold\fontsize{8}{8}\selectfont\color{popcream}\thepage\,/\,\pageref*{LastPage}};}

% ============================================================
% Titres
% ============================================================
\newcommand{\chaptitre}[2]{%
  \begin{tcolorbox}[enhanced,colback=poporange,colframe=popink,boxrule=2.6pt,arc=11pt,
      drop shadow=popink,left=15pt,right=15pt,top=12pt,bottom=11pt,before skip=3pt,after skip=18pt]
    \poppill{#2}{popink}{popcream}\par\vspace{9pt}
    {\raggedright\hyphenpenalty=10000\exhyphenpenalty=10000\popdisplay\fontsize{22}{24}\selectfont\color{popcream}\MakeUppercase{#1}\par}
  \end{tcolorbox}
}
% Section numérotée : pastille orange avec le numéro + titre Archivo
\newcounter{popsec}
\newcommand{\coursec}[1]{%
  \stepcounter{popsec}\setcounter{popsub}{0}%
  \par\vspace{16pt}\noindent
  \begin{minipage}{\linewidth}
  \tikz[baseline=-0.7ex]\node[draw=popink,line width=1.6pt,fill=poporange,rounded corners=4pt,minimum size=22pt,inner sep=0]{\popdisplay\fontsize{12}{12}\selectfont\color{popcream}\thepopsec};%
  \hspace{8pt}{\popdisplay\fontsize{15}{17}\selectfont\color{popink}\MakeUppercase{#1}}\par
  \vspace{4pt}\noindent{\color{poporange}\rule{36pt}{3.6pt}}
  \end{minipage}\par\nopagebreak\vspace{8pt}
}
\newcounter{popsub}
\newcommand{\courssub}[1]{%
  \stepcounter{popsub}%
  \par\vspace{9pt}\noindent{\popxbold\fontsize{11.5}{13}\selectfont\color{popink}{\color{poporange}\thepopsec.\thepopsub}\hspace{6pt}#1}\par\nopagebreak\vspace{4pt}
}

% ============================================================
% Boîtes pédagogiques — même recette : fond blanc, bordure épaisse colorée,
% ombre dure encre, badge plein. Titre optionnel [#1].
% ============================================================
\tcbset{popbase/.style={breakable,enhanced,colback=white,boxrule=2.3pt,arc=9pt,
  shadow={3pt}{-3pt}{0pt}{popink},left=14pt,right=14pt,top=11pt,bottom=11pt,before skip=11pt,after skip=15pt,
  fontupper=\popsemi\fontsize{10.6}{14.5}\selectfont\color{popink},
  fontlower=\popsemi\fontsize{10.6}{14.5}\selectfont\color{popink}}}
% Coche dessinée (le glyphe ✓ n'existe pas dans les polices du document)
\newcommand{\popcheck}[1]{\tikz[baseline=-0.5ex]\draw[#1,line width=1.6pt,line cap=round,line join=round](-0.13,0.0)--(-0.04,-0.09)--(0.13,0.1);}
\providecommand{\coche}{\popcheck{popgreen}}
\providecommand{\resultat}[1]{\par\smallskip\noindent\fcolorbox{popgreen}{popgreenL}{\parbox{\dimexpr\linewidth-2\fboxsep-2\fboxrule\relax}{\textbf{Résultat.} #1}}\par\smallskip}
\newcommand{\popboxhead}[3]{\poppill{#1}{#2}{white}\ifx\relax#3\relax\else\hspace{6pt}{\popxbold\fontsize{10.6}{12}\selectfont\color{popink}#3}\fi\par\vspace{7pt}}

\newtcolorbox{aretenir}[1][]{popbase,colframe=popgreen,before upper={\popboxhead{À retenir}{popgreen}{#1}}}
\newtcolorbox{exemplebox}[1][]{popbase,colframe=poporange,before upper={\popboxhead{Exemple}{poporange}{#1}}}
\newtcolorbox{definition}[1][]{popbase,colframe=popblue,before upper={\popboxhead{Définition}{popblue}{#1}}}
\newtcolorbox{propriete}[1][]{popbase,colframe=poppurple,before upper={\popboxhead{Propriété}{poppurple}{#1}}}
\newtcolorbox{methode}[1][]{popbase,colframe=popgold,before upper={\popboxhead{Méthode}{popgold}{#1}}}
\newtcolorbox{attention}[1][]{popbase,colframe=poppink,before upper={\popboxhead{Attention}{poppink}{#1}}}
\newtcolorbox{experience}[1][]{popbase,colframe=popblue,colback=popblueL,before upper={\popboxhead{Expérience}{popblue}{#1}}}
% « Le savais-tu ? » : carte violette pleine (contexte, culture, Cameroun)
\newtcolorbox{savaistu}[1][]{popbase,colback=poppurple,colframe=popink,
  fontupper=\popsemi\fontsize{10.4}{14.2}\selectfont\color{white},
  before upper={\poppill{Le savais-tu\,?}{white}{poppurple}\ifx\relax#1\relax\else\hspace{6pt}{\popxbold\color{white}#1}\fi\par\vspace{7pt}}}
% Exercice résolu : énoncé en haut, \tcblower puis la solution sur fond crème
\newcounter{popexo}\newcounter{popexr}
\newtcolorbox{exoresolu}[1][]{popbase,colframe=poporange,colbacklower=popcream,
  segmentation style={popink,line width=1.2pt,dashed},
  before upper={\stepcounter{popexr}\popboxhead{Exercice résolu \thepopexr}{poporange}{#1}},
  before lower={\poppill{Solution}{popink}{popcream}\par\vspace{6pt}}}
% Exercice (non résolu en place) — la correction est dans la partie Corrigés
\newtcolorbox{exercice}[1][]{popbase,colframe=popink,
  before upper={\stepcounter{popexo}\popboxhead{Exercice \thepopexo}{popink}{#1}}}
% Corrigé
\newtcolorbox{corrige}[1][]{popbase,colframe=popgreen,colback=popgreenL,
  before upper={\popboxhead{Corrigé}{popgreen}{#1}}}
% Objectifs / prérequis / bilan
\newtcolorbox{objectifs}{popbase,colframe=popink,colback=poporangeL,
  before upper={\popboxhead{Objectifs de la leçon}{popink}{}}}
\newtcolorbox{prerequis}{popbase,colframe=popink,colback=popblueL,
  before upper={\popboxhead{Prérequis}{popblue}{}}}
\newtcolorbox{bilan}{popbase,colframe=popink,colback=white,boxrule=2.8pt,
  before upper={\popboxhead{Fiche bilan}{poporange}{L'essentiel à connaître pour l'examen}}}
% Figure encadrée avec légende
\newtcolorbox{popfigure}[1][]{enhanced,breakable=false,colback=white,colframe=popline,boxrule=1.4pt,arc=8pt,
  left=10pt,right=10pt,top=10pt,bottom=8pt,before skip=10pt,after skip=12pt,halign=center,
  lower separated=false,halign lower=center,
  fontlower=\popsemi\fontsize{9}{11.5}\selectfont\color{popmuted},
  #1}
\newcommand{\legende}[1]{\par\vspace{6pt}{\popsemi\fontsize{9}{11.5}\selectfont\color{popmuted}#1}}

% ============================================================
% Signature OnBuch+
% ============================================================
\newcommand{\poplogoinline}[1]{{\poplogo\fontsize{#1}{#1}\selectfont\color{popink}OnBuch\color{poporange}+}}
\newcommand{\signatureonbuch}{%
  \begin{tcolorbox}[enhanced,colback=popink,colframe=popink,boxrule=2.2pt,arc=10pt,
      drop shadow=poporange,left=14pt,right=14pt,top=10pt,bottom=10pt,before skip=0pt,after skip=0pt]
    \tikz[baseline=-0.7ex]\node[circle,fill=popgreen,draw=popcream,line width=1.3pt,minimum size=22pt,inner sep=0]{\popcheck{white}};%
    \hspace{9pt}\begin{minipage}[c]{0.62\linewidth}
      {\popxbold\fontsize{12}{13}\selectfont\color{popcream}Signé\ \textbullet\ {\poplogo OnBuch\color{poporange}+}}\par\vspace{2pt}
      {\raggedright\popsemi\fontsize{8}{10.5}\selectfont\color{popline}\MakeUppercase{Équipe pédagogique OnBuch+}\\\MakeUppercase{Conforme au programme officiel MINESEC}\par}
    \end{minipage}\hfill
    {\poplogo\fontsize{22}{22}\selectfont\color{popcream}OnBuch\color{poporange}+}
  \end{tcolorbox}%
}

% ============================================================
% Page de garde
% #1 titre · #2 matière · #3 niveau · #4 module · #5 n° leçon · #6 durée ·
% #7 description / objectifs (texte ou itemize)
% ============================================================
\newcommand{\pagegarde}[7]{%
  \thispagestyle{empty}
  \newgeometry{a4paper,margin=1.7cm,top=1.6cm,bottom=1.4cm}%
  \begin{tikzpicture}[remember picture,overlay]
    \fill[poporange!18] ([xshift=-3.2cm,yshift=-3.6cm]current page.north east) circle (4.6cm);
    \fill[poppurple!14] ([xshift=2.4cm,yshift=7.6cm]current page.south west) circle (3.8cm);
  \end{tikzpicture}%
  {\poplogo\fontsize{30}{30}\selectfont\color{popink}OnBuch\color{poporange}+}\hfill
  \raisebox{0.6ex}{\poppill{Cours signé \textbullet\ Premium}{popink}{popcream}}\par
  \vspace{1pt}\popsquiggle{poppurple}\par
  \vspace{8pt}
  \begin{tcolorbox}[enhanced,colback=poporange,colframe=popink,boxrule=2.8pt,arc=13pt,
      drop shadow=popink,left=18pt,right=18pt,top=12pt,bottom=13pt,after skip=10pt]
    \poppill{#2 \textbullet\ #3}{popink}{popcream}\hspace{5pt}\poppill{#4}{white}{popink}\par\vspace{8pt}
    {\popdisplay\fontsize{14}{14}\selectfont\color{popink}LEÇON #5}\par\vspace{4pt}
    {\raggedright\hyphenpenalty=10000\exhyphenpenalty=10000\popdisplay\fontsize{25}{27}\selectfont\color{popcream}\MakeUppercase{#1}\par}
    \vspace{8pt}
    {\popxbold\fontsize{9.5}{9.5}\selectfont\color{popink}\MakeUppercase{Durée conseillée : #6}}
  \end{tcolorbox}
  \begin{tcolorbox}[enhanced,colback=white,colframe=popink,boxrule=2.2pt,arc=10pt,
      drop shadow=popink,left=16pt,right=16pt,top=10pt,bottom=10pt,after skip=8pt]
    \poppill{Dans cette leçon}{poppurple}{white}\par\vspace{5pt}
    \popsemi\fontsize{9.3}{12.6}\selectfont\color{popink}#7
  \end{tcolorbox}
  \noindent\begin{minipage}[t]{0.487\linewidth}
    \begin{tcolorbox}[enhanced,colback=popgreen,colframe=popink,boxrule=2.2pt,arc=10pt,
        drop shadow=popink,left=11pt,right=11pt,top=10pt,bottom=10pt,height=3.1cm,valign=center]
      \tcbox[colback=white,colframe=popink,boxrule=1.3pt,arc=5pt,boxsep=0pt,nobeforeafter,
        left=3pt,right=3pt,top=3pt,bottom=3pt,valign=center]{\qrcode[height=1.9cm]{https://play.google.com/store/apps/details?id=com.luvvix.onbuch&hl=fr}}%
      \hspace{8pt}\begin{minipage}[c]{0.5\linewidth}
        {\popdisplay\fontsize{11}{12.5}\selectfont\color{white}\MakeUppercase{Télécharge OnBuch}}\par\vspace{4pt}
        {\raggedright\popsemi\fontsize{8.4}{11}\selectfont\color{white}Gratuit \textbullet\ Google Play\\Tuteur IA, annales, concours, résultats\par}
      \end{minipage}
    \end{tcolorbox}
  \end{minipage}\hfill
  \begin{minipage}[t]{0.487\linewidth}
    \begin{tcolorbox}[enhanced,colback=poppurple,colframe=popink,boxrule=2.2pt,arc=10pt,
        drop shadow=popink,left=11pt,right=11pt,top=10pt,bottom=10pt,height=3.1cm,valign=center]
      \tikz[baseline=-0.7ex]\node[circle,fill=white,draw=popink,line width=1.3pt,minimum size=1.6cm,inner sep=0]{\popdisplay\fontsize{24}{24}\selectfont\color{poppurple}+};%
      \hspace{8pt}\begin{minipage}[c]{0.6\linewidth}
        {\popdisplay\fontsize{11}{12.5}\selectfont\color{white}\MakeUppercase{Passe à OnBuch+}}\par\vspace{4pt}
        {\raggedright\popsemi\fontsize{8.4}{11}\selectfont\color{white}Tous les cours signés, Tuteur IA illimité, fiches PDF — onglet OnBuch+ de l'app\par}
      \end{minipage}
    \end{tcolorbox}
  \end{minipage}
  \vspace{8pt}
  \popcontenu
  \vfill
  \signatureonbuch
  \restoregeometry
  \newpage
}

% Rangée « Ce cours contient » (page de garde)
\newcommand{\poptile}[3]{% #1 couleur #2 gros texte #3 légende
  \begin{tcolorbox}[enhanced,colback=#1,colframe=popink,boxrule=2pt,arc=9pt,shadow={2.5pt}{-2.5pt}{0pt}{popink},
      left=6pt,right=6pt,top=9pt,bottom=9pt,halign=center,valign=center,height=1.75cm,width=\linewidth,nobeforeafter]
    {\popdisplay\fontsize{12.5}{13}\selectfont\color{white}\MakeUppercase{#2}}\par\vspace{3pt}
    {\centering\popsemi\fontsize{7.8}{9.5}\selectfont\color{white}#3\par}
  \end{tcolorbox}}
\newcommand{\popcontenu}{%
  \par\noindent{\popxboldls\fontsize{7.5}{7.5}\selectfont\color{popmuted}CE COURS SIGNÉ CONTIENT}\par\vspace{7pt}
  \noindent\begin{minipage}[t]{0.235\linewidth}\poptile{popblue}{Cours}{complet, illustré, pas à pas}\end{minipage}\hfill
  \begin{minipage}[t]{0.235\linewidth}\poptile{poporange}{Méthodes}{et exercices résolus}\end{minipage}\hfill
  \begin{minipage}[t]{0.235\linewidth}\poptile{poppink}{Exercices}{type examen + corrigés}\end{minipage}\hfill
  \begin{minipage}[t]{0.235\linewidth}\poptile{popgreen}{Fiche bilan}{l'essentiel à retenir}\end{minipage}\par}

% Sommaire
\newtcolorbox{sommaire}{enhanced,colback=white,colframe=popink,boxrule=2.2pt,arc=10pt,
  drop shadow=popink,left=18pt,right=18pt,top=13pt,bottom=13pt,before skip=6pt,after skip=20pt,
  before upper={\poppill{Sommaire}{poppurple}{white}\par\vspace{9pt}\popsemi\fontsize{10.6}{14.5}\selectfont\color{popink}}}

% Signature de fin de cours — poussée tout en bas de la dernière page
\newcommand{\findecours}{%
  \par\vfill\nopagebreak
  \begin{center}\popsquiggle{poporange}\end{center}\vspace{4pt}
  \signatureonbuch
}
\AtBeginDocument{\def\llbracket{\mathopen{[\mkern-3.5mu[}}\def\rrbracket{\mathclose{]\mkern-3.5mu]}}} % intervalles d'entiers (glyphe ⟦ absent de la police)
% Glyphes absents de Fira Math : redéfinis à partir de glyphes disponibles
\AtBeginDocument{%
  \def\setminus{\mathbin{\backslash}}\let\smallsetminus\setminus
  \def\vdots{\mathord{\vbox{\baselineskip3.2pt\lineskiplimit0pt\kern4pt\hbox{.}\hbox{.}\hbox{.}}}}%
  \def\ddots{\mathinner{\mkern1mu\raise6.5pt\hbox{.}\mkern2mu\raise3.6pt\hbox{.}\mkern2mu\raise0.7pt\hbox{.}\mkern1mu}}%
  \def\triangle{\mathord{\scalebox{0.9}{$\symup{\Delta}$}}}%
  \def\nabla{\mathord{\raisebox{\depth}{\scalebox{1}[-1]{$\symup{\Delta}$}}}}%
  \def\checkmark{\popcheck{popdark}}%
  \def\star{\mathbin{\ast}}\def\bigstar{\mathord{\ast}}%
  \def\diamond{\mathbin{\scalebox{0.6}{$\Diamond$}}}%
  \def\bigcup{\mathop{\mathchoice{\raisebox{-0.3em}{\scalebox{1.7}{$\cup$}}}{\scalebox{1.25}{$\cup$}}{\cup}{\cup}}\displaylimits}%
  \def\bigcap{\mathop{\mathchoice{\raisebox{-0.3em}{\scalebox{1.7}{$\cap$}}}{\scalebox{1.25}{$\cap$}}{\cap}{\cap}}\displaylimits}%
  \def\lesseqqgtr{\mathrel{\lessgtr}}\def\bigoplus{\mathop{\oplus}\displaylimits}%
}
\providecommand{\ssi}{si et seulement si} % abréviation fréquente
\AtBeginDocument{\providecommand{\wideparen}{\overparen}} % arc de cercle

\begin{document}

\pagegarde{Solutions aqueuses}{PCT}{3ème}{Module 1 : Chimie}{N06}{2 séances (1 h chacune)}{Cette leçon étudie les solutions aqueuses : dissolution d'un soluté dans un solvant, notion de solubilité, calcul de la concentration molaire et mesure du pH pour distinguer les solutions acides, neutres et basiques. Elle relie ces notions à des situations concrètes de la vie courante camerounaise (préparation de boissons, traitement de l'eau, agriculture).
\vspace{4pt}
\begin{multicols}{2}\small
\begin{itemize}
\item À la fin de cette leçon, je sais définir les termes dissolution, soluté, solvant, solution et solubilité.
\item À la fin de cette leçon, je sais identifier le soluté et le solvant dans des situations concrètes de la vie courante (eau sucrée, eau salée, potion médicamenteuse).
\item À la fin de cette leçon, je sais réaliser et schématiser une dissolution en respectant les consignes de sécurité.
\item À la fin de cette leçon, je sais interpréter une expérience de saturation pour dégager la notion de solubilité.
\item À la fin de cette leçon, je sais calculer une concentration molaire à l'aide de la relation C = n/V.
\item À la fin de cette leçon, je sais mesurer le pH d'une solution avec du papier pH et classer la solution en acide, neutre ou basique.
\end{itemize}\end{multicols}}

\begin{sommaire}
\begin{multicols}{2}
\begin{enumerate}
\item Dissolution, soluté, solvant, solution
\item Solubilité et saturation
\item Concentration molaire C = n/V
\item Notion de pH : solutions acides, neutres, basiques
\item Sécurité et applications concrètes
\item Activité d'intégration
\item Méthodes et exercices résolus
\item Exercices
\item Corrigés des exercices
\item Fiche bilan
\end{enumerate}
\end{multicols}
\end{sommaire}

\begin{objectifs}
\begin{itemize}
\item À la fin de cette leçon, je sais définir les termes dissolution, soluté, solvant, solution et solubilité.
\item À la fin de cette leçon, je sais identifier le soluté et le solvant dans des situations concrètes de la vie courante (eau sucrée, eau salée, potion médicamenteuse).
\item À la fin de cette leçon, je sais réaliser et schématiser une dissolution en respectant les consignes de sécurité.
\item À la fin de cette leçon, je sais interpréter une expérience de saturation pour dégager la notion de solubilité.
\item À la fin de cette leçon, je sais calculer une concentration molaire à l'aide de la relation C = n/V.
\item À la fin de cette leçon, je sais mesurer le pH d'une solution avec du papier pH et classer la solution en acide, neutre ou basique.
\item À la fin de cette leçon, je sais interpréter la valeur du pH de solutions usuelles (jus de citron, eau de forage, eau de Javel) pour prendre des décisions pratiques.
\item À la fin de cette leçon, je sais appliquer les consignes de sécurité lors de la manipulation de solutions acides et basiques.
\end{itemize}
\end{objectifs}

\begin{prerequis}
\begin{itemize}
\item Notion de mélange et de corps pur (classe de 4ème)
\item Quantité de matière : la mole et la masse molaire (leçon précédente du Module 1)
\item Écriture et lecture de formules chimiques simples (H2O, NaCl)
\item Unités de volume (litre, cm3) et conversions simples
\item Utilisation de la verrerie de base : bécher, agitateur, balance
\end{itemize}
\end{prerequis}

\newpage

\coursec{Dissolution, soluté, solvant, solution}

Au marché Mokolo, une vendeuse prépare une solution de potasse pour laver les légumes : elle dissout une certaine masse de cendres dans un litre d'eau. À l'hôpital de district, l'infirmier vérifie le pH de l'eau de boisson avant de la distribuer aux malades. Pourquoi dit-on que le sel « disparaît » dans l'eau alors qu'il est toujours là ? Comment savoir si une solution est trop concentrée ou trop acide pour être consommée ? Cette leçon répond à ces questions de la vie quotidienne.

\courssub{Expérience introductive : le sel qui « disparaît »}

\begin{experience}[Dissolution du sel de cuisine dans l'eau]
\textbf{Matériel :} un bécher, de l'eau du robinet, du sel de cuisine (NaCl), une cuillère, une balance de précision.
\textbf{Protocole :}
\begin{enumerate}
    \item Mesurer la masse $m_{\text{eau}}$ de l'eau versée dans le bécher (par exemple \SI{100}{g}).
    \item Mesurer la masse $m_{\text{sel}}$ d'une cuillère de sel (par exemple \SI{5}{g}).
    \item Verser le sel dans l'eau et remuer avec la cuillère jusqu'à ce qu'il n'y ait plus de cristaux visibles au fond.
    \item Observer l'aspect du mélange. Goûter \textbf{avec une extrême prudence} (une goutte sur le bout du doigt) pour vérifier la saveur salée.
    \item Peser la masse $m_{\text{mélange}}$ du bécher + solution obtenue.
\end{enumerate}
\textbf{Observations :}
\begin{itemize}
    \item Au début, on voit des cristaux blancs au fond de l'eau.
    \item En remuant, les cristaux deviennent de plus en plus petits, puis invisibles. Le liquide reste transparent et incolore.
    \item Le goût salé persiste.
    \item $m_{\text{mélange}} = m_{\text{eau}} + m_{\text{sel}}$ (masse conservée).
\end{itemize}
\textbf{Interprétation :} Les cristaux de sel ne sont pas détruits. Ils se sont \textbf{dispersés} à l'échelle microscopique entre les molécules d'eau. Le mélange obtenu est \textbf{homogène} : on ne distingue plus qu'une seule phase.
\end{experience}

\begin{popfigure}
\begin{tikzpicture}[scale=0.9, every node/.style={font=\small}]
    % Styles
    \tikzset{
        water/.style={draw=popblue, fill=popblue!15, thick},
        ionNa/.style={draw=poporange, fill=poporange!80, circle, minimum size=6pt, inner sep=0},
        ionCl/.style={draw=popgreen, fill=popgreen!80, circle, minimum size=7pt, inner sep=0},
        crystal/.style={draw=popdark, fill=popgray!30, thick},
        beaker/.style={draw=popdark, thick},
        arrow/.style={->, >=Stealth, thick, popdark},
        labelbox/.style={rectangle, draw=popdark, fill=popcream, rounded corners, inner sep=4pt, align=center, font=\small}
    }
    % Etape 1 : Cristaux dans l'eau
    \begin{scope}[xshift=0cm]
        \draw[beaker] (0,0) -- (0,3) -- (2.5,3) -- (2.5,0);
        \draw[water] (0.1,0.1) rectangle (2.4,2.0);
        % Cristaux au fond
        \foreach \x/\y in {0.4/0.3, 0.8/0.25, 1.3/0.35, 1.8/0.2, 2.1/0.4} {
            \draw[crystal] (\x,\y) rectangle ++(0.3,0.3);
        }
        \node[labelbox, anchor=north] at (1.25, 3.3) {Étape 1 : Avant agitation};
        \node[anchor=north, font=\tiny, align=center] at (1.25, -0.3){Cristaux de NaCl (solide) + Eau (liquide)\\Mélange \textbf{hétérogène} (2 phases)};
    \end{scope}
    % Flèche
    \draw[arrow] (3.0, 1.5) -- (4.0, 1.5) node[midway, above] {Agitation};
    % Etape 2 : Dispersion en cours
    \begin{scope}[xshift=5.5cm]
        \draw[beaker] (0,0) -- (0,3) -- (2.5,3) -- (2.5,0);
        \draw[water] (0.1,0.1) rectangle (2.4,2.8);
        % Ions en dispersion
        \foreach \i in {1,...,15} {
            \pgfmathsetmacro{\rx}{0.2+2.0*rnd}
            \pgfmathsetmacro{\ry}{0.2+2.4*rnd}
            \ifodd\i
                \node[ionNa] at (\rx,\ry) {$\ce{Na+}$};
            \else
                \node[ionCl] at (\rx,\ry) {$\ce{Cl-}$};
            \fi
        }
        % Quelques cristaux résiduels
        \draw[crystal] (0.5,0.3) rectangle ++(0.25,0.25);
        \node[labelbox, anchor=north] at (1.25, 3.3) {Étape 2 : En cours};
        \node[anchor=north, font=\tiny, align=center] at (1.25, -0.3){Ions $\ce{Na+}$ et $\ce{Cl-}$ se dispersent\\dans l'eau};
    \end{scope}
    % Flèche
    \draw[arrow] (8.5, 1.5) -- (9.5, 1.5) node[midway, above] {Temps};
    % Etape 3 : Solution homogène
    \begin{scope}[xshift=11cm]
        \draw[beaker] (0,0) -- (0,3) -- (2.5,3) -- (2.5,0);
        \draw[water] (0.1,0.1) rectangle (2.4,2.8);
        % Ions uniformément répartis
        \foreach \i in {1,...,30} {
            \pgfmathsetmacro{\rx}{0.2+2.0*rnd}
            \pgfmathsetmacro{\ry}{0.2+2.4*rnd}
            \ifodd\i
                \node[ionNa] at (\rx,\ry) {$\ce{Na+}$};
            \else
                \node[ionCl] at (\rx,\ry) {$\ce{Cl-}$};
            \fi
        }
        \node[labelbox, anchor=north] at (1.25, 3.3) {Étape 3 : Solution finale};
        \node[anchor=north, font=\tiny, align=center] at (1.25, -0.3){Solution aqueuse de NaCl\\Mélange \textbf{homogène} (1 phase)};
    \end{scope}
    % Légende ions
    \begin{scope}[yshift=-4.5cm, xshift=2cm]
        \node[ionNa] at (0,0) {}; \node[anchor=west] at (0.4,0) {Ion sodium $\ce{Na+}$};
        \node[ionCl] at (3,0) {}; \node[anchor=west] at (3.4,0) {Ion chlorure $\ce{Cl-}$};
        \draw[water] (6,-0.2) rectangle (7.5,0.2); \node[anchor=west] at (7.7,0) {Molécule d'eau $\ce{H2O}$};
    \end{scope}
\end{tikzpicture}
\legende{Dissolution du chlorure de sodium (sel de cuisine) dans l'eau : passage d'un mélange hétérogène (cristaux + eau) à une solution aqueuse homogène (ions $\ce{Na+}$ et $\ce{Cl-}$ dispersés).}
\end{popfigure}

\courssub{Notions fondamentales : dissolution, soluté, solvant, solution}

\begin{definition}[Dissolution]
Une \textbf{dissolution} est le phénomène physique au cours duquel les particules (atomes, molécules ou ions) d'un corps, appelé \textbf{soluté}, se dispersent de manière uniforme au milieu des particules d'un autre corps, appelé \textbf{solvant}, pour former un mélange homogène appelé \textbf{solution}.
\end{definition}

\begin{definition}[Soluté, solvant, solution]
\begin{itemize}
    \item \textbf{Soluté :} Corps qui se dissout (il est en général en plus petite quantité). Exemple : le sel, le sucre, le gaz carbonique, la potasse.
    \item \textbf{Solvant :} Corps qui dissout le soluté (il est en général en plus grande quantité). Dans cette leçon, le solvant est \textbf{l'eau} : on parle de \textbf{solution aqueuse}.
    \item \textbf{Solution :} Mélange homogène obtenu après dissolution. Une solution est transparente (mais pas nécessairement incolore) et ne se sépare pas au repos.
\end{itemize}
\end{definition}

\begin{aretenir}[Solution aqueuse]
Une \textbf{solution aqueuse} est une solution dont le solvant est l'eau. On note souvent l'eau comme solvant par l'indice $(aq)$ après la formule du soluté dissous.
\textbf{Exemple :} $\ce{NaCl(s) ->[eau] NaCl(aq)}$ (en réalité $\ce{Na+(aq) + Cl-(aq)}$).
\end{aretenir}

\begin{attention}[Piège : « Le sel a disparu » ?]
\textbf{FAUX.} Le soluté ne disparaît pas, il se \textbf{disperse} à l'échelle microscopique (niveau des particules).
\begin{itemize}
    \item La \textbf{masse totale} est conservée : $m_{\text{solution}} = m_{\text{solvant}} + m_{\text{soluté}}$.
    \item Les propriétés du soluté persistent (goût salé, conductivité électrique pour le sel, couleur pour le sirop de menthe).
    \item Au BEPC, écris toujours : « Les particules du soluté se dispersent entre les molécules du solvant. »
\end{itemize}
\end{attention}

\courssub{Distinction : mélange homogène vs mélange hétérogène}

\begin{experience}[Comparaison : eau + sel vs eau + huile vs eau + sable]
\textbf{Protocole :} Dans trois tubes à essai identiques, mettre \SI{5}{mL} d'eau. Ajouter : (1) une pincée de sel, (2) \SI{2}{mL} d'huile de cuisine, (3) une cuillère de sable. Agiter énergiquement chaque tube pendant 30 s. Observer immédiatement et après 5 min de repos.
\textbf{Observations et interprétation :}
\begin{center}
\begin{tabularx}{\linewidth}{|l|X|X|X|}\hline
\textbf{Mélange} & \textbf{Aspect immédiat} & \textbf{Aspect après repos} & \textbf{Type de mélange} \\ \hline
Eau + Sel &Transparent, incolore, pas de dépôt & Identique (stable) & \textbf{Homogène} (1 phase) \\ \hline
Eau + Huile & Trouble, gouttelettes & Séparation en 2 couches (huile au-dessus) & \textbf{Hétérogène} (2 phases) \\ \hline
Eau + Sable & Trouble, particules en suspension & Sable au fond (dépôt rapide) & \textbf{Hétérogène} (2 phases) \\ \hline
\end{tabularx}
\end{center}
\end{experience}

\begin{propriete}[Critères de distinction]
\begin{itemize}
    \item \textbf{Mélange homogène (solution) :} Une seule phase visible, aspect uniforme, transparent (ou coloré uniformément), stable au repos, ne se filtre pas (les particules sont trop petites).
    \item \textbf{Mélange hétérogène :} Au moins deux phases visibles (ex : liquide + solide, ou deux liquides non miscibles), aspect non uniforme, instable au repos (dépôt ou séparation), peut souvent se séparer par filtration ou décantation.
\end{itemize}
\end{propriete}

\begin{popfigure}
\begin{tikzpicture}
    \begin{axis}[
        width=\linewidth, height=6cm,
        hide axis,
        xmin=0, xmax=10, ymin=0, ymax=4,
    ]
    % Tableau Homogène
    \node[draw=popblue, thick, fill=popblueL, rounded corners, minimum width=4.5cm, minimum height=3.5cm, anchor=north west] (homog) at (0.5, 3.5) {};
    \node[font=\bfseries\large, popblue] at (2.75, 3.2) {Mélange HOMOGÈNE};
    \node[align=center, font=\small] at (2.75, 2.6) {1 seule phase visible\\Transparent ou coloré uniformément\\Stable au repos\\Ex: Eau+Sel, Eau+Sucre, Sirop,\\Eau de Javel, Air, Lait (apparence)};
    \node[align=center, font=\tiny, popmuted] at (2.75, 1.2) {Particules : molécules/ions\\Taille < 1 nm\\Non filtrables};
    
    % Tableau Hétérogène
    \node[draw=poporange, thick, fill=poporangeL, rounded corners, minimum width=4.5cm, minimum height=3.5cm, anchor=north east] (heterog) at (9.5, 3.5) {};
    \node[font=\bfseries\large, poporange] at (7.25, 3.2) {Mélange HÉTÉROGÈNE};
    \node[align=center, font=\small] at (7.25, 2.6) {Au moins 2 phases visibles\\Aspect non uniforme (trouble, couches)\\Instable au repos (dépôt/séparation)\\Ex: Eau+Huile, Eau+Sable, Jus de fruit\\avec pulpe, Granit, Vinaigrette};
    \node[align=center, font=\tiny, popmuted] at (7.25, 1.2) {Particules : agrégats/gouttes\\Taille > 100 nm\\Souvent filtrables/décantables};
    
    % Flèche centrale
    \draw[<->, >=Stealth, very thick, poppurple] (4.8, 2.5) -- (5.2, 2.5) node[midway, above, font=\small, poppurple] {Frontière};
    \draw[<->, >=Stealth, very thick, poppurple] (4.8, 1.8) -- (5.2, 1.8) node[midway, above, font=\small, poppurple] {Taille particules};
    \end{axis}
\end{tikzpicture}
\legende{Comparaison visuelle entre mélange homogène (solution) et mélange hétérogène avec exemples du quotidien camerounais.}
\end{popfigure}

\courssub{Exemples concrets de solutions aqueuses au Cameroun}

\begin{exemplebox}[Identifier soluté et solvant]
\begin{enumerate}
    \item \textbf{Eau de Javel diluée (eau de Javel + eau) :} C'est une solution aqueuse. Le solvant est l'\textbf{eau}. Le soluté est l'\textbf{hypochlorite de sodium} ($\ce{NaClO}$), agent blanchissant et désinfectant.
    \item \textbf{Jus de citron sucré (jus de citron + sucre + eau) :} C'est une solution aqueuse. Le solvant est l'\textbf{eau} (majoritaire). Il y a \textbf{deux solutés} : l'\textbf{acide citrique} (goût acide) et le \textbf{saccharose} (goût sucré).
    \item \textbf{Boisson gazeuse (Coca, Youki, etc.) :} Solvant = \textbf{eau}. Solutés = \textbf{sucres}, \textbf{acide phosphorique}, \textbf{arômes}, \textbf{dioxyde de carbone} ($\ce{CO2}$) dissous sous pression.
    \item \textbf{Sérum physiologique (sachet ORS / sel de réhydratation) :} Solvant = \textbf{eau bouillie refroidie}. Solutés = \textbf{chlorure de sodium} (NaCl), \textbf{chlorure de potassium} (KCl), \textbf{glucose}, \textbf{citrate de sodium}.
\end{enumerate}
\end{exemplebox}

\begin{savaistu}[L'eau, « solvant universel » ?]
On surnomme l'eau « \textbf{solvant universel} » car elle dissout une plus grande variété de substances (sels, sucres, acides, bases, gaz, alcools) que tout autre liquide connu. Cette propriété exceptionnelle vient de sa molécule \textbf{polaires} ($\ce{H2O}$) : elle attire les ions positifs par son atome d'oxygène (charge partielle négative) et les ions négatifs par ses atomes d'hydrogène (charges partielles positives), arrachant ainsi les ions du cristal (solvatation).
\textbf{Attention :} L'eau ne dissout \textbf{PAS} les corps gras (huiles, essence, graisses, plastiques) : ils sont \textbf{apolaires} et ne sont pas attirés par les molécules d'eau. C'est pourquoi l'huile flotte sur l'eau et ne se mélange pas.
\end{savaistu}

\courssub{Exercice résolu : analyser une étiquette}

\begin{exoresolu}[Composition d'une boisson énergisante]
\textbf{Énoncé :} Sur une canette de boisson énergisante vendue à Yaoundé, on lit : « Eau, sucre, acide citrique, taurine, caféine, arômes, gaz carbonique, colorant caramel. Volume : \SI{330}{mL}. Masse de sucre : \SI{37}{g}. »
\begin{enumerate}
    \item Quel est le solvant de cette boisson ?
    \item Citez trois solutés solides et un soluté gazeux.
    \item Cette boisson est-elle un mélange homogène ou hétérogène ? Justifiez.
    \item Si on laisse la canette ouverte une nuit, que devient le gaz carbonique ? La masse de la boisson change-t-elle ?
\end{enumerate}
\tcblower
\textbf{Solution détaillée :}
\begin{enumerate}
    \item Le solvant est l'\textbf{eau} (premier ingrédient, majoritaire). C'est une solution aqueuse.
    \item Solutés solides : \textbf{sucre} (saccharose), \textbf{acide citrique}, \textbf{taurine}, \textbf{caféine}, \textbf{colorant caramel}. (Trois suffisent). Soluté gazeux : \textbf{dioxyde de carbone} ($\ce{CO2}$, « gaz carbonique »).
    \item C'est un \textbf{mélange homogène} (solution). Justification : un seul liquide visible, transparent (coloré par le caramel), pas de dépôt au fond, pas de séparation en couches au repos.
    \item Le $\ce{CO2}$ dissous s'échappe dans l'air (dégazage) : la boisson devient « plate ». La \textbf{masse totale de la canette ouverte diminue} car la masse de $\ce{CO2}$ perdue n'est plus dans le système. La masse du liquide seul diminue aussi.
\end{enumerate}
\end{exoresolu}

\begin{methode}[Identifier les constituants d'une solution aqueuse]
\begin{enumerate}
    \item Lire la liste des ingrédients (le premier est généralement le solvant).
    \item Identifier le composant présent en plus grande quantité : c'est le \textbf{solvant} (souvent l'eau).
    \item Tous les autres composants dissous (solides, liquides ou gazeux) sont des \textbf{solutés}.
    \item Vérifier l'homogénéité : aspect unique, pas de séparation spontanée $\rightarrow$ solution.
\end{enumerate}
\end{methode}

\begin{exercice}[À toi de jouer !]
\begin{enumerate}
    \item Un élève verse \SI{10}{g} de sucre dans \SI{100}{g} d'eau et remue jusqu'à dissolution complète. Quelle est la masse de la solution obtenue ? Le sucre a-t-il disparu ? Justifie.
    \item Parmi les mélanges suivants, lesquels sont des solutions aqueuses homogènes ?
    \begin{itemize}
        \item Lait entier frais
        \item Eau de mer
        \item Vinaigre (eau + acide acétique)
        \item Mélange eau + essence (dans un bidon de moto-malaxeur)
        \item Jus de tamarin filtré (eau + pulpe + sucre)
    \end{itemize}
    \item Complète : Dans une solution aqueuse de sulfate de cuivre (bleue), le solvant est \ldots, le soluté est \ldots. Les particules du soluté sont des \ldots (ions $\ce{Cu^{2+}}$ et $\ce{SO4^{2-}}$) dispersés entre les molécules d'\ldots
\end{enumerate}
\end{exercice}

\begin{corrige}[Exercice n°1]
\begin{enumerate}
    \item Masse solution = \SI{100}{g} + \SI{10}{g} = \SI{110}{g}. Non, le sucre ne disparaît pas : ses molécules se dispersent entre les molécules d'eau (conservation de la masse et du goût sucré).
    \item Solutions aqueuses homogènes : \textbf{Eau de mer} (sels dissous), \textbf{Vinaigre} (acide acétique dissous), \textbf{Jus de tamarin filtré} (sucre et acides dissous, pulpe retirée par filtration).
    \begin{itemize}
        \item Lait : émulsion (gouttelettes de gras en suspension) $\rightarrow$ hétérogène (trouble, instable).
        \item Eau + essence : deux liquides non miscibles $\rightarrow$ hétérogène (deux couches).
    \end{itemize}
    \item Solvant : \textbf{eau}. Soluté : \textbf{sulfate de cuivre} ($\ce{CuSO4}$). Particules : \textbf{ions}. Molécules d'\textbf{eau}.
\end{enumerate}
\end{corrige}

\coursec{Solubilité et saturation}

Dans la section précédente, tu as découvert ce qu'est une \cle{solution} : un mélange homogène obtenu en dissolvant un \cle{soluté} dans un \cle{solvant}. Tu sais maintenant que le sel disparaît dans l'eau pour former de l'eau salée. Mais une question naturelle se pose : peut-on dissoudre \emph{n'importe quelle} quantité de sel dans un verre d'eau ? Si tu verses un kilogramme de sel dans un petit verre d'eau, tout va-t-il disparaître ? Évidemment non ! C'est exactement ce que nous allons étudier dans cette section : il existe une \textbf{limite} à la dissolution, et cette limite s'appelle la \cle{solubilité}.

\courssub{Expérience : jusqu'où peut-on dissoudre du sel dans l'eau ?}

\begin{experience}[Dissolution du sel jusqu'à saturation]
\textbf{Matériel :} un bécher, une balance électronique, une baguette de verre (ou une cuillère propre), de l'eau distillée, du sel de cuisine (chlorure de sodium \ce{NaCl}), une éprouvette graduée.

\textbf{Protocole :}
\begin{enumerate}
\item Verser \SI{100}{mL} d'eau distillée à température ambiante (environ \SI{20}{\celsius}) dans le bécher.
\item Ajouter le sel par petites portions de \SI{5}{g}, en agitant énergiquement après chaque ajout jusqu'à dissolution complète.
\item Noter la masse totale de sel ajoutée et observer le bécher après chaque ajout.
\item Poursuivre jusqu'à ce que le sel ne se dissolve plus, même après une agitation prolongée.
\end{enumerate}

\textbf{Observations :}
\begin{itemize}
\item Au début, chaque portion de sel disparaît complètement : la solution reste limpide et homogène.
\item À partir d'environ \SI{36}{g} de sel ajoutés (soit \SI{360}{g} par litre d'eau), des grains de sel restent visibles au fond du bécher, même après plusieurs minutes d'agitation.
\item Le sel qui reste au fond forme un \textbf{dépôt} solide.
\end{itemize}

\textbf{Interprétation :} tant que la solution peut encore dissoudre du sel, on dit qu'elle est \emph{non saturée}. Lorsque le sel ajouté ne se dissout plus, c'est que l'eau a atteint sa \textbf{capacité maximale de dissolution} à cette température : la solution est dite \emph{saturée}. Le dépôt au fond prouve que le soluté en excès ne peut plus se disperser dans le solvant.

\textbf{Conclusion :} il existe une masse maximale de soluté qu'un volume donné de solvant peut dissoudre à une température donnée.
\end{experience}

\begin{popfigure}
\begin{center}
\begin{tikzpicture}[scale=0.9]
% Bécher
\draw[thick,popink] (-2,0) -- (-2,4.2);
\draw[thick,popink] (2,0) -- (2,4.2);
\draw[thick,popink] (-2,0) -- (2,0);
% Eau
\fill[popblue!25] (-1.95,0.05) rectangle (1.95,3.0);
\draw[popblue] (-1.95,3.0) -- (1.95,3.0);
% Dépôt de sel
\fill[popgold!70] (-1.95,0.05) -- (1.95,0.05) -- (1.95,0.45) -- (1.2,0.6) -- (0.3,0.42) -- (-0.5,0.62) -- (-1.3,0.45) -- cycle;
\foreach \x in {-1.6,-1.1,-0.6,-0.1,0.4,0.9,1.4} {\fill[popgold!90] (\x,0.18) circle (0.06);}
% Grains en suspension (peu nombreux près du dépôt)
\foreach \p in {(-1.2,2.3),(0.5,2.6),(1.1,1.8),(-0.4,1.5),(0.9,2.9)} {\fill[popgold!60] \p circle (0.045);}
% Agitateur
\draw[thick,popmuted] (1.3,4.8) -- (0.4,0.9);
% Annotations
\draw[->,thick,popdark] (2.6,3.0) -- (2.1,3.0);
\node[right,font=\small,popdark] at (2.6,3.0) {surface de la solution};
\draw[->,thick,popdark] (2.6,0.35) -- (1.6,0.35);
\node[right,font=\small,popdark,align=left] at (2.6,0.35) {dépôt de sel\\non dissous};
\node[font=\small,popink] at (0,4.6) {Bécher contenant une solution \textbf{saturée}};
\end{tikzpicture}
\end{center}
\legende{Schéma d'un bécher contenant une solution saturée de chlorure de sodium : le sel ajouté en excès ne se dissout plus et forme un dépôt solide au fond.}
\end{popfigure}

\courssub{Solution saturée : définition}

\begin{definition}[Solution saturée]
Une solution est dite \cle{saturée} lorsque, à une température donnée, le solvant ne peut plus dissoudre de soluté supplémentaire : tout soluté ajouté en plus reste à l'état solide et forme un dépôt.

Une solution qui peut encore dissoudre du soluté est dite \cle{non saturée} (ou insaturée).
\end{definition}

\begin{attention}[Saturation et température]
La saturation dépend de la \textbf{température} : une solution saturée à \SI{20}{\celsius} peut devenir non saturée si on la chauffe, car elle pourra alors dissoudre davantage de soluté. Il faut donc \emph{toujours} préciser la température quand on parle de saturation. Dire simplement « la solution est saturée » sans préciser la température est une réponse incomplète au BEPC.
\end{attention}

\courssub{La solubilité : définition et unité}

L'expérience précédente nous a montré qu'il existe une masse maximale de sel dissous. Donnons un nom précis à cette grandeur.

\begin{definition}[Solubilité]
La \cle{solubilité} d'un soluté dans un solvant, notée $s$, est la \textbf{masse maximale} de ce soluté que l'on peut dissoudre dans \textbf{un litre de solvant} (généralement de l'eau), à une \textbf{température donnée}.

Elle s'exprime en grammes par litre : \SI{}{g/L} (ou g$\cdot$L$^{-1}$).
\end{definition}

\begin{exemplebox}[Solubilités de quelques espèces dans l'eau à 20 °C]
\begin{center}
\begin{tabularx}{\linewidth}{|l|X|}\hline
\rowcolor{popblueL} \textbf{Soluté} & \textbf{Solubilité à \SI{20}{\celsius}} \\ \hline
Chlorure de sodium \ce{NaCl} (sel de cuisine) & environ \SI{360}{g/L} \\ \hline
Nitrate de potassium \ce{KNO3} & environ \SI{320}{g/L} \\ \hline
Saccharose (sucre) \ce{C12H22O11} & environ \SI{2000}{g/L} \\ \hline
Carbonate de calcium \ce{CaCO3} (calcaire) & environ \SI{0,014}{g/L} \\ \hline
\end{tabularx}
\end{center}
On constate que la solubilité dépend de la \textbf{nature du soluté} : le sucre est beaucoup plus soluble que le sel, et le calcaire est pratiquement insoluble dans l'eau.
\end{exemplebox}

\begin{aretenir}[Les trois facteurs qui influencent la solubilité]
\begin{enumerate}
\item La \textbf{nature du soluté} (le sucre est plus soluble que le sel) ;
\item La \textbf{nature du solvant} (le sel se dissout dans l'eau mais pas dans l'huile) ;
\item La \textbf{température} (voir ci-dessous).
\end{enumerate}
\end{aretenir}

\courssub{Influence de la température sur la solubilité}

\begin{experience}[Eau chaude, eau froide]
\textbf{Protocole :} préparer deux béchers identiques contenant chacun \SI{100}{mL} d'eau : l'un avec de l'eau froide (\SI{10}{\celsius}), l'autre avec de l'eau chaude (\SI{60}{\celsius}). Ajouter du nitrate de potassium \ce{KNO3} par petites portions dans chaque bécher en agitant, jusqu'à saturation.

\textbf{Observation :} le bécher d'eau chaude dissout une masse de \ce{KNO3} nettement plus grande que le bécher d'eau froide avant d'atteindre la saturation.

\textbf{Conclusion :} pour la plupart des solides, la solubilité \textbf{augmente} lorsque la température augmente.
\end{experience}

\begin{propriete}[Solubilité et température]
La solubilité des solides dans l'eau \textbf{augmente généralement avec la température}. C'est pourquoi on chauffe l'eau pour dissoudre plus vite et plus de sucre ou de sel en cuisine.

Remarque : pour les \textbf{gaz}, c'est l'inverse ! La solubilité des gaz dans l'eau \emph{diminue} quand la température augmente. C'est pour cela qu'une boisson gazeuse chaude perd rapidement ses bulles.
\end{propriete}

\courssub{Lecture d'une courbe de solubilité}

Les résultats de mesures de solubilité à différentes températures sont souvent présentés sous forme d'une \textbf{courbe de solubilité} : $s = f(\text{température})$.

\begin{popfigure}
\begin{center}
\begin{tikzpicture}
\begin{axis}[
width=0.9\linewidth,
height=7.5cm,
xlabel={Température (\si{\celsius})},
ylabel={Solubilité $s$ (\si{g/L})},
xmin=0, xmax=100,
ymin=0, ymax=2500,
xtick={0,20,40,60,80,100},
ytick={0,500,1000,1500,2000,2500},
grid=both,
grid style={popline!40},
legend pos=north west,
legend style={font=\small, fill=white, draw=popline},
axis line style={popink},
label style={font=\small},
tick label style={font=\small}
]
% KNO3 : points expérimentaux approximatifs (g/L dans l'eau)
\addplot[poporange,very thick,smooth] coordinates {
(0,130) (10,210) (20,320) (30,460) (40,640) (50,850) (60,1050) (70,1250) (80,1400) (90,1800) (100,2400)
};
\addlegendentry{\ce{KNO3}}
% NaCl : quasi constant
\addplot[popblue,very thick,smooth] coordinates {
(0,357) (20,360) (40,365) (60,370) (80,380) (100,391)
};
\addlegendentry{\ce{NaCl}}
% Point de lecture NaCl à 20°C
\addplot[only marks,mark=*,popdark,mark size=2.5pt] coordinates {(20,360)};
\draw[dashed,popdark,thin] (axis cs:20,0) -- (axis cs:20,360) -- (axis cs:0,360);
\node[popdark,font=\small,anchor=south west] at (axis cs:22,380) {$s \approx \SI{360}{g/L}$ à \SI{20}{\celsius}};
% Point de lecture KNO3 à 60°C
\addplot[only marks,mark=*,poporange,mark size=2.5pt] coordinates {(60,1050)};
\draw[dashed,poporange,thin] (axis cs:60,0) -- (axis cs:60,1050) -- (axis cs:0,1050);
\node[poporange,font=\small,anchor=south west] at (axis cs:62,1100) {$s \approx \SI{1050}{g/L}$ à \SI{60}{\celsius}};
\end{axis}
\end{tikzpicture}
\end{center}
\legende{Courbes de solubilité du nitrate de potassium \ce{KNO3} et du chlorure de sodium \ce{NaCl} en fonction de la température. La solubilité du \ce{KNO3} augmente fortement avec la température, alors que celle du \ce{NaCl} varie très peu.}
\end{popfigure}

L'examen de ces deux courbes est très instructif :
\begin{itemize}
\item La courbe du \ce{KNO3} \textbf{monte fortement} : sa solubilité passe d'environ \SI{320}{g/L} à \SI{20}{\celsius} à plus de \SI{2400}{g/L} à \SI{100}{\celsius}. Ce sel est beaucoup plus soluble à chaud.
\item La courbe du \ce{NaCl} est presque \textbf{horizontale} : sa solubilité reste voisine de \SI{360}{g/L} quelle que soit la température. Chauffer l'eau ne permet donc pas de dissoudre beaucoup plus de sel de cuisine.
\end{itemize}

\begin{methode}[Exploiter une courbe de solubilité]
Pour lire une solubilité et prévoir une masse dissoute, suis ces étapes :
\begin{enumerate}
\item \textbf{Repérer} la température donnée sur l'axe horizontal (abscisses).
\item \textbf{Monter} verticalement jusqu'à la courbe du soluté étudié.
\item \textbf{Lire} la solubilité $s$ sur l'axe vertical (ordonnées), en \si{g/L}.
\item \textbf{Calculer} la masse maximale dissoute dans le volume $V$ de solvant (eau) donné par la proportionnalité :
\[ m_{\max} = s \times V \quad \text{avec } s \text{ en \si{g/L} et } V \text{ en L.} \]
\item \textbf{Convertir} le volume en litres si nécessaire (rappel : \SI{250}{mL} = \SI{0,250}{L}).
\end{enumerate}
\end{methode}

\begin{exoresolu}[Masse maximale de sel dissoute dans 250 mL d'eau]
\textbf{Énoncé.} La solubilité du chlorure de sodium dans l'eau à \SI{20}{\celsius} vaut $s \approx \SI{360}{g/L}$. Quelle masse maximale de sel peut-on dissoudre dans \SI{250}{mL} d'eau à cette température ? Que se passe-t-il si on ajoute \SI{100}{g} de sel ?
\tcblower
\textbf{Solution détaillée.}

Étape 1 : conversion du volume en litres.
\[ V = \SI{250}{mL} = \SI{0,250}{L} \]

Étape 2 : application de la relation $m_{\max} = s \times V$.
\[ m_{\max} = 360 \times 0{,}250 = \SI{90}{g} \]

Conclusion : on peut dissoudre au maximum \SI{90}{g} de sel dans \SI{250}{mL} d'eau à \SI{20}{\celsius}.

Si on ajoute \SI{100}{g} de sel : seuls \SI{90}{g} se dissolvent ; il reste $100 - 90 = \SI{10}{g}$ de sel solide au fond du récipient. La solution obtenue est \textbf{saturée}.
\end{exoresolu}

\begin{exoresolu}[Lecture graphique avec le KNO3]
\textbf{Énoncé.} En utilisant la courbe de solubilité du \ce{KNO3}, déterminer la masse maximale de \ce{KNO3} que l'on peut dissoudre dans \SI{200}{mL} d'eau à \SI{60}{\celsius}.
\tcblower
\textbf{Solution détaillée.}

Étape 1 : lecture graphique. À \SI{60}{\celsius}, la courbe du \ce{KNO3} donne $s \approx \SI{1050}{g/L}$.

Étape 2 : conversion. $V = \SI{200}{mL} = \SI{0,200}{L}$.

Étape 3 : calcul.
\[ m_{\max} = s \times V = 1050 \times 0{,}200 = \SI{210}{g} \]

Conclusion : on peut dissoudre environ \SI{210}{g} de \ce{KNO3} dans \SI{200}{mL} d'eau à \SI{60}{\celsius}.
\end{exoresolu}

\begin{attention}[Ne pas confondre solubilité et concentration]
C'est l'erreur la plus fréquente au BEPC !
\begin{itemize}
\item La \cle{solubilité} $s$ s'exprime en \textbf{g/L} : c'est une \emph{masse maximale} par litre de \emph{solvant} (eau) pour obtenir une solution saturée.
\item La \cle{concentration molaire} $C = \dfrac{n}{V}$ s'exprime en \textbf{mol/L} : c'est une \emph{quantité de matière} par litre de \emph{solution}, qu'elle soit saturée ou non.
\end{itemize}
Ce sont deux grandeurs différentes, avec des unités différentes. Une solution non saturée a une concentration, mais on ne peut pas parler de \emph{sa} solubilité : la solubilité est une caractéristique du couple soluté-solvant à une température donnée, pas d'une solution particulière.
\end{attention}

\courssub{Application concrète : la saumure, une solution saturée pour conserver les aliments}

\begin{savaistu}[La saumure, une vieille astuce de conservation]
Bien avant les réfrigérateurs, nos grands-parents conservaient la viande, le poisson fumé et certains légumes dans une \cle{saumure} : une solution d'eau très concentrée en sel, souvent proche de la saturation. Le sel en grande quantité empêche le développement des microbes responsables de la pourriture, car il « aspire » l'eau hors de leurs cellules (phénomène d'osmose). Au marché, le poisson salé et séché que l'on trouve partout au Cameroun doit sa longue conservation à ce principe. La saumure est aussi utilisée pour préparer les olives et certains fromages.
\end{savaistu}

Pour préparer une saumure efficace, on dissout du sel dans de l'eau jusqu'à obtenir une solution proche de la saturation : environ \SI{360}{g} de sel par litre d'eau à température ambiante. Si un dépôt de sel apparaît au fond, c'est la preuve que la solution est bien saturée : on est alors certain d'avoir atteint la concentration maximale possible.

\begin{exercice}[À toi de jouer]
\begin{enumerate}
\item On dissout \SI{50}{g} de sel dans \SI{200}{mL} d'eau à \SI{20}{\celsius}. La solution est-elle saturée ? Justifier avec un calcul ($s = \SI{360}{g/L}$).
\item À l'aide de la courbe, comparer les masses de \ce{KNO3} dissoutes dans \SI{1}{L} d'eau à \SI{20}{\celsius} puis à \SI{80}{\celsius}. Conclure.
\item Expliquer pourquoi, en cuisine, on chauffe l'eau pour préparer un sirop de sucre très concentré.
\end{enumerate}
\end{exercice}

\begin{corrige}[Exercice : À toi de jouer]
\begin{enumerate}
\item Masse maximale dissoute dans \SI{200}{mL} : $m_{\max} = 360 \times 0{,}200 = \SI{72}{g}$. Or on n'a dissous que \SI{50}{g}, ce qui est inférieur à \SI{72}{g} : la solution est \textbf{non saturée}.
\item À \SI{20}{\celsius} : environ \SI{320}{g} ; à \SI{80}{\celsius} : environ \SI{1400}{g}. La solubilité du \ce{KNO3} augmente fortement avec la température : on dissout plus de quatre fois plus de sel à chaud.
\item La solubilité du sucre augmente avec la température : en chauffant, on peut dissoudre une masse de sucre bien plus grande qu'à froid, puis en refroidissant on obtient un sirop très concentré.
\end{enumerate}
\end{corrige}

\begin{aretenir}[L'essentiel de la section]
\begin{itemize}
\item Une solution \cle{saturée} ne peut plus dissoudre de soluté à une température donnée ; l'excès forme un dépôt.
\item La \cle{solubilité} $s$ est la masse maximale de soluté dissous par litre de \textbf{solvant (eau)}, en \textbf{g/L}, à une température précise.
\item La solubilité des solides \textbf{augmente généralement avec la température}.
\item La masse maximale dissoute dans un volume $V$ de solvant se calcule par $m_{\max} = s \times V$ (avec $V$ en litres).
\item Ne jamais confondre solubilité (g/L) et concentration molaire (mol/L).
\end{itemize}
\end{aretenir}

Maintenant que tu maîtrises la notion de quantité maximale dissoute, il reste à apprendre à exprimer avec précision « combien » de soluté se trouve dans une solution donnée : c'est l'objet de la section suivante, consacrée à la \textbf{concentration molaire} $C = \dfrac{n}{V}$.

\coursec{Concentration molaire C = n/V}

Dans la section précédente, nous avons appris qu'un solvant ne peut dissoudre qu'une quantité limitée de soluté : c'est la \cle{solubilité}. Au-delà de cette limite, la solution est dite \cle{saturée}. Mais entre « presque rien dissous » et « saturation », il existe toute une gamme de solutions plus ou moins « chargées » en soluté. Comment dire, de façon précise, combien de soluté se trouve dans un litre de solution ? C'est exactement le rôle de la \cle{concentration molaire}, la grandeur que nous allons découvrir maintenant. Elle est indispensable : au laboratoire, à l'hôpital, dans l'industrie, on prépare chaque jour des solutions de concentration parfaitement connue.

\courssub{Pourquoi mesurer la concentration ?}

Imagine deux verres d'eau. Dans le premier, tu dissous une cuillère de sel ; dans le second, cinq cuillères. Les deux liquides sont limpides et se ressemblent, pourtant ils ne sont pas identiques : le second est beaucoup plus salé. Pour décrire cette différence, on dit que la deuxième solution est plus \cle{concentrée} que la première.

En chimie, « beaucoup » ou « peu » ne suffit pas. Le chimiste veut un \textbf{nombre} : la quantité de matière de soluté dissous dans un volume donné de solution. Ce nombre, c'est la concentration molaire.

\begin{exemplebox}[La potion du pharmacien]
Quand le pharmacien prépare un sirop ou quand l'infirmier prépare une perfusion à l'hôpital de district, il doit respecter exactement la concentration prescrite par le médecin. Trop concentrée, la solution devient dangereuse ; trop diluée, elle est inefficace. La concentration est donc une question de sécurité.
\end{exemplebox}

\courssub{Rappel : la quantité de matière n}

Avant de définir la concentration, rappelons une notion vue en début d'année : la \cle{quantité de matière}.

Compter les atomes ou les molécules un par un est impossible : ils sont beaucoup trop nombreux. Les chimistes les comptent donc par « paquets » appelés \cle{moles}. Une mole contient environ $6,022 \times 10^{23}$ particules : c'est le nombre d'Avogadro.

\begin{definition}[Quantité de matière et masse molaire]
La \cle{quantité de matière} $n$ d'une espèce chimique se mesure en \textbf{moles} (symbole : mol). La \cle{masse molaire} $M$ d'une espèce est la masse d'une mole de cette espèce ; elle s'exprime en grammes par mole (g/mol). On les relie par la relation :
\[ n = \frac{m}{M} \]
où $m$ est la masse de l'échantillon en grammes (g).
\end{definition}

\begin{exemplebox}[Exemples de masses molaires]
On calcule la masse molaire d'un composé en additionnant les masses molaires atomiques de tous les atomes de sa formule (données : $M(\ce{H}) = 1$ g/mol ; $M(\ce{C}) = 12$ g/mol ; $M(\ce{O}) = 16$ g/mol ; $M(\ce{Na}) = 23$ g/mol ; $M(\ce{Cl}) = 35,5$ g/mol).
\begin{itemize}
\item Chlorure de sodium (sel de cuisine) \ce{NaCl} : $M = 23 + 35,5 = 58,5$ g/mol ;
\item Glucose \ce{C6H12O6} : $M = 6 \times 12 + 12 \times 1 + 6 \times 16 = 180$ g/mol ;
\item Eau \ce{H2O} : $M = 2 \times 1 + 16 = 18$ g/mol.
\end{itemize}
\end{exemplebox}

\courssub{Définition de la concentration molaire}

\begin{experience}[Observer pour comprendre]
\textbf{Matériel :} deux fioles jaugées de 1 L, du sulfate de cuivre \ce{CuSO4} (poudre bleue), une balance, de l'eau distillée.

\textbf{Protocole :} dans la fiole A, on dissout 0,1 mol de sulfate de cuivre et on complète à 1 L. Dans la fiole B, on dissout 0,5 mol et on complète à 1 L.

\textbf{Observation :} la solution A est bleu clair, la solution B est bleu foncé. Plus il y a de soluté par litre, plus la couleur est intense.

\textbf{Interprétation :} la couleur traduit la concentration : la fiole B contient 5 fois plus de soluté par litre que la fiole A ($0,5 = 5 \times 0,1$). Sa concentration molaire est donc 5 fois plus grande.

\textbf{Conclusion :} pour comparer des solutions, on compare la quantité de matière de soluté rapportée au volume de solution : c'est la concentration molaire.
\end{experience}

\begin{definition}[Concentration molaire]
La \cle{concentration molaire} $C$ d'un soluté dans une solution est le quotient de la quantité de matière $n$ de soluté dissous par le volume $V$ de la solution :
\[ C = \frac{n}{V} \]
avec :
\begin{itemize}
\item $n$ en moles (mol) ;
\item $V$ en litres (L) ;
\item $C$ en moles par litre (mol/L).
\end{itemize}
De cette définition, on tire deux relations utiles : $n = C \times V$ et $V = \dfrac{n}{C}$.
\end{definition}

\begin{aretenir}[Les trois formules à connaître par cœur]
\begin{align*}
n &= \frac{m}{M} & C &= \frac{n}{V} & C_m &= \frac{m}{V}
\end{align*}
La première relie masse et quantité de matière ; la deuxième définit la concentration molaire ; la troisième définit la concentration massique (voir plus bas). Elles seront utilisées dans presque tous les exercices de chimie du BEPC.
\end{aretenir}

\courssub{Tableau récapitulatif des grandeurs}

\begin{center}
\begin{tabularx}{\linewidth}{|l|c|c|X|}
\hline
\rowcolor{popblueL} \textbf{Grandeur} & \textbf{Symbole} & \textbf{Unité} & \textbf{Signification} \\
\hline
Masse & $m$ & gramme (g) & Masse de soluté pesée sur la balance \\
\hline
Masse molaire & $M$ & g/mol & Masse d'une mole de l'espèce chimique \\
\hline
Quantité de matière & $n$ & mol & Nombre de « paquets » de particules \\
\hline
Volume de solution & $V$ & litre (L) & Volume total de la solution préparée \\
\hline
Concentration molaire & $C$ & mol/L & Moles de soluté par litre de solution \\
\hline
Concentration massique & $C_m$ & g/L & Grammes de soluté par litre de solution \\
\hline
\end{tabularx}
\end{center}

\courssub{Préparer une solution de concentration connue}

Au laboratoire, on ne prépare pas une solution « au hasard ». On suit un protocole précis, en trois étapes : la \textbf{pesée}, la \textbf{dissolution}, puis l'\textbf{ajustage du volume} dans une fiole jaugée.

\begin{popfigure}
\begin{center}
\begin{tikzpicture}[scale=0.95, every node/.style={font=\small}]
% Étape 1 : balance
\draw[fill=popline!40] (0,0) rectangle (2.4,0.5);
\draw[fill=popline!20] (0.7,0.5) rectangle (1.7,0.7);
\draw[fill=popcream] (0.55,0.7) .. controls (0.55,1.15) and (1.85,1.15) .. (1.85,0.7) -- cycle;
\draw[fill=popblue!60] (0.8,0.7) .. controls (0.8,1.0) and (1.6,1.0) .. (1.6,0.7) -- cycle;
\node at (1.2,-0.35) {\textbf{1. Pesée}};
\node[align=center] at (1.2,-0.95) {on pèse la masse $m$\\de soluté};
% flèche 1
\draw[-{Stealth[length=3mm]},thick,poporange] (2.6,0.6) -- (3.5,0.6);
% Étape 2 : bécher
\draw[thick] (4.0,0) -- (4.0,1.5) (5.8,0) -- (5.8,1.5);
\draw[thick] (4.0,0) -- (5.8,0);
\draw[fill=popblue!30] (4.02,0.02) rectangle (5.78,0.85);
\draw[fill=popblue!60] (4.6,0.85) circle (0.09);
\draw[fill=popblue!60] (5.1,0.85) circle (0.07);
\draw[fill=popblue!60] (4.85,0.85) circle (0.05);
\draw[-{Stealth[length=2.5mm]},popdark] (5.6,1.9) -- (5.3,1.0);
\node[popdark] at (6.35,2.0) {eau distillée};
\node at (4.9,-0.35) {\textbf{2. Dissolution}};
\node[align=center] at (4.9,-0.95) {on dissout le soluté\\dans un peu d'eau};
% flèche 2
\draw[-{Stealth[length=3mm]},thick,poporange] (6.6,0.6) -- (7.5,0.6);
% Étape 3 : fiole jaugée
\draw[thick] (8.4,0) -- (8.4,1.1) -- (8.75,1.7) -- (8.75,2.3);
\draw[thick] (10.0,0) -- (10.0,1.1) -- (9.65,1.7) -- (9.65,2.3);
\draw[thick] (8.4,0) -- (10.0,0);
\draw[fill=popblue!30] (8.42,0.02) -- (8.42,1.08) -- (8.76,1.66) -- (9.64,1.66) -- (9.98,1.08) -- (9.98,0.02) -- cycle;
\draw[thick,poporange] (8.62,1.45) -- (9.78,1.45);
\node[poporange,right] at (10.05,1.45) {trait de jauge};
\node at (9.2,-0.35) {\textbf{3. Ajustage}};
\node[align=center] at (9.2,-0.95) {on complète à 1 L\\dans la fiole jaugée};
\end{tikzpicture}
\end{center}
\legende{Préparation d'une solution de concentration connue : (1) on pèse la masse $m$ de soluté ; (2) on le dissout dans un peu d'eau distillée ; (3) on transvase dans une fiole jaugée et on complète avec de l'eau jusqu'au trait de jauge pour obtenir un volume $V$ exact.}
\end{popfigure}

\begin{attention}[Pourquoi la fiole jaugée et pas le bécher ?]
Le bécher et l'éprouvette donnent des volumes \emph{approximatifs}. La \cle{fiole jaugée} possède un trait de jauge gravé sur son col : quand le bas du ménisque du liquide affleure ce trait, le volume est exactement celui annoncé (100 mL, 250 mL, 500 mL, 1 L...). C'est elle qui garantit la précision de la concentration. On ne dissout jamais le soluté directement dans la fiole remplie à ras bord : on le dissout d'abord dans un peu d'eau, puis on complète.
\end{attention}

\begin{methode}[Calculer une concentration molaire en 3 étapes]
Pour calculer la concentration molaire $C$ d'une solution préparée en dissolvant une masse $m$ de soluté dans un volume $V$ de solution :
\begin{enumerate}
\item \textbf{Calculer la quantité de matière} : $n = \dfrac{m}{M}$ (avec $m$ en g et $M$ en g/mol) ;
\item \textbf{Convertir le volume en litres} si nécessaire : $V$ (L) $= V$ (mL) $\div 1000$ ;
\item \textbf{Appliquer la formule} : $C = \dfrac{n}{V}$ (avec $n$ en mol et $V$ en L). Le résultat s'exprime en mol/L.
\end{enumerate}
\end{methode}

\begin{exoresolu}[Préparation d'une solution de sel de cuisine]
\textbf{Énoncé.} Au laboratoire du collège, on dissout $m = 11,7$ g de chlorure de sodium \ce{NaCl} (masse molaire $M = 58,5$ g/mol) dans de l'eau, puis on ajuste le volume total de la solution à $V = 500$ mL. Calculer la concentration molaire $C$ de la solution obtenue.
\tcblower
\textbf{Solution détaillée.}

\textbf{Étape 1 : calcul de la quantité de matière.}
\[ n = \frac{m}{M} = \frac{11,7}{58,5} = 0,2 \text{ mol} \]

\textbf{Étape 2 : conversion du volume en litres.}
\[ V = 500 \text{ mL} = \frac{500}{1000} = 0,5 \text{ L} \]

\textbf{Étape 3 : calcul de la concentration molaire.}
\[ C = \frac{n}{V} = \frac{0,2}{0,5} = 0,4 \text{ mol/L} \]

\textbf{Conclusion :} la solution de chlorure de sodium a une concentration molaire $C = 0,4$ mol/L.
\end{exoresolu}

\begin{exoresolu}[Une solution de glucose]
\textbf{Énoncé.} Pour une séance de TP, on dissout 9 g de glucose \ce{C6H12O6} ($M = 180$ g/mol) dans une fiole jaugée et on complète à 250 mL. Quelle est la concentration molaire du glucose dans cette solution ?
\tcblower
\textbf{Solution détaillée.}

\textbf{Étape 1 :} quantité de matière de glucose :
\[ n = \frac{m}{M} = \frac{9}{180} = 0,05 \text{ mol} \]

\textbf{Étape 2 :} conversion du volume :
\[ V = 250 \text{ mL} = 0,25 \text{ L} \]

\textbf{Étape 3 :} concentration molaire :
\[ C = \frac{n}{V} = \frac{0,05}{0,25} = 0,2 \text{ mol/L} \]

\textbf{Conclusion :} la solution de glucose a une concentration $C = 0,2$ mol/L.
\end{exoresolu}

\begin{attention}[L'erreur n° 1 : oublier de convertir le volume en litres]
Dans la formule $C = \dfrac{n}{V}$, le volume $V$ doit \textbf{toujours} être exprimé en \textbf{litres}. Si l'énoncé donne $V = 500$ mL et que tu écris $C = \dfrac{0,2}{500}$, tu trouves 0,0004 mol/L : un résultat mille fois trop petit !

\textbf{Réflexe à acquérir :} dès que tu lis un volume en mL, convertis-le immédiatement : divise par 1000. Autre piège : ne pas confondre le volume d'eau ajoutée et le volume \emph{total} de la solution — c'est le volume total, mesuré dans la fiole jaugée, qui entre dans la formule.
\end{attention}

\courssub{Solution diluée et solution concentrée}

Ces deux mots du langage courant ont un sens précis en chimie : ils servent à \textbf{comparer} deux solutions d'un même soluté.

\begin{definition}[Diluée ou concentrée ?]
Entre deux solutions d'un même soluté, celle qui a la \textbf{plus grande} concentration molaire est dite plus \cle{concentrée} ; celle qui a la \textbf{plus petite} concentration est dite plus \cle{diluée}. \textbf{Diluer} une solution, c'est ajouter du solvant (souvent de l'eau) : le volume $V$ augmente, la quantité $n$ de soluté ne change pas, donc la concentration $C = n/V$ diminue.
\end{definition}

\begin{exemplebox}[Le jus de foléré]
Quand maman prépare du jus de bissap (foléré) trop fort, elle ajoute de l'eau : le goût devient plus léger. Elle vient de \emph{diluer} la solution. La quantité de « soluté » (les extraits de fleur, le sucre) n'a pas changé, mais le volume a augmenté, donc la concentration a diminué. C'est exactement la même opération au laboratoire !
\end{exemplebox}

\begin{attention}[Diluer n'est pas « enlever du soluté »]
Erreur fréquente au BEPC : croire qu'en ajoutant de l'eau, on « enlève » du soluté. Faux ! Lors d'une dilution, la quantité de matière $n$ de soluté reste \textbf{constante} ; seul le volume $V$ augmente. C'est pourquoi $C = n/V$ diminue. Retiens : \emph{diluer = ajouter du solvant, rien d'autre}.
\end{attention}

\courssub{Lien avec la concentration massique}

Il existe une autre façon d'exprimer la concentration : en grammes par litre au lieu de moles par litre.

\begin{definition}[Concentration massique]
La \cle{concentration massique} $C_m$ d'un soluté est le quotient de la masse $m$ de soluté par le volume $V$ de solution :
\[ C_m = \frac{m}{V} \]
avec $m$ en grammes (g), $V$ en litres (L) et $C_m$ en grammes par litre (g/L).
\end{definition}

\begin{propriete}[Lien entre $C$ et $C_m$]
Comme $n = \dfrac{m}{M}$, on a $m = n \times M$, donc :
\[ C_m = \frac{m}{V} = \frac{n \times M}{V} = C \times M \]
La concentration massique est égale à la concentration molaire multipliée par la masse molaire : $C_m = C \times M$. Cette relation sera très utilisée au lycée ; en 3ème, il suffit de savoir qu'elle existe et de l'appliquer dans des cas simples.
\end{propriete}

\begin{exemplebox}[Vérification avec notre solution de sel]
Pour la solution du premier exercice résolu : $C_m = C \times M = 0,4 \times 58,5 = 23,4$ g/L. Vérifions directement : $C_m = \dfrac{m}{V} = \dfrac{11,7}{0,5} = 23,4$ g/L. Les deux méthodes donnent le même résultat : la relation $C_m = C \times M$ est bien cohérente.
\end{exemplebox}

\courssub{Application : lire une étiquette}

Les concentrations sont partout dans la vie quotidienne : sachets de sérum, sirops, désinfectants, engrais liquides. Savoir lire une étiquette, c'est savoir interpréter une concentration.

\begin{exemplebox}[Le sérum physiologique]
Sur la poche de sérum physiologique utilisée à l'hôpital, on lit « \ce{NaCl} à 0,9 \% ». Cela signifie qu'il y a 0,9 g de chlorure de sodium pour 100 mL de solution, soit 9 g par litre. Sa concentration massique est donc $C_m = 9$ g/L. Calculons sa concentration molaire :
\[ C = \frac{C_m}{M} = \frac{9}{58,5} \approx 0,154 \text{ mol/L} \]
Cette valeur, proche de 0,15 mol/L, correspond à la concentration en sel du liquide qui baigne nos cellules.
\end{exemplebox}

\begin{savaistu}[Pourquoi le sérum physiologique est-il si bien toléré ?]
Le sang et le liquide qui entourent nos cellules contiennent des sels dissous à une concentration d'environ 0,15 mol/L. Le sérum physiologique a été réglé exactement à cette concentration : on dit qu'il est \emph{isotonique} par rapport au sang. Injecté en perfusion, il ne provoque ni gonflement ni rétraction des cellules, car l'eau n'a aucune raison d'entrer ou de sortir massivement. C'est pour cela qu'il est si bien toléré par l'organisme. Un sérum trop concentré (« hypertonique ») déshydraterait les cellules ; un sérum trop dilué (« hypotonique ») les ferait gonfler. La concentration n'est donc pas un détail : c'est une question de sécurité du patient !
\end{savaistu}

\begin{exercice}[À toi de jouer]
\textbf{Exercice 1.} On dissout 20 g d'hydroxyde de sodium \ce{NaOH} ($M = 40$ g/mol) dans de l'eau et on complète à 2 L. Calculer la concentration molaire de la solution.

\textbf{Exercice 2.} Une solution de sucre a une concentration $C = 0,5$ mol/L. Quelle quantité de matière de sucre contient un verre de 200 mL de cette solution ? (Indication : utilise $n = C \times V$.)

\textbf{Exercice 3.} Une infirmière doit préparer 500 mL d'une solution de glucose à 0,1 mol/L ($M = 180$ g/mol). Quelle masse de glucose doit-elle peser ?
\end{exercice}

\begin{corrige}[Exercice 1]
$n = \dfrac{m}{M} = \dfrac{20}{40} = 0,5$ mol ; $V = 2$ L ; $C = \dfrac{n}{V} = \dfrac{0,5}{2} = 0,25$ mol/L.
\end{corrige}

\begin{corrige}[Exercice 2]
Conversion : $V = 200$ mL $= 0,2$ L. Alors $n = C \times V = 0,5 \times 0,2 = 0,1$ mol. Le verre contient 0,1 mol de sucre.
\end{corrige}

\begin{corrige}[Exercice 3]
Quantité nécessaire : $n = C \times V = 0,1 \times 0,5 = 0,05$ mol. Masse à peser : $m = n \times M = 0,05 \times 180 = 9$ g. L'infirmière pèse 9 g de glucose, le dissout dans un peu d'eau distillée, puis complète à 500 mL dans une fiole jaugée.
\end{corrige}

\begin{aretenir}[L'essentiel de la section]
\begin{itemize}
\item La \cle{concentration molaire} mesure la quantité de soluté dissous par litre de solution : $C = \dfrac{n}{V}$, avec $n$ en mol, $V$ en \textbf{litres} et $C$ en mol/L.
\item Pour trouver $n$ à partir de la masse pesée : $n = \dfrac{m}{M}$. On peut aussi utiliser les formules inverses : $n = C \times V$ et $m = n \times M$.
\item La \cle{concentration massique} est $C_m = \dfrac{m}{V} = C \times M$ (g/L).
\item Diluer une solution, c'est ajouter du solvant : $n$ reste constante, $V$ augmente, donc $C$ diminue.
\item Une solution de concentration précise se prépare en fiole jaugée : pesée, dissolution, ajustage au trait de jauge.
\item La lecture des étiquettes (sérum physiologique à 0,9 \%, médicaments) repose sur ces notions : la concentration est une information de sécurité.
\end{itemize}
\end{aretenir}

\coursec{Leçon 6 : Solutions aqueuses}

\courssub{Dissolution, soluté, solvant, solution}

\begin{exemplebox}[Préparer un sirop de bissap]
Imagine que tu prépares du jus de bissap (feuilles d'hibiscus) pour la famille. Tu fais bouillir les fleurs dans l'eau : l'eau se colore en rouge foncé et prend le goût du fruit. Ensuite, tu ajoutes du sucre en poudre et tu remues : le sucre disparaît, le liquide reste limpide. Tu viens de réaliser une \textbf{dissolution}. Ce phénomène est partout : le café du matin, l'eau de mer à Kribi, le sérum physiologique à l'hôpital, l'essence dans le réservoir d'une moto-taxi (mélange d'hydrocarbures).
\end{exemplebox}

\begin{experience}[Observer une dissolution : le sucre dans l'eau]
\textbf{Matériel :} bécher, eau du robinet, sucre en poudre (saccharose), agitateur en verre, balance, spatule.

\textbf{Protocole :}
\begin{enumerate}
\item Verser \SI{100}{mL} d'eau dans le bécher.
\item Peser \SI{20}{g} de sucre.
\item Verser le sucre dans l'eau et remuer avec l'agitateur jusqu'à disparition complète des cristaux.
\item Observer l'aspect final du mélange.
\end{enumerate}

\textbf{Observations :} Au début, on voit les grains de sucre au fond du bécher. En remuant, ils deviennent invisibles. Le liquide final est \textbf{limpide} (transparent), \textbf{incolore} (ou légèrement coloré si l'eau l'était) et \textbf{homogène} (même aspect partout). Le volume final est légèrement supérieur à \SI{100}{mL} (environ \SI{115}{mL}), mais \textbf{la masse totale est conservée} (masse eau + masse sucre).

\textbf{Interprétation :} Les grains de sucre (solide) se sont dispersés au niveau microscopique en molécules individuelles \ce{C12H22O11} qui se sont insérées entre les molécules d'eau \ce{H2O}. Il n'y a plus de frontière visible entre le sucre et l'eau.

\textbf{Conclusion :} La dissolution est la dispersion d'une espèce chimique (le soluté) dans une autre (le solvant) pour former un mélange homogène appelé \textbf{solution}.
\end{experience}

\begin{definition}[Vocabulaire de la dissolution]
\begin{itemize}
\item \cle{Solvant} : l'espèce chimique \textbf{majoritaire} qui « accueille » l'autre. Dans une solution aqueuse, le solvant est l'\textbf{eau} (\ce{H2O}).
\item \cle{Soluté} : l'espèce chimique \textbf{minoritaire} qui est dispersée. Cela peut être un solide (sucre, sel), un gaz (\ce{CO2} dans les boissons gazeuses, \ce{O2} dans l'eau des rivières pour les poissons) ou un liquide (alcool dans le vin, eau de Javel).
\item \cle{Solution} : le mélange homogène obtenu (solvant + soluté).
\item \cle{Dissolution} : le phénomène physique (pas de réaction chimique, les espèces gardent leur identité) qui conduit à la solution.
\end{itemize}
\end{definition}

\begin{propriete}[Conservation de la masse lors d'une dissolution]
Lors d'une dissolution, la \textbf{masse de la solution} est égale à la \textbf{somme des masses} du solvant et du soluté :
\[ m_{\text{solution}} = m_{\text{solvant}} + m_{\text{soluté}} \]
Le volume, lui, \textbf{n'est pas conservé} (souvent $V_{\text{solution}} < V_{\text{solvant}} + V_{\text{soluté}}$ pour les solides).
\end{propriete}

\begin{attention}[Ne pas confondre dissolution et réaction chimique]
Quand tu verses de l'acide chlorhydrique sur du calcaire (\ce{CaCO3}), ça mousse : c'est une \textbf{réaction chimique} (formation de \ce{CO2}), pas une simple dissolution. Quand tu verses du sel dans l'eau, il n'y a pas de bulles, pas de chauffage spontané, le sel « disparaît » : c'est une \textbf{dissolution}. Au BEPC, on te demandera souvent de distinguer les deux.
\end{attention}

\begin{aretenir}[Critères d'une solution]
Une solution est un mélange \textbf{homogène} (une seule phase visible), \textbf{limpide} (pas de trouble, ne diffuse pas la lumière — pas d'effet Tyndall), et \textbf{stable} (le soluté ne se dépose pas au fond au repos).
\end{aretenir}

\courssub{Solubilité et saturation}

\begin{exemplebox}[Combien de sucre dans ton thé ?]
Tu suces ton thé chaud : tu peux y dissoudre beaucoup de sucre (jusqu'à \SI{200}{g} par \SI{100}{mL} d'eau à \SI{20}{\celsius} !). Mais si tu essaies de dissoudre du sel dans de l'eau froide, tu bloques vers \SI{36}{g} par \SI{100}{mL}. Et l'huile dans l'eau ? \textbf{Zéro} : ça ne se mélange pas, on dit qu'elles sont \textbf{immiscibles}. Pourquoi ces différences ?
\end{exemplebox}

\begin{experience}[Mettre en évidence la saturation (sel dans l'eau)]
\textbf{Matériel :} bécher, eau à \SI{20}{\celsius}, sel fin (NaCl), balance, spatule, agitateur, thermomètre.

\textbf{Protocole :}
\begin{enumerate}
\item Verser \SI{100}{mL} d'eau dans le bécher, noter la température.
\item Ajouter \SI{5}{g} de sel, remuer jusqu'à dissolution complète.
\item Répéter l'opération par portions de \SI{5}{g} en notant la masse totale ajoutée.
\item S'arrêter quand il reste du sel non dissous au fond malgré l'agitation prolongée.
\item Filtrer la solution et peser la masse de sel dissous (masse initiale - masse résidu sec).
\end{enumerate}

\textbf{Observations :} Tant qu'on n'a pas dépassé environ \SI{36}{g}, tout se dissout. Au-delà, l'excédent reste au fond. La solution obtenue est \textbf{saturée}. Si on chauffe le bécher, le sel restant se dissout : la solubilité augmente avec la température.

\textbf{Interprétation :} Il existe une \textbf{limite} à la quantité de soluté qu'un solvant peut dissoudre à une température donnée. Cette limite est la \textbf{solubilité}.
\end{experience}

\begin{definition}[Solubilité et solution saturée]
La \cle{solubilité} d'un soluté dans un solvant à une température donnée est la \textbf{masse maximale} de ce soluté (en \SI{}{g} ou \SI{}{g/L}) qui peut se dissoudre dans une quantité donnée de solvant (souvent \SI{100}{g} ou \SI{1}{L}) pour former une \cle{solution saturée}. Au-delà, l'excédent de soluté reste non dissous (dépôt au fond).
\begin{itemize}
\item Solution \textbf{non saturée} : $m_{\text{soluté dissous}} < \text{solubilité}$. Tout le soluté a disparu.
\item Solution \textbf{saturée} : $m_{\text{soluté dissous}} = \text{solubilité}$. Équilibre entre soluté dissous et soluté solide en excès.
\item Solution \textbf{supersaturée} (instable) : obtenue en refroidissant une solution saturée chaude sans agitation ; le moindre choc déclenche la cristallisation.
\end{itemize}
\end{definition}

\begin{propriete}[Facteurs influençant la solubilité des solides]
\begin{itemize}
\item \textbf{Nature du soluté et du solvant :} « Qui se ressemble s'assemble ». L'eau (polaire) dissout les solides ioniques (sel) et polaires (sucre), mais pas l'huile (apolaire).
\item \textbf{Température :} Pour \textbf{la quasi-totalité des solides}, la solubilité \textbf{augmente} avec la température. (Exception rare : quelques sels comme le sulfate de cérium).
\item \textbf{Pression :} N'influence \textbf{pas} la solubilité des solides (contrairement aux gaz).
\end{itemize}
\end{propriete}

\begin{popfigure}
\begin{tikzpicture}[scale=0.9, every node/.style={font=\small}]
\begin{axis}[
    width=\linewidth, height=7cm,
    xlabel={Température (\si{\celsius})}, ylabel={Solubilité (\si{g/100g\_eau})},
    xmin=0, xmax=100, ymin=0, ymax=250,
    grid=major, grid style={popline},
    axis lines=middle, enlargelimits={abs=5},
    tick style={popdark}, label style={popdark},
    legend style={at={(0.02,0.98)}, anchor=north west, fill=popcream, draw=popmuted}
]
\addplot[popcurve, domain=0:100, samples=50, thick] {35.7 + 0.04*x + 0.001*x^2}; % NaCl approx
\addlegendentry{NaCl (sel) : faible variation}
\addplot[poporange, domain=0:100, samples=50, thick] {179 + 2.1*x + 0.03*x^2}; % Sucrose approx
\addlegendentry{Saccharose (sucre) : forte augmentation}
\addplot[popgreen, domain=0:100, samples=50, thick, dashed] {139 - 0.2*x}; % Ce2(SO4)3 approx
\addlegendentry{Ce$_2$(SO$_4$)$_3$ : exception (diminue)}
\end{axis}
\end{tikzpicture}
\legende{Évolution qualitative de la solubilité de quelques solides dans l'eau en fonction de la température. Le sucre (courbe orange) voit sa solubilité fortement augmenter, le sel (courbe bleue) varie peu, le sulfate de cérium (pointillés verts) diminue.}
\end{popfigure}

\begin{exemplebox}[Applications au Cameroun]
\begin{itemize}
\item \textbf{Production de sel marin (Kribi, Douala) :} L'eau de mer est canalisée dans des bassins d'évaporation. Sous l'effet du soleil, l'eau s'évapore, la solution devient saturée puis supersaturée, et le sel \textbf{cristallise}. On récupère le sel par récolte mécanique.
\item \textbf{Boissons gazeuses (usines de Yaoundé/Douala) :} On dissout du \ce{CO2} dans l'eau sous \textbf{haute pression} (loi de Henry : solubilité des gaz $\propto$ pression). À l'ouverture, la pression chute, le gaz devient supersaturé et s'échappe (bulles).
\item \textbf{Eau de rivière et poissons :} L'\ce{O2} dissous (solubilité $\approx \SI{9}{mg/L}$ à \SI{20}{\celsius}) permet la respiration des poissons. L'eau chaude (rejet industriel, déforestation) contient moins d'\ce{O2} $\rightarrow$ asphyxie (mortalité piscicole).
\end{itemize}
\end{exemplebox}

\begin{methode}[Préparer une solution saturée]
\begin{enumerate}
\item Choisir le soluté et le solvant (ex: sel + eau).
\item Verser un volume connu de solvant dans un bécher.
\item Ajouter le soluté par petites quantités en agitant énergiquement entre chaque ajout.
\item S'arrêter dès qu'un dépôt persistant apparaît au fond malgré l'agitation.
\item (Optionnel) Filtrer pour récupérer la solution saturée sans excédent solide.
\end{enumerate}
\end{methode}

\begin{attention}[Pièges BEPC sur la solubilité]
\begin{itemize}
\item Confondre \textbf{solubilité} (valeur maximale, en g/100g ou g/L) et \textbf{concentration} (valeur réelle d'une solution donnée, en mol/L ou g/L). Une solution non saturée a une concentration \textbf{inférieure} à la solubilité.
\item Oublier de préciser la \textbf{température} quand on donne une solubilité. « La solubilité du sel est 36 g/100g » est incomplet ; il faut dire « à \SI{20}{\celsius} ».
\item Croire que « saturé » veut dire « concentré ». Une solution saturée en sel est concentrée (\SI{360}{g/L}), mais une solution saturée en \ce{CaCO3} (calcaire) est très peu concentrée (\SI{0.013}{g/L}).
\end{itemize}
\end{attention}

\courssub{Concentration molaire $C = n/V$}

\begin{exemplebox}[Le dosage du sérum physiologique]
À l'hôpital, on te met une perfusion de « sérum physiologique » : c'est de l'eau avec du sel (\ce{NaCl}) à \SI{0,9}{\%} en masse, soit \SI{9}{g/L}. Pourquoi cette valeur précise ? Parce qu'elle correspond à une concentration molaire d'environ \SI{0,15}{mol/L}, isotone avec le sang (même pression osmotique). Si l'infirmière se trompe et met \SI{20}{g/L}, tes globules rouges éclatent (hémolyse) ; si elle met \SI{2}{g/L}, ils gonflent et éclatent aussi. La \textbf{concentration molaire} est une question de vie ou de mort.
\end{exemplebox}

\begin{definition}[Quantité de matière $n$ (rappel)]
La \cle{quantité de matière} $n$ (unité : \textbf{mole}, symbole \textbf{mol}) compte le nombre d'entités (atomes, molécules, ions) dans un échantillon.
\[ n = \frac{m}{M} \]
avec $m$ la masse en \textbf{gramme (g)} et $M$ la \textbf{masse molaire} en \textbf{g/mol} (valeur numérique de la masse atomique/moléculaire trouvée dans le tableau périodique).
\end{definition}

\begin{definition}[Concentration molaire $C$]
La \cle{concentration molaire} (ou concentration en quantité de matière) d'une solution est le quotient de la quantité de matière de soluté $n$ (en \textbf{mol}) par le volume de la solution $V$ (en \textbf{litre, L}) :
\[ C = \frac{n}{V} \qquad \text{Unité : \textbf{mol/L} (ou \textbf{mol$\cdot$L$^{-1}$})} \]
On note souvent $C(\text{espèce})$ ou $[\text{espèce}]$.
\end{definition}

\begin{methode}[Calculer une concentration molaire / Préparer une solution]
\textbf{Données :} masse $m$ de soluté, masse molaire $M$, volume final $V$ de solution.
\begin{enumerate}
\item Calculer la quantité de matière : $n = \dfrac{m}{M}$ (en mol).
\item Calculer la concentration : $C = \dfrac{n}{V}$ (en mol/L).
\end{enumerate}
\textbf{Inversement (préparation) :} Pour préparer $V$ litres de solution à $C$ mol/L :
\begin{enumerate}
\item Calculer $n = C \times V$ (mol).
\item Calculer $m = n \times M$ (g).
\item Peser $m$ grammes de soluté (balance de précision).
\item Dissoudre dans un bécher avec un peu de solvant.
\item Transférer \textbf{quantitativement} (tout rincer) dans une \textbf{fiole jaugée} de volume $V$.
\item Compléter au trait de jauge avec le solvant (pipette ou burette pour la précision), boucher, homogénéiser.
\end{enumerate}
\end{methode}

\begin{exemplebox}[Préparation de \SI{500}{mL} de solution de NaCl à \SI{0,10}{mol/L}]
$M(\ce{NaCl}) = 23,0 + 35,5 = \SI{58,5}{g/mol}$.
\begin{enumerate}
\item $n = C \times V = 0,10 \times 0,500 = \SI{0,050}{mol}$.
\item $m = n \times M = 0,050 \times 58,5 = \SI{2,925}{g} \approx \SI{2,93}{g}$.
\item Peser \SI{2,93}{g} de NaCl, dissoudre, transférer en fiole jaugée \SI{500}{mL}, compléter à l'eau distillée.
\end{enumerate}
\end{exemplebox}

\begin{propriete}[Dilution : conservation de la quantité de matière]
Lors d'une dilution (ajout de solvant), la quantité de matière de soluté \textbf{ne change pas} : $n_1 = n_2$.
Donc $C_1 V_1 = C_2 V_2$ (avec $V$ en L).
C'est la relation fondamentale pour préparer une solution fille à partir d'une solution mère plus concentrée.
\end{propriete}

\begin{exoresolu}[Dilution d'une solution mère d'acide chlorhydrique]
Énoncé : On dispose d'une solution mère d'acide chlorhydrique $S_0$ de concentration $C_0 = \SI{1,0}{mol/L}$. On veut préparer \SI{100}{mL} d'une solution fille $S_1$ de concentration $C_1 = \SI{0,10}{mol/L}$.
\begin{enumerate}
\item Calculer le volume $V_0$ de solution mère à prélever.
\item Décrire la manipulation en toute sécurité.
\end{enumerate}
\tcblower
Solution détaillée :
\begin{enumerate}
\item Relation de dilution : $C_0 V_0 = C_1 V_1$.
$V_0 = \dfrac{C_1 V_1}{C_0} = \dfrac{0,10 \times 0,100}{1,0} = \SI{0,010}{L} = \SI{10,0}{mL}$.
\item \textbf{Sécurité :} Acide corrosif $\rightarrow$ gants, lunettes, blouse. \textbf{Matériel :} pipette jaugée \SI{10}{mL} (ou graduée), propipette (jamais la bouche !), fiole jaugée \SI{100}{mL}, bécher, eau distillée. \textbf{Geste :} Prélever $V_0$ avec la pipette, verser dans la fiole jaugée, rincer la pipette dans la fiole, compléter à l'eau distillée au trait de jauge, homogénéiser. \textbf{Règle d'or :} On verse \textbf{l'acide dans l'eau} (jamais l'inverse : projection dangereuse par ébullition locale).
\end{enumerate}
\end{exoresolu}

\begin{attention}[Erreurs fréquentes sur $C = n/V$]
\begin{itemize}
\item Unités : $V$ \textbf{doit être en litres (L)}. Si $V = \SI{250}{mL}$, il faut écrire $V = 0,250$ L dans la formule.
\item Confondre $V_{\text{solvant}}$ et $V_{\text{solution}}$. La formule utilise le volume \textbf{final de la solution} (trait de jauge de la fiole), pas le volume d'eau versé au début.
\item Oublier la conversion $n = m/M$ si l'énoncé donne une masse.
\end{itemize}
\end{attention}

\courssub{Notion de pH : solutions acides, neutres, basiques}

Dans la section précédente, nous avons appris à caractériser une solution par sa concentration molaire $C = \dfrac{n}{V}$, c'est-à-dire par la \emph{quantité} de soluté dissous par litre de solution. Mais la concentration ne dit pas tout ! Deux solutions peuvent avoir la même concentration et pourtant des effets très différents : le jus de citron pique la langue, l'eau pure est sans saveur particulière, et l'eau savonneuse est glissante et amère. Il existe une autre grandeur qui décrit le \emph{caractère} d'une solution aqueuse : c'est le \cle{pH}. C'est lui qui permet de distinguer une solution acide d'une solution basique. Cette notion est indispensable dans la vie courante : contrôle de l'eau de boisson (Camerounaise des Eaux), choix des engrais selon le sol (cacao, café), soins médicaux, entretien des batteries de moto-taxi.

\courssub{Expérience : tester des solutions usuelles avec du papier pH}

\begin{experience}[Tester des solutions de la vie courante]
\textbf{Matériel :} papier pH (bandelettes), échelle de couleurs de référence, agitateur en verre, quatre béchers, jus de citron, eau pure (eau distillée), solution de savon, eau de Javel diluée, gants et lunettes de protection.

\textbf{Protocole :}
\begin{enumerate}
\item Verser un peu de chaque solution dans un bécher étiqueté.
\item Prélever une goutte de la solution avec l'agitateur en verre propre.
\item Déposer la goutte sur un morceau de papier pH. \textbf{Ne jamais tremper le papier dans le flacon} (cela polluerait toute la solution).
\item Observer la couleur prise par le papier et la comparer à l'échelle de couleurs.
\item Rincer l'agitateur à l'eau distillée entre chaque mesure.
\end{enumerate}

\textbf{Observations :}
\begin{center}
\begin{tabularx}{\linewidth}{|l|X|c|}\hline
\rowcolor{popblueL} \textbf{Solution testée} & \textbf{Couleur du papier pH} & \textbf{pH lu} \\ \hline
Jus de citron & Rouge-orangé & $\approx 2$ \\ \hline
Eau pure & Vert & $= 7$ \\ \hline
Solution de savon & Bleu & $\approx 10$ \\ \hline
Eau de Javel diluée & Bleu foncé-violet & $\approx 12$ \\ \hline
\end{tabularx}
\end{center}

\textbf{Interprétation :} le papier pH change de couleur selon la solution testée. Le jus de citron donne une couleur rouge-orangée (petit pH), l'eau pure donne du vert (pH = 7), le savon et l'eau de Javel donnent des couleurs bleues (grand pH). Le pH semble donc lié au caractère acide ou non de la solution.

\textbf{Conclusion :} chaque solution aqueuse possède un pH que l'on peut mesurer simplement avec du papier pH. Les solutions acides (citron) ont un petit pH, les solutions comme le savon ont un grand pH, et l'eau pure se situe au milieu de l'échelle.
\end{experience}

\begin{popfigure}
\begin{tikzpicture}[scale=1, every node/.style={font=\small}]
% Flacon et agitateur
\draw[thick, popblue] (0,0) -- (0,1.8) -- (1.6,1.8) -- (1.6,0) -- cycle;
\fill[popblue!25] (0.05,0.05) rectangle (1.55,0.9);
\draw[thick, popblue] (0.5,1.8) -- (0.5,2.1) -- (1.1,2.1) -- (1.1,1.8);
\node[below, popdark] at (0.8,-0.15) {Solution à tester};
% Agitateur
\draw[thick, poporange] (0.8,3.1) -- (0.8,0.6);
\fill[poporange] (0.8,0.55) circle (0.06);
\node[right, popdark] at (0.9,2.9) {Agitateur en verre};
% Goutte
\fill[popblue] (3.2,1.6) circle (0.08);
\draw[->, thick, popmuted] (1.0,0.9) .. controls (2.0,1.5) .. (3.1,1.7);
\node[above, popdark] at (2.1,1.6) {1 goutte};
% Bandelette de papier pH
\fill[popgold!50] (2.4,0.6) rectangle (6.4,1.1);
\draw[thick, popgold!80!black] (2.4,0.6) rectangle (6.4,1.1);
\fill[popgreen] (3.0,0.6) rectangle (3.5,1.1);
\node[below, popdark] at (4.4,0.45) {Bandelette de papier pH};
\node[above, popdark] at (3.25,1.15) {zone mouillée};
% Echelle de couleurs
\foreach \i/\c in {0/poppink,1/poppink!70!poporange,2/poporange,3/poporange!60!popgold,4/popgold,5/popgold!70!popgreen,6/popgreen!70!popgold,7/popgreen,8/popgreen!60!popblue,9/popblue!60!popgreen,10/popblue,11/popblue!75!poppurple,12/poppurple,13/poppurple!75!popdark,14/popdark}{
  \fill[\c] (2.4+0.2857*\i,2.6) rectangle (2.6857+0.2857*\i,3.0);
}
\draw[thick, popdark] (2.4,2.6) rectangle (6.4,3.0);
\node[above, popdark] at (4.4,3.05) {Échelle de couleurs de référence};
% Interdiction
\draw[thick, poppink] (7.6,0.4) -- (7.6,1.9) -- (8.6,1.9) -- (8.6,0.4) -- cycle;
\fill[popblue!25] (7.65,0.45) rectangle (8.55,1.1);
\fill[popgold!50] (7.8,1.2) rectangle (8.4,1.7);
\draw[very thick, poppink] (7.4,0.2) -- (8.8,2.1);
\draw[very thick, poppink] (8.1,1.15) circle (1.05);
\node[below, poppink, align=center] at (8.1,0.0) {INTERDIT : tremper le\\papier dans le flacon};
\end{tikzpicture}
\legende{Le bon geste de mesure au papier pH : on prélève une goutte de solution avec un agitateur en verre propre et on la dépose sur la bandelette, puis on compare la couleur obtenue à l'échelle de référence. On ne trempe jamais le papier dans le flacon, pour ne pas contaminer la solution.}
\end{popfigure}

\courssub{Définition du pH}

L'expérience précédente montre qu'une seule grandeur, lue sur une échelle de couleurs, suffit à classer nos quatre solutions. Cette grandeur porte un nom : le pH.

\begin{definition}[pH d'une solution aqueuse]
Le \cle{pH} (potentiel hydrogène) d'une solution aqueuse est une \textbf{grandeur sans unité} qui indique le \textbf{caractère acide, neutre ou basique} de cette solution. Il se lit sur une échelle graduée de \textbf{0 à 14}.
\end{definition}

\begin{definition}[Solutions acides, neutres et basiques]
\begin{itemize}
\item Une solution est \cle{acide} si son pH est \textbf{inférieur à 7} ($\text{pH} < 7$).
\item Une solution est \cle{neutre} si son pH est \textbf{égal à 7} ($\text{pH} = 7$).
\item Une solution est \cle{basique} (ou alcaline) si son pH est \textbf{supérieur à 7} ($\text{pH} > 7$).
\end{itemize}
\end{definition}

\begin{propriete}[Propriété de l'échelle de pH (admise)]
Toute solution aqueuse a un pH compris entre \textbf{0 et 14}. De plus :
\begin{itemize}
\item plus le pH est \textbf{faible} (proche de 0), plus la solution est \textbf{acide} ;
\item plus le pH est \textbf{élevé} (proche de 14), plus la solution est \textbf{basique} ;
\item le pH = 7 correspond à une solution \textbf{neutre} (comme l'eau pure à \SI{25}{\celsius}).
\end{itemize}
Cette propriété est admise en classe de 3ème : aucune démonstration n'est exigible au BEPC.
\end{propriete}

\begin{attention}[Piège classique : pH = 0 ne veut pas dire « pas d'acidité »]
Beaucoup d'élèves croient qu'un pH égal à 0 signifie « absence d'acidité », comme une note de zéro. C'est l'inverse ! Un pH = 0 correspond à une \textbf{acidité maximale}. Retiens bien le sens de l'échelle : \textbf{petit pH = très acide}, \textbf{grand pH = très basique}. De même, une solution de pH = 6,5 est déjà acide (faiblement), car son pH est inférieur à 7.
\end{attention}

\begin{popfigure}
\begin{tikzpicture}[scale=1]
% Dégradé de l'échelle
\foreach \i/\c in {0/poppink,1/poppink!70!poporange,2/poporange,3/poporange!60!popgold,4/popgold,5/popgold!70!popgreen,6/popgreen!70!popgold,7/popgreen,8/popgreen!60!popblue,9/popblue!60!popgreen,10/popblue,11/popblue!75!poppurple,12/poppurple,13/poppurple!75!popdark,14/popdark}{
  \fill[\c] ({0.85*\i},0) rectangle ({0.85*\i+0.85},0.7);
}
\draw[thick, popdark] (0,0) rectangle (11.9,0.7);
% Graduations
\foreach \i in {0,...,14}{
  \draw[thick, popdark] ({0.85*\i},0) -- ({0.85*\i},-0.15);
  \node[below, popdark, font=\small] at ({0.85*\i},-0.2) {\i};
}
% Zones
\draw[decorate, decoration={brace, amplitude=6pt}, poppink, thick] (0,0.85) -- (5.95,0.85);
\node[above, poppink, font=\bfseries] at (2.975,1.15) {ACIDE (pH $<$ 7)};
\draw[decorate, decoration={brace, amplitude=6pt}, popgreen!60!black, thick] (5.95,0.85) -- (6.8,0.85);
\node[above, popgreen!50!black, font=\bfseries] at (6.375,1.15) {Neutre};
\draw[decorate, decoration={brace, amplitude=6pt}, popblue, thick] (6.8,0.85) -- (11.9,0.85);
\node[above, popblue, font=\bfseries] at (9.35,1.15) {BASIQUE (pH $>$ 7)};
% Exemples positionnés
\draw[->, thick, popdark] (1.7,-0.9) -- (1.7,-0.45);
\node[below, popdark, font=\small] at (1.7,-0.95) {Citron ($\approx 2$)};
\draw[->, thick, popdark] (3.4,-1.7) -- (3.4,-0.45);
\node[below, popdark, font=\small] at (3.4,-1.75) {Vinaigre ($\approx 3$)};
\draw[->, thick, popdark] (5.95,-0.9) -- (5.95,-0.45);
\node[below, popdark, font=\small] at (5.95,-0.95) {Eau pure ($= 7$)};
\draw[->, thick, popdark] (8.5,-0.9) -- (8.5,-0.45);
\node[below, popdark, font=\small] at (8.5,-0.95) {Savon ($\approx 10$)};
\draw[->, thick, popdark] (10.2,-1.7) -- (10.2,-0.45);
\node[below, popdark, font=\small] at (10.2,-1.75) {Eau de Javel ($\approx 12$)};
% Flèches de tendance
\draw[->, very thick, poppink] (5.5,-2.6) -- (0.5,-2.6);
\node[left, poppink, font=\small\bfseries] at (0.4,-2.6) {Acidité croissante};
\draw[->, very thick, popblue] (6.3,-2.6) -- (11.4,-2.6);
\node[right, popblue, font=\small\bfseries] at (11.5,-2.6) {Basicité croissante};
\end{tikzpicture}
\legende{L'échelle de pH, de 0 à 14. En dessous de 7, la solution est acide (rouge-orange) ; à 7, elle est neutre (vert) ; au-dessus de 7, elle est basique (bleu-violet). Quelques solutions usuelles sont positionnées à titre d'exemple.}
\end{popfigure}

\courssub{Mesure du pH}

\begin{methode}[Mesurer un pH avec du papier pH, en toute sécurité]
\begin{enumerate}
\item Mettre des \textbf{gants} et des \textbf{lunettes de protection}, surtout avec des produits comme l'eau de Javel (irritant, corrosif).
\item Verser un petit volume de la solution dans un bécher propre.
\item Prélever \textbf{une goutte} de solution avec un agitateur en verre propre et sec.
\item Déposer la goutte sur une bandelette de papier pH posée sur une surface propre (tuile, verre de montre). \textbf{Ne jamais plonger le papier dans le flacon} : cela contaminerait toute la solution.
\item Comparer \textbf{immédiatement} la couleur obtenue à l'échelle de référence et lire la valeur du pH.
\item Rincer l'agitateur à l'eau distillée avant de tester une autre solution, et se laver les mains à la fin.
\end{enumerate}
\end{methode}

\begin{aretenir}[Instruments de mesure du pH]
\begin{itemize}
\item Le \cle{papier pH} (ou bandelettes) : méthode \textbf{simple, rapide, peu coûteuse}, mais peu précise (lecture à l'unité près, $\pm 1$). Idéal pour le terrain (agriculture, aquarium, piscine).
\item Le \cle{pH-mètre} : instrument électronique de \textbf{précision}, muni d'une sonde (électrode de verre) que l'on plonge dans la solution ; il affiche le pH au dixième près ($\pm 0,1$). Nécessite étalonnage (solutions tampons pH 4, 7, 10). Utilisé dans les laboratoires (SONARA, CDE, hôpitaux, industries agroalimentaires).
\item \textbf{Indicateurs colorés} (phénolphtaléine, bleu de bromothymol, rouge de méthyle, chou rouge) : utiles en titrage (dosage) pour détecter l'équivalence (changement de couleur brutal).
\end{itemize}
\end{aretenir}

\begin{exemplebox}[Quelques valeurs de pH usuelles (contexte camerounais)]
\begin{center}
\begin{tabularx}{\linewidth}{|l|c|X|}\hline
\rowcolor{popblueL} \textbf{Solution} & \textbf{pH} & \textbf{Caractère / Contexte} \\ \hline
Acide batterie (moto-taxi) & $\approx 0,5$ & Très acide (H$_2$SO$_4$ concentré) \\ \hline
Jus de citron / Bissap & $\approx 2$--$3$ & Acide (vitamine C, acides organiques) \\ \hline
Vinaigre & $\approx 3$ & Acide (acide acétique) \\ \hline
Eau de pluie (ville, pollution) & $\approx 5$--$6$ & Faiblement acide (CO$_2$, SO$_2$, NO$_x$) \\ \hline
Eau pure (distillée) & $= 7$ & Neutre (à \SI{25}{\celsius}) \\ \hline
Sang humain & $\approx 7,4$ & Très légèrement basique (tampon bicarbonate) \\ \hline
Eau de mer (Kribi) & $\approx 8,2$ & Légèrement basique (carbonates) \\ \hline
Solution de savon / Lessive & $\approx 10$ & Basique (carbonates, phosphates) \\ \hline
Eau de Javel (concentrée) & $\approx 12$--$13$ & Très basique (hypochlorite NaClO + soude) \\ \hline
Soude caustique (déboucheur) & $\approx 14$ & Extrêmement basique (NaOH) \\ \hline
\end{tabularx}
\end{center}
\end{exemplebox}

\begin{exoresolu}[Classer des solutions d'après leur pH]
Énoncé : on mesure au papier pH trois solutions de la maison. On obtient : solution A (jus d'orange) : pH = 4 ; solution B (eau minérale) : pH = 7 ; solution C (lessive) : pH = 11.
\begin{enumerate}
\item Indiquer le caractère (acide, neutre ou basique) de chaque solution.
\item Classer ces solutions de la plus acide à la plus basique.
\end{enumerate}
\tcblower
Solution détaillée :
\begin{enumerate}
\item On compare chaque pH à 7 :
\begin{itemize}
\item Solution A : pH = 4 $<$ 7, donc la solution A est \textbf{acide}.
\item Solution B : pH = 7, donc la solution B est \textbf{neutre}.
\item Solution C : pH = 11 $>$ 7, donc la solution C est \textbf{basique}.
\end{itemize}
\item Plus le pH est faible, plus la solution est acide ; plus il est élevé, plus elle est basique. Le classement de la plus acide à la plus basique est donc :
\[ \text{A (pH = 4)} \;<\; \text{B (pH = 7)} \;<\; \text{C (pH = 11)} \]
c'est-à-dire : jus d'orange, puis eau minérale, puis lessive.
\end{enumerate}
\end{exoresolu}

\begin{attention}[Erreurs fréquentes au BEPC]
\begin{itemize}
\item Écrire que le pH a une unité : c'est faux, le pH est une grandeur \textbf{sans unité}. On écrit « pH = 7 » et jamais « pH = 7 g/L » ou « pH = 7 mol/L ».
\item Dire qu'une solution de pH = 8 est acide « parce que c'est proche de 7 » : non, dès que le pH dépasse 7, la solution est \textbf{basique}.
\item Tremper la bandelette dans le flacon : geste interdit, il fausse la mesure et salit le produit.
\item Confondre pH et concentration : une solution peut être très concentrée et neutre (eau très salée \ce{NaCl} saturée, pH = 7), ou très diluée et acide (acide chlorhydrique \SI{10^{-5}}{mol/L}, pH = 5).
\end{itemize}
\end{attention}

\courssub{Importance pratique du pH}

Le pH n'est pas une notion abstraite de laboratoire : il intervient chaque jour dans notre vie.

\begin{itemize}
\item \textbf{Eau de boisson (Camerounaise des Eaux - CDE) :} L'eau potable distribuée doit avoir un pH proche de 7, en général entre 6,5 et 8,5 (norme OMS). Une eau trop acide (pH < 6,5) attaque les canalisations métalliques (corrosion, goût métallique, plomb/cuivre dans l'eau). Une eau trop basique (pH > 8,5) a mauvais goût (savonneuse), précipite le calcaire et diminue l'efficacité du chlore (désinfection). Les techniciens mesurent le pH régulièrement avec un pH-mètre aux stations de traitement et sur le réseau.
\item \textbf{Agriculture (cacao, café, bananier) :} Le pH du sol conditionne la disponibilité des éléments nutritifs (azote, phosphore, potassium). Le cacao, culture majeure d'exportation au Cameroun, pousse bien dans des sols légèrement acides (pH entre 5,0 et 6,5). Si le sol est trop acide (pH < 4,5, fréquent sous forêt après déboisement), l'aluminium devient toxique pour les racines. L'agriculteur épand alors de la \textbf{chaux} (carbonate de calcium \ce{CaCO3}, basique) pour faire remonter le pH : c'est le \textbf{chaulage}.
\item \textbf{Santé (estomac, médicaments) :} L'estomac sécrète de l'acide chlorhydrique (pH $\approx$ 1,5--2) pour digérer les protéines et tuer les bactéries. Quand cette acidité remonte vers l'œsophage (reflux), on ressent des brûlures (pyrosis). On prend un \textbf{antiacide} (souvent hydrogénocarbonate de sodium \ce{NaHCO3} ou hydroxyde d'aluminium \ce{Al(OH)3}), de nature basique, qui neutralise l'excès d'acidité : $\ce{H+ + HCO3- -> H2O + CO2}\,\uparrow$.
\item \textbf{Hygiène (savon, peau) :} Les savons traditionnels sont basiques (pH $\approx$ 10), ce qui leur permet de saponifier (dissoudre) les graisses. Mais un savon trop basique irrite la peau, dont le pH naturel (film hydrolipidique) est voisin de 5,5 (légèrement acide pour freiner les bactéries). Les « savons sans savon » (syndets) ont un pH neutre ou physiologique (5,5).
\item \textbf{Industrie (SONARA, brasseries, savonneries) :} Contrôle du pH des effluents avant rejet en milieu naturel (loi sur l'environnement), contrôle du pH lors du raffinage du pétrole (traitement des coupes acides), brassage (pH du moût $\approx$ 5,2--5,6 pour l'activité enzymatique).
\end{itemize}

\begin{savaistu}[Le pH du sang : une question de vie ou de mort]
Le sang humain a un pH maintenu très précisément autour de \textbf{7,40} (intervalle vital : 7,35 -- 7,45). Notre organisme dispose de systèmes de régulation ultra-rapides (respiration : élimination \ce{CO2} par les poumons en quelques minutes) et lents (reins : élimination/reabsorption \ce{H+} et \ce{HCO3-} en heures/jours). Une chute du pH sanguin de seulement \textbf{0,4 unité} (pH $\approx$ 7,0, acidose) ou une hausse de 0,4 unité (pH $\approx$ 7,8, alcalose) peut provoquer le coma, des troubles cardiaques graves, voire la mort. C'est pourquoi les médecins surveillent le pH du sang (gaz du sang) des patients en réanimation. Le corps humain est une véritable usine de contrôle du pH !
\end{savaistu}

\courssub{Sécurité et applications concrètes}

\begin{aretenir}[Pictogrammes de danger (règlement CLP) — À reconnaître au BEPC]
\begin{center}
\begin{tabularx}{\linewidth}{|c|X|c|}\hline
\rowcolor{popblueL} \textbf{Pictogramme} & \textbf{Signification} & \textbf{Exemples courants} \\ \hline
\begin{tikzpicture}[scale=0.3]\draw[poppink,thick] (0,0) circle (1); \draw[poppink,thick] (-0.7,-0.7) -- (0.7,0.7); \draw[poppink,thick] (0.7,-0.7) -- (-0.7,0.7);\end{tikzpicture} & \textbf{Corrosif} (attaque métaux, peau, yeux) & Acide sulfurique (batterie), Soude caustique, Eau de Javel concentrée \\ \hline
\begin{tikzpicture}[scale=0.3]\draw[poporange,thick] (0,0) circle (1); \fill[poporange] (0,0.3) circle (0.15); \draw[poporange,thick] (-0.4,-0.3) -- (0.4,-0.3);\end{tikzpicture} & \textbf{Irritant / Nocif} (irritation peau/yeux, toxicité légère) & Eau de Javel diluée, Ammoniaque, Lessives, Détartrant \\ \hline
\begin{tikzpicture}[scale=0.3]\draw[poporange,thick] (0,0) circle (1); \draw[poporange,thick] (-0.5,0.5) -- (0.5,-0.5); \draw[poporange,thick] (0.5,0.5) -- (-0.5,-0.5);\end{tikzpicture} & \textbf{Comburant} (active la combustion) & Eau oxygénée concentrée, Nitrates, Peroxydes \\ \hline
\begin{tikzpicture}[scale=0.3]\draw[popgreen,thick] (0,0) circle (1); \draw[popgreen,thick] (0,-0.8) -- (0,0.5); \draw[popgreen,thick] (-0.4,-0.2) -- (0.4,-0.2);\end{tikzpicture} & \textbf{Dangereux pour l'environnement} & Pesticides, Hydrocarbures (essence, gasoil SONARA), Métaux lourds \\ \hline
\end{tabularx}
\end{center}
\end{aretenir}

\begin{methode}[Règles d'or de sécurité avec les acides et bases (laboratoire / maison)]
\begin{enumerate}
\item \textbf{Protection :} Lunettes (obligatoires), gants (nitrile), blouse, chaussures fermées. Pas de lentilles de contact (risque de piégeage du produit).
\item \textbf{Manipulation :} Toujours verser \textbf{l'acide (ou la base concentrée) DANS l'eau} (ex: dilution acide batterie, préparation lessive). Jamais l'inverse (éclaboussures brûlantes par ébullition locale).
\item \textbf{Étiquetage :} Tout récipient (bécher, fiole, bouteille) doit porter une étiquette : nom, concentration, date, pictogramme danger, nom du manipulateur.
\item \textbf{Déversement :} Petit volume : neutraliser (argile absorbante, bicarbonate pour acide, acide citrique pour base) puis ramasser. Grand volume : évacuer, aérer, appeler les secours (pompiers 112).
\item \textbf{Projection sur peau/yeux :} Rincer \textbf{immédiatement et abondamment} à l'eau courante pendant \textbf{au moins 15 minutes} (douche de sécurité, lava-yeux). Retirer vêtements contaminés. Consulter un médecin.
\item \textbf{Stockage :} Acides et bases séparés. Produits corrosifs en bas d'armoire (bac de rétention). Incompatibles séparés (ex: acide + hypochlorite = gaz chlore toxique !).
\end{enumerate}
\end{methode}

\begin{exemplebox}[Neutralisation domestique : déboucher un évier]
Un évier bouché par des graisses et cheveux. On verse \SI{100}{g} de cristaux de soude (\ce{Na2CO3}, basique) puis \SI{200}{mL} de vinaigre blanc (acide acétique \ce{CH3COOH}, pH $\approx$ 3).
Réaction : $\ce{CO3^{2-} + 2 CH3COOH -> 2 CH3COO- + H2O + CO2}\,\uparrow$.
L'effervescence (\ce{CO2}) mécanique + l'action basique (saponification des graisses) débouchent le tuyau. \textbf{Attention :} Ne jamais mélanger eau de Javel (basique, chlore) et acide (vinaigre, détartrant) $\rightarrow$ dégage du \textbf{dichlore \ce{Cl2} (gaz de combat, mortel)}.
\end{exemplebox}

\begin{exercice}[Bilan de la leçon — Préparation au BEPC]
\begin{enumerate}
\item \textbf{Vocabulaire :} Compléter : Dans une solution aqueuse de sucre, l'eau est le \_\_\_\_\_\_\_\_ et le sucre est le \_\_\_\_\_\_\_\_. Le mélange obtenu est \_\_\_\_\_\_\_\_ et \_\_\_\_\_\_\_\_.
\item \textbf{Solubilité :} À \SI{20}{\celsius}, la solubilité du chlorure de sodium est de \SI{36}{g/100g} d'eau. On introduit \SI{50}{g} de sel dans \SI{100}{g} d'eau à \SI{20}{\celsius}. Quelle masse de sel reste-t-elle non dissoute ? Quel est le caractère de la solution obtenue ?
\item \textbf{Concentration molaire :} On dissout \SI{5,85}{g} de NaCl ($M = \SI{58,5}{g/mol}$) dans l'eau et on complète à \SI{250}{mL} en fiole jaugée. Calculer la concentration molaire $C$ de cette solution.
\item \textbf{pH :} On mesure le pH de quatre solutions : $S_1$ (pH = 1), $S_2$ (pH = 7), $S_3$ (pH = 9), $S_4$ (pH = 13).
\begin{itemize}
\item[a)] Classer par caractère (acide, neutre, basique).
\item[b)] Laquelle est la plus acide ? La plus basique ?
\item[c)] $S_1$ est une solution d'acide chlorhydrique \SI{0,1}{mol/L}. $S_4$ est une solution de soude \SI{0,1}{mol/L}. Que se passe-t-il si on les mélange en volumes égaux ? (pH final ?)
\end{itemize}
\item \textbf{Sécurité :} Citer 3 équipements de protection individuelle (EPI) obligatoires pour manipuler de l'acide sulfurique concentré (batterie de moto). Quelle est la règle d'or pour diluer un acide concentré ?
\end{enumerate}
\end{exercice}

\begin{corrige}[Exercice : Bilan de la leçon]
\begin{enumerate}
\item Solvant / soluté / homogène / limpide (ou transparent).
\item Masse dissoute max = \SI{36}{g}. Masse non dissoute = \SI{50}{g} - \SI{36}{g} = \SI{14}{g}. La solution est \textbf{saturée}.
\item $n = \frac{5,85}{58,5} = \SI{0,10}{mol}$. $V = \SI{0,250}{L}$. $C = \frac{0,10}{0,250} = \SI{0,40}{mol/L}$.
\item 
\begin{itemize}
\item[a)] $S_1$ : acide ; $S_2$ : neutre ; $S_3$ : basique ; $S_4$ : basique.
\item[b)] Plus acide : $S_1$ (pH le plus petit). Plus basique : $S_4$ (pH le plus grand).
\item[c)] Mélange équimolaire acide fort + base forte $\rightarrow$ neutralisation totale. Solution finale neutre : \textbf{pH = 7} (sel \ce{NaCl} + eau).
\end{itemize}
\item EPI : Lunettes de protection, gants (nitrile/caoutchouc), blouse (coton). Règle d'or : \textbf{Verser l'acide dans l'eau} (jamais l'eau dans l'acide).
\end{enumerate}
\end{corrige}

\begin{aretenir}[L'essentiel de la Leçon 6 : Solutions aqueuses]
\begin{itemize}
\item \textbf{Dissolution :} Soluté (minoritaire) + Solvant (majoritaire, eau pour solution aqueuse) $\rightarrow$ Solution (homogène, limpide, stable). Conservation de la masse.
\item \textbf{Solubilité :} Masse max de soluté dissous par \SI{100}{g} (ou \SI{1}{L}) de solvant à $T$ fixée. Solution saturée = limite atteinte. Solubilité des solides $\uparrow$ avec $T$.
\item \textbf{Concentration molaire :} $C = n/V$ (mol/L). $n = m/M$. Dilution : $C_1 V_1 = C_2 V_2$ ($n$ constant). Préparation en fiole jaugée.
\item \textbf{pH :} Grandeur sans unité (0--14). pH < 7 acide ; pH = 7 neutre ; pH > 7 basique. Petit pH = très acide. Mesure : papier pH (rapide) ou pH-mètre (précis). Sécurité : gants, lunettes, verser acide dans l'eau.
\item \textbf{Applications :} Eau potable (CDE), agriculture (chaulage cacao), santé (sang pH 7,4, antiacides), hygiène (pH peau 5,5), industrie (SONARA, effluents).
\end{itemize}
\end{aretenir}

\coursec{Sécurité et applications concrètes}

Nous savons maintenant identifier une solution acide, neutre ou basique grâce au pH, et calculer sa concentration molaire. Mais manipuler des solutions aqueuses — surtout des acides et des bases — exige une rigueur absolue. Au laboratoire comme à la maison, un geste imprudent peut provoquer de graves brûlures ou dégager des gaz mortels. Cette section te donne les réflexes vitaux et te montre comment la chimie des solutions s'applique concrètement à ton quotidien : eau de forage, cultures de cacao ou de maïs, produits ménagers.

\courssub{Consignes de sécurité et pictogrammes de danger}

\begin{aretenir}[Règles d'or au laboratoire et à la maison]
\begin{itemize}
    \item \textbf{Ne jamais goûter} un produit chimique, même s'il ressemble à de l'eau ou à une boisson (l'acide sulfurique dilué est incolore et inodore).
    \item \textbf{Porter des gants} (nitrile ou latex épais) et des lunettes de protection pour manipuler les acides et bases \emph{concentrés} (acide chlorhydrique, soude caustique, eau de Javel concentrée).
    \item \textbf{Ne jamais verser d'eau dans un acide concentré} : risque de projection violente par échauffement brutal. On verse \emph{l'acide dans l'eau} en remuant.
    \item \textbf{Interdiction formelle de mélanger l'eau de Javel (solution d'hypochlorite de sodium $\ce{NaClO}$) avec un produit acide} (détartrant, vinaigre blanc, acide chlorhydrique) : cela libère du dichlore $\ce{Cl2}$, un gaz très toxique, suffocant, utilisé comme arme chimique lors de la Première Guerre mondiale.
\end{itemize}
\end{aretenir}

\begin{attention}[Danger mortel — Eau de Javel + Acide = Gaz dichlore]
Le mélange \ce{NaClO} (eau de Javel) + \ce{2H+} (acide) $\rightarrow$ \ce{Cl2} (gaz) + \ce{Na+} + \ce{H2O}.
Le dichlore $\ce{Cl2}$ irrite violemment les voies respiratoires, provoque des œdèmes pulmonaires et peut tuer en quelques minutes dans un local mal ventilé. \textbf{Ne mélange jamais deux produits ménagers différents.} Utilise-les l'un après l'autre en rinçant abondamment.
\end{attention}

Les produits chimiques (bidons d'acide, flacons de soude, bouteilles d'eau de Javel, lessives) portent des \textbf{pictogrammes de danger} normalisés internationalement (système SGH). Savoir les lire protège ta vie.

\begin{popfigure}
\begin{tikzpicture}[scale=0.9, every node/.style={font=\small}]
% Couleurs
\colorlet{picred}{red!80!black}
\colorlet{picblack}{black}
\colorlet{picwhite}{white}
\colorlet{picblue}{blue!60!black}
\colorlet{bgcorr}{red!10!white}
\colorlet{bgirrit}{orange!10!white}
\colorlet{bgfirstaid}{green!10!white}

% --- COLONNE 1 : CORROSIF ---
\begin{scope}[xshift=0cm]
\fill[bgcorr] (0,0) rectangle (5.5,7.5);
\draw[thick, picred] (0,0) rectangle (5.5,7.5);
% Pictogramme Corrosif (GHS05)
\begin{scope}[shift={(2.75,5.5)}, scale=0.6]
\draw[thick, picblack] (0,0) circle (1.5);
% Tube qui verse
\draw[thick, picblack, fill=picblack] (-0.3,1.2) rectangle (0.3,2.0);
\draw[thick, picblack] (-0.5,1.2) -- (-0.5,0.5) -- (0.5,0.5) -- (0.5,1.2);
% Gouttes
\foreach \y in {0.2, -0.2, -0.6} {
  \draw[thick, picblack, fill=picblue!50] (0,\y) ellipse (0.25 and 0.15);
}
% Main
\draw[thick, picblack] (-1.0,-0.8) .. controls (-1.5,-1.5) and (-0.5,-2.0) .. (0.0,-1.8);
\draw[thick, picblack] (1.0,-0.8) .. controls (1.5,-1.5) and (0.5,-2.0) .. (0.0,-1.8);
% Surface corrodée
\draw[thick, picblack, fill=picred!30] (-1.2,-1.8) -- (1.2,-1.8) -- (1.0,-2.2) -- (-1.0,-2.2) -- cycle;
\draw[thick, picblack] (-0.8,-1.8) -- (-0.6,-2.5) (-0.4,-1.8) -- (-0.2,-2.5) (0,-1.8) -- (0,-2.5) (0.4,-1.8) -- (0.2,-2.5) (0.8,-1.8) -- (0.6,-2.5);
\end{scope}
\node[align=center, text width=5cm] at (2.75, 3.8) {\textbf{Pictogramme : Corrosif (GHS05)}\\(Crâne + os croisés \emph{ou} tube sur main/métal)};
\node[align=center, text width=5cm] at (2.75, 2.8) {\textbf{Mention de danger $H_{314}$} : « Provoque de graves brûlures de la peau et des lésions oculaires graves. »};
\node[align=center, text width=5cm] at (2.75, 1.8) {\textbf{Exemples} : Acide chlorhydrique concentré, Soude caustique (NaOH), Acide sulfurique (batteries).};
\node[align=center, text width=5cm, fill=bgfirstaid, draw=green!60!black, rounded corners, inner sep=4pt] at (2.75, 0.6) {\textbf{Conduite à tenir} : \textbf{Peau/Yeux} $\rightarrow$ Rincer \textbf{15 min} min à l'eau courante. \textbf{Inhalation} $\rightarrow$ Air frais. \textbf{Ingestion} $\rightarrow$ Ne pas faire vomir, rincer la bouche, appeler le 112 / Centre antipoison.};
\end{scope}

% --- COLONNE 2 : IRRITANT / NOCIF ---
\begin{scope}[xshift=6.5cm]
\fill[bgirrit] (0,0) rectangle (5.5,7.5);
\draw[thick, orange!80!black] (0,0) rectangle (5.5,7.5);
% Pictogramme Irritant (GHS07) - Point d'exclamation
\begin{scope}[shift={(2.75,5.5)}, scale=0.7]
\draw[thick, picblack] (0,0) circle (1.5);
\node[font=\bfseries\Huge, text=picblack] at (0,0) {!};
\end{scope}
\node[align=center, text width=5cm] at (2.75, 3.8) {\textbf{Pictogramme : Irritant / Nocif (GHS07)}\\(Point d'exclamation)};
\node[align=center, text width=5cm] at (2.75, 2.8) {\textbf{Mentions de danger $H_{315}$/$H_{319}$} : « Provoque une irritation cutanée / une sévère irritation des yeux. »};
\node[align=center, text width=5cm] at (2.75, 1.8) {\textbf{Exemples} : Eau de Javel (diluée), Lessives, Détartrants doux, Ammoniaque dilué.};
\node[align=center, text width=5cm, fill=bgfirstaid, draw=green!60!black, rounded corners, inner sep=4pt] at (2.75, 0.6) {\textbf{Conduite à tenir} : \textbf{Peau} $\rightarrow$ Laver à l'eau et au savon. \textbf{Yeux} $\rightarrow$ Rincer \textbf{10-15 min} (paupières écartées). Si gêne persiste $\rightarrow$ Médecin.};
\end{scope}

% --- COLONNE 3 : CONSEILS DE PRUDENCE (P) ---
\begin{scope}[xshift=13cm]
\fill[popblueL] (0,0) rectangle (5.5,7.5);
\draw[thick, popblue] (0,0) rectangle (5.5,7.5);
\node[align=center, text width=5cm, font=\bfseries] at (2.75, 7.0) {Conseils de prudence (Codes P)};
\node[align=left, text width=5cm] at (2.75, 6.3) {
\textbf{$P_{101}$} : En cas de consultation d'un médecin, garder l'étiquette.\\[2pt]
\textbf{$P_{102}$} : Tenir hors de portée des enfants.\\[2pt]
\textbf{$P_{260}$} : Ne pas respirer les vapeurs.\\[2pt]
\textbf{$P_{280}$} : Porter des gants / protection oculaire.\\[2pt]
\textbf{$P_{301}+P_{330}+P_{331}$} : SI INGÉRÉ : rincer la bouche. NE PAS faire vomir.\\[2pt]
\textbf{$P_{305}+P_{351}+P_{338}$} : SI DANS LES YEUX : rincer avec précaution pendant plusieurs minutes. Retirer les lentilles si possible. Continuer à rincer.\\[2pt]
\textbf{$P_{309}+P_{310}$} : SI exposition ou malaise : appeler un CENTRE ANTIPOISON / médecin.\\[2pt]
\textbf{$P_{405}$} : Garder sous clef.\\[2pt]
\textbf{$P_{501}$} : Éliminer le contenu/récipient selon la réglementation locale.
};
\end{scope}
\end{tikzpicture}
\legende{Tableau récapitulatif des deux pictogrammes les plus fréquents sur les produits ménagers et réactifs de laboratoire (Corrosif GHS05 et Irritant GHS07) avec mentions de danger (H) et conduites à tenir. À droite, les conseils de prudence (P) standards à mémoriser.}
\end{popfigure}

\begin{methode}[Lire une étiquette de produit chimique (Méthode 3 étapes)]
\begin{enumerate}
    \item \textbf{Repérer le(s) pictogramme(s)} (losange rouge sur fond blanc) : Corrosif ? Inflammable ? Toxique ? Irritant ? Danger pour l'environnement ?
    \item \textbf{Lire la mention de danger (phrase H)} : Elle décrit \emph{la nature} du risque (ex. « $H_{314}$ : Provoque de graves brûlures »).
    \item \textbf{Appliquer les conseils de prudence (phrases P)} : Ils disent \emph{quoi faire} (protection, réaction en cas d'accident, stockage, élimination). Ex. « $P_{280}$ : Porter des gants », « $P_{305}+P_{351}+P_{338}$ : Si dans les yeux, rincer... ».
\end{enumerate}
\end{methode}

\courssub{Conduite à tenir en cas de contact accidentel}

Malgré la prévention, un accident peut survenir. La règle unique : \textbf{l'eau, tout de suite, longtemps}.

\begin{propriete}[Protocole d'urgence — Contact peau ou yeux]
\begin{enumerate}
    \item \textbf{Alerter} immédiatement l'enseignant (ou un adulte / les secours 112 / Centre Antipoison du Cameroun : \textbf{22 20 11 22}).
    \item \textbf{Rincer à grande eau courante} (robinet, douche de sécurité, fontaine oculaire) \textbf{pendant 15 minutes minimum} (chronométrer).
    \begin{itemize}
        \item \emph{Peau} : Enlever les vêtements contaminés pendant le rinçage (ne pas tirer par la tête si le visage est touché, couper le vêtement).
        \item \emph{Yeux} : Écarter les paupières avec les doigts, rouler le regard dans tous les sens. Si tu portes des lentilles : les retirer \emph{après} avoir commencé le rinçage.
    \end{itemize}
    \item \textbf{Ne rien appliquer d'autre} (pas de pommade, pas de neutralisant chimique : la réaction de neutralisation chauffe et aggrave la brûlure).
    \item \textbf{Consulter un médecin} systématiquement après un contact oculaire ou une brûlure cutanée étendue.
\end{enumerate}
\end{propriete}

\courssub{Applications concrètes au Cameroun}

La chimie des solutions aqueuses n'est pas abstraite : elle gère ta santé, ton alimentation et ton porte-monnaie.

\courssub{Traitement de l'eau de forage et de source}

Au Cameroun, de nombreux foyers (Yaoundé, Douala, zones rurales) boivent l'eau de forage ou de source. Cette eau traverse des couches géologiques qui la chargent en ions ($\ce{Ca^{2+}}$, $\ce{Mg^{2+}}$, $\ce{Fe^{2+}}$, $\ce{HCO3-}$) et peut être acide (pH $<$ 6,5) ou basique (pH $>$ 8,5). Une eau trop acide corrode les canalisations (plomb, cuivre) et libère des métaux toxiques ; une eau trop basique a un goût savonneux et précipite le calcaire.

\begin{exemplebox}[Qualité de l'eau de boisson — Normes OMS / MINSANTÉ]
L'Organisation Mondiale de la Santé (OMS) et le Ministère de la Santé Publique recommandent pour l'eau destinée à la consommation humaine :
\[
\boxed{6,5 \le \text{pH} \le 8,5}
\]
\begin{itemize}
    \item \textbf{pH $<$ 6,5} : Eau agressive, corrosion tuyauteries, risque métallique. \emph{Traitement} : filtration sur calcite (carbonate de calcium $\ce{CaCO3}$) pour remonter le pH.
    \item \textbf{pH $>$ 8,5} : Eau incrustante, goût désagréable, efficacité du chlore réduite. \emph{Traitement} : injection d'acide (souvent $\ce{CO2}$ ou $\ce{H2SO4}$ dilué) ou adoucisseur.
\end{itemize}
\textbf{Contrôle} : On utilise des bandelettes pH ou un pH-mètre portable (étalonné avec solutions tampons pH 4, 7, 10) avant toute consommation ou après traitement.
\end{exemplebox}

\begin{exoresolu}[Correction du pH d'une eau de forage]
\textbf{Énoncé.} Une eau de forage a un pH = 5,8 (acide). On décide de la faire passer dans un filtre à calcite ($\ce{CaCO3}$) pour la neutraliser partiellement. La réaction est :
\ce{CaCO3(s) + H2O(l) + CO2(aq) <=> Ca^{2+}(aq) + 2HCO3-(aq)}
On mesure une concentration en ions hydrogénocarbonate $[\ce{HCO3-}] = \SI{4,0e-3}{mol.L^{-1}}$ après traitement. Le volume d'eau traité est \SI{1000}{L}. Calculer la quantité de matière de $\ce{CaCO3}$ dissoute (en mol) et la masse correspondante (en g). $M(\ce{CaCO3}) = \SI{100}{g.mol^{-1}}$.
\tcblower
\textbf{Solution détaillée.}
\begin{enumerate}
    \item \textbf{Analyse de l'équation} : 1 mol $\ce{CaCO3}$ produit 2 mol $\ce{HCO3-}$.
    \item \textbf{Quantité de matière de $\ce{HCO3-}$ formée} :
    $n(\ce{HCO3-}) = C \times V = \SI{4,0e-3}{mol.L^{-1}} \times \SI{1000}{L} = \SI{4,0}{mol}$.
    \item \textbf{Quantité de matière de $\ce{CaCO3}$ dissoute} (stoïchiométrie) :
    $n(\ce{CaCO3}) = \frac{1}{2} \times n(\ce{HCO3-}) = \frac{1}{2} \times 4,0 = \SI{2,0}{mol}$.
    \item \textbf{Masse de $\ce{CaCO3}$} :
    $m(\ce{CaCO3}) = n \times M = \SI{2,0}{mol} \times \SI{100}{g.mol^{-1}} = \SI{200}{g}$.
\end{enumerate}
\textbf{Conclusion :} Il a fallu dissoudre \textbf{200 g de calcite} pour traiter 1000 L d'eau. Ce filtre doit être rechargé régulièrement.
\end{exoresolu}

\courssub{Choix des engrais selon le pH du sol : cacao et maïs}

L'agriculture est le poumon de l'économie camerounaise. Le cacaoyer et le maïs ont des exigences précises en pH du sol pour absorber les éléments nutritifs (azote N, phosphore P, potassium K).

\begin{savaistu}[Agronomie et pH du sol au Cameroun]
Dans les zones de culture du cacao (Sud, Centre, Est, Sud-Ouest), les sols sont souvent \textbf{acides} (pH 4,5 -- 5,5) à cause du lessivage intense par les pluies (ions $\ce{H+}$ remplacent $\ce{Ca^{2+}}$, $\ce{Mg^{2+}}$ sur l'argile). Un sol trop acide bloque le phosphore (forme des phosphates de fer/aluminium insolubles) et libère de l'aluminium toxique ($\ce{Al^{3+}}$) qui brûle les racines.
\begin{itemize}
    \item \textbf{Cacaoyer} : pH optimal \textbf{5,5 -- 6,5}. On apporte de la \textbf{chaux agricole} (carbonate de calcium $\ce{CaCO3}$ ou oxyde de calcium $\ce{CaO}$) pour monter le pH. Réaction : $\ce{CaO + 2H+ -> Ca^{2+} + H2O}$.
    \item \textbf{Maïs} : pH optimal \textbf{5,5 -- 7,0}. Sensible à l'acidité (toxicité Al) et à l'excès de chaux (blocage du phosphore, du zinc, du bore).
\end{itemize}
Les agronomes de l'IRAD (Institut de Recherche Agricole pour le Développement) et de la SODECAO font des \textbf{analyses de sol} (méthode KCl 1M ou eau) avant de recommander la dose de chaux (tonnes/ha) et le type d'engrais (NPK 20-10-10, urée, etc.). Un engrais mal choisi sur sol acide = argent jeté + pollution.
\end{savaistu}

\courssub{Dilution des produits ménagers concentrés : notice et calcul}

L'eau de Javel concentrée (souvent \SI{12}{\%} chlore actif, soit $\approx \SI{1,6}{mol.L^{-1}}$ en $\ce{ClO-}$) ou l'acide chlorhydrique commercial (\SI{23}{\%}, $\approx \SI{7,6}{mol.L^{-1}}$) sont dangereux purs. Les notices imposent des dilutions précises.

\begin{propriete}[Loi de dilution (conservation de la quantité de matière du soluté)]
Lors d'une dilution, le soluté ne change pas, seul le volume de solvant augmente.
\[
n_{\text{initial}} = n_{\text{final}} \quad \Rightarrow \quad C_i \times V_i = C_f \times V_f
\]
$C_i, V_i$ : concentration et volume de la solution mère (concentrée).
$C_f, V_f$ : concentration et volume de la solution fille (diluée).
\end{propriete}

\begin{exoresolu}[Préparation d'une solution d'eau de Javel pour désinfection]
\textbf{Énoncé.} La notice d'une eau de Javel concentrée ($C_{\text{mère}} = \SI{1,6}{mol.L^{-1}}$ en $\ce{ClO-}$) indique : « Pour la désinfection des sols : 1 verre (200 mL) dans 10 L d'eau ». Vérifier la concentration finale en $\ce{ClO-}$ et exprimer le taux de chlore actif en g/L (sachant que 1 mol $\ce{ClO-}$ $\equiv$ 1 équivalent chlore, $M_{\text{Cl}_2} = \SI{71}{g.mol^{-1}}$).
\tcblower
\textbf{Solution détaillée.}
\begin{enumerate}
    \item \textbf{Données} : $C_i = \SI{1,6}{mol.L^{-1}}$ ; $V_i = \SI{200}{mL} = \SI{0,200}{L}$ ; $V_f = \SI{10}{L} + \SI{0,200}{L} \approx \SI{10,2}{L}$ (on néglige souvent $V_i$ devant $V_f$, mais soyons rigoureux).
    \item \textbf{Concentration finale molaire} :
    $C_f = \frac{C_i \times V_i}{V_f} = \frac{1,6 \times 0,200}{10,2} = \frac{0,32}{10,2} \approx \SI{0,0314}{mol.L^{-1}}$.
    \item \textbf{Taux de chlore actif (g/L)} : La concentration en chlore actif s'exprime souvent en g/L de $\ce{Cl2}$ équivalent.
    $C_{\text{massique}} = C_f \times M(\ce{Cl2}) = 0,0314 \times 71 \approx \SI{2,23}{g.L^{-1}}$.
\end{enumerate}
\textbf{Résultat :} La solution de désinfection est à $\approx \SI{0,031}{mol.L^{-1}}$ soit \textbf{2,2 g/L de chlore actif}. C'est conforme aux recommandations (souvent 0,1\% à 0,5\% chlore actif, soit 1 à 5 g/L).
\end{exoresolu}

\courssub{Exemple : lecture d'une fiche de données de sécurité (FDS) simplifiée}

Tout produit chimique professionnel (et tout produit ménager dangereux) doit être accompagné d'une FDS (16 sections). En voici une version simplifiée pour un bidon de \textbf{détartrant acide concentré (acide chlorhydrique \SI{10}{\%})}.

\begin{popfigure}
\begin{tikzpicture}[scale=0.85, every node/.style={font=\small}]
% Styles
\tikzset{
  box/.style={draw=popdark, thick, rounded corners=4pt, fill=white, minimum height=1.2cm, text width=5.5cm, align=left, inner sep=6pt},
  titlebox/.style={draw=popdark, thick, rounded corners=4pt, fill=popblue, text=white, font=\bfseries, minimum height=1cm, text width=17.5cm, align=center, inner sep=4pt},
  sectionbox/.style={draw=poporange, thick, rounded corners=4pt, fill=poporangeL, font=\bfseries, minimum height=0.8cm, text width=5.5cm, align=center, inner sep=4pt},
  arrow/.style={-Stealth, thick, popdark}
}

% Titre global
\node[titlebox] (title) at (8.75, 11.5) {FICHE DE DONNÉES DE SÉCURITÉ SIMPLIFIÉE — DÉTARTRANT ACIDE CONCENTRÉ (HCl 10\%)};

% Ligne 1 : Identification + Danger
\node[sectionbox] (id) at (3, 10.2) {1. IDENTIFICATION};
\node[box, below=0.2cm of id.south, anchor=north, align=center] (idtxt){
\textbf{Produit} : Détartrant sanitaire fort\\
\textbf{Composition} : Acide chlorhydrique \SI{10}{\%} masse ($\approx \SI{3}{mol.L^{-1}}$)\\
\textbf{Usage} : Élimination calcaire WC, carrelage\\
\textbf{Fournisseur} : Entreprise locale (ex. SONARA Chimie)\\
\textbf{Tél. Urgence} : 112 / Centre Antipoison 22 20 11 22
};

\node[sectionbox] (haz) at (8.75, 10.2) {2. IDENTIFICATION DES DANGERS};
\node[box, below=0.2cm of haz.south, anchor=north, align=center] (haztxt){
\textbf{Classification} : Corrosif (Cat. 1B)\\
\textbf{Pictogramme} : \tikz[baseline=-2pt]{\draw[red,thick] (0,0) circle (0.2); \node[red,font=\bfseries] at (0,0) {\texttimes};} GHS05\\
\textbf{Mentions H} : $H_{314}$ (Brûlures graves), $H_{290}$ (Peut être corrosif pour les métaux)\\
\textbf{Mots de signal} : \textbf{DANGER}
};

\node[sectionbox] (comp) at (14.5, 10.2) {3. COMPOSITION};
\node[box, below=0.2cm of comp.south, anchor=north, align=center] (comptxt){
\textbf{Substance} : Chlorure d'hydrogène (HCl)\\
\textbf{N° CAS} : 7647-01-0\\
\textbf{Concentration} : 10 \% masse (3 M)\\
\textbf{Autres} : Eau, inhibiteur de corrosion, colorant, parfum
};

% Ligne 2 : Premiers secours + Lutte incendie + Mesures accidentelles
\node[sectionbox] (fa) at (3, 7.0) {4. PREMIERS SECOURS};
\node[box, below=0.2cm of fa.south, anchor=north, align=center] (fatxt){
\textbf{Inhalation} : Air frais, position demi-assise, O2 si besoin, médecin.\\
\textbf{Peau} : \textbf{Rincer 15 min eau courante}, enlever vêtements, médecin.\\
\textbf{Yeux} : \textbf{Rincer 15 min} (paupières ouvertes), médecin \textbf{urgent}.\\
\textbf{Ingestion} : \textbf{Ne pas faire vomir}, rincer bouche, boire eau/lait, médecin urgent.
};

\node[sectionbox] (fire) at (8.75, 7.0) {5. LUTTE CONTRE L'INCENDIE};
\node[box, below=0.2cm of fire.south, anchor=north, align=center] (firetxt){
\textbf{Moyens} : Eau en spray, poudre ABC, $\ce{CO2}$.\\
\textbf{Risques spécifiques} : Dégage $\ce{HCl}$ gazeux toxique/corrosif si chauffé.\\
\textbf{Protection pompiers} : ARI (Appareil Respiratoire Isolant), tenue étanche.
};

\node[sectionbox] (spill) at (14.5, 7.0) {6. MESURES ACCIDENTELLES};
\node[box, below=0.2cm of spill.south, anchor=north, align=center] (spilltxt){
\textbf{Précautions} : Éviter contact, ventiler, porter EPI (gants, lunettes, bottes, tablier).\\
\textbf{Nettoyage} : Confiner, absorber (sable, vermiculite), neutraliser \emph{lentement} avec $\ce{NaHCO3}$ (bicarbonate) ou chaux. Ne pas rejeter à l'égout pur. Recueillir en fût étiqueté.
};

% Ligne 3 : Manipulation/Stockage + Propriétés + Stabilité
\node[sectionbox] (hand) at (3, 3.8) {7. MANIPULATION / STOCKAGE};
\node[box, below=0.2cm of hand.south, anchor=north, align=center] (handtxt){
\textbf{Manipulation} : Hotte ou local ventilé. Verser acide \textbf{dans} l'eau. Pas de contact yeux/peau.\\
\textbf{Stockage} : Local frais, ventilé, verrouillé ($P_{405}$). Récipient fermé, en PEHD. \textbf{Séparer des bases, hypochlorites (Javel), métaux, aliments}. Bac de rétention obligatoire.
};

\node[sectionbox] (prop) at (8.75, 3.8) {9. PROPRIÉTÉS PHYSICO-CHIMIQUES};
\node[box, below=0.2cm of prop.south, anchor=north, align=center] (proptxt){
\textbf{Aspect} : Liquide incolore à jaune pâle\\
\textbf{Odeur} : Piquante, irritante (HCl gazeux)\\
\textbf{pH} : $< 1$ (solution 10\%)\\
\textbf{Densité} : $\approx \SI{1,05}{g.cm^{-3}}$\\
\textbf{Solubilité} : Miscible à l'eau (exothermique)
};

\node[sectionbox] (stab) at (14.5, 3.8) {10. STABILITÉ / RÉACTIVITÉ};
\node[box, below=0.2cm of stab.south, anchor=north, align=center] (stabtxt){
\textbf{Stable} : Oui, aux conditions normales.\\
\textbf{À éviter} : Chaleur, contact métaux (dégage $\ce{H2}$ inflammable), \textbf{bases fortes}, \textbf{oxidants (hypochlorites $\rightarrow$ $\ce{Cl2}$)}.\\
\textbf{Produits décomposition} : Chlorure d'hydrogène gazeux.
};

% Ligne 4 : Toxico + Éco + Transport + Réglementation
\node[sectionbox] (tox) at (3, 0.6) {11. TOXICOLOGIE};
\node[box, below=0.2cm of tox.south, anchor=north, align=center] (toxtxt){
\textbf{Voies} : Inhalation, cutanée, oculaire, ingestion.\\
\textbf{Effets aigus} : Brûlures immédiates, nécrose, œdème glotte/poumon.\\
\textbf{Chroniques} : Érosion dentaire, bronchite chronique si expositions répétées.
};

\node[sectionbox] (eco) at (8.75, 0.6) {12. ÉCOLOGIE};
\node[box, below=0.2cm of eco.south, anchor=north, align=center] (ecotxt){
\textbf{Écotoxicité} : Très toxique pour organismes aquatiques (pH bas).\\
\textbf{Biodégradabilité} : Inorganique, ne biodégrade pas. Neutralisation = sel ($\ce{NaCl}$, $\ce{CaCl2}$).\\
\textbf{Élimination} : Neutraliser puis rejeter selon arrêté préfectoral / réglementation camerounaise (MINEPDED).
};

\node[sectionbox] (trans) at (14.5, 0.6) {14. TRANSPORT / 15. RÉGLEMENTATION};
\node[box, below=0.2cm of trans.south, anchor=north, align=center] (transtxt){
\textbf{ONU} : 1789 (Acide chlorhydrique)\\
\textbf{Classe} : 8 (Corrosif)\\
\textbf{Groupe emballage} : II ou III\\
\textbf{Étiquette} : Corrosif (GHS05)\\
\textbf{Réglementation CM} : Arrêté MINEPDED sur produits chimiques, Code du travail (hygiène/sécurité).
};

% Flèches de lecture
\draw[arrow] (idtxt.south east) -- ++(0.5,0) |- (haztxt.north west);
\draw[arrow] (haztxt.south east) -- ++(0.5,0) |- (comptxt.north west);
\draw[arrow] (idtxt.south) -- ++(0,-0.3) -| (fatxt.north);
\draw[arrow] (fatxt.south east) -- ++(0.5,0) |- (firetxt.north west);
\draw[arrow] (firetxt.south east) -- ++(0.5,0) |- (spilltxt.north west);
\draw[arrow] (fatxt.south) -- ++(0,-0.3) -| (handtxt.north);
\draw[arrow] (handtxt.south east) -- ++(0.5,0) |- (proptxt.north west);
\draw[arrow] (proptxt.south east) -- ++(0.5,0) |- (stabtxt.north west);
\draw[arrow] (handtxt.south) -- ++(0,-0.3) -| (toxtxt.north);
\draw[arrow] (toxtxt.south east) -- ++(0.5,0) |- (ecotxt.north west);
\draw[arrow] (ecotxt.south east) -- ++(0.5,0) |- (transtxt.north west);

\end{tikzpicture}
\legende{Exemple de Fiche de Données de Sécurité (FDS) simplifiée pour un détartrant acide (HCl 10\%). Les 16 sections officielles sont condensées en 9 blocs clés. Repère les pictogrammes, les phrases H (danger) et P (prudence). Ce document doit être accessible à tout utilisateur.}
\end{popfigure}

\begin{exoresolu}[Exploitation d'une FDS — Détartrant acide]
\textbf{Énoncé.} À partir de la FDS ci-dessus, répondre aux questions :
\begin{enumerate}
    \item Quel pictogramme figure sur le bidon ? Quelle est la mention de danger H associée ?
    \item Tu renverses 50 mL de ce produit sur le plan de travail. Quelle conduite tenir ? (Cite 3 actions précises).
    \item Pourquoi est-il \emph{interdit} de stocker ce bidon à côté de l'eau de Javel ?
    \item La notice dit : « Diluer 100 mL dans 5 L d'eau pour nettoyage courant ». Calculer la concentration molaire finale de la solution diluée. ($C_{\text{mère}} = \SI{3,0}{mol.L^{-1}}$).
\end{enumerate}
\tcblower
\textbf{Solution détaillée.}
\begin{enumerate}
    \item \textbf{Pictogramme} : GHS05 (Corrosif — tube qui verse sur main et métal). \textbf{Mention H} : $H_{314}$ (« Provoque de graves brûlures de la peau et des lésions oculaires graves »).
    \item \textbf{Conduite tenue} (Section 6) :
    \begin{enumerate}
        \item Aérer la pièce / ventiler (éviter inhalation vapeurs HCl).
        \item Porter gants + lunettes (EPI) avant d'intervenir.
        \item Absorber avec matière inerte (sable, chiffons) puis neutraliser \emph{lentement} avec du bicarbonate de soude ($\ce{NaHCO3}$) — ça mousse ($\ce{CO2}$) — jusqu'à cessation de réaction. Recueillir les déchets dans un seau pour élimination filière adaptée. Rincer la surface à grande eau.
    \end{enumerate}
    \item \textbf{Incompatibilité majeure} (Section 10) : Acide ($\ce{H+}$) + Hypochlorite ($\ce{ClO-}$ de l'eau de Javel) $\rightarrow$ \textbf{Dichlore $\ce{Cl2}$ gazeux} (très toxique, mortel). Risque de réaction accidentelle si fuites ou vapeurs se mélangent.
    \item \textbf{Calcul dilution} : $C_f = \frac{C_i V_i}{V_f} = \frac{3,0 \times 0,100}{5,100} \approx \SI{0,0588}{mol.L^{-1}}$. (On prend $V_f = 5,1$ L pour être exact, ou $5,0$ L si approximation $V_i \ll V_f$ : $C_f \approx 0,060$ M).
\end{enumerate}
\end{exoresolu}

\begin{aretenir}[Bilan — Sécurité et Applications]
\begin{itemize}
    \item \textbf{Sécurité} : Zéro tolérance — pas de goût, gants/lunettes pour concentrés, \textbf{jamais Javel + Acide} (gaz $\ce{Cl2}$ mortel). Pictogrammes GHS05 (Corrosif) et GHS07 (Irritant) à reconnaître instantanément.
    \item \textbf{Accident} : \textbf{Rincer 15 min à l'eau courante} (peau/yeux), alerter, ne pas neutraliser chimiquement sur la peau.
    \item \textbf{Eau de forage} : Contrôle pH (cible 6,5--8,5). Correction par calcite si acide, $\ce{CO2}$/acide si basique.
    \item \textbf{Agriculture} : pH sol conditionne disponibilité NPK. Cacao/maïs : pH 5,5--6,5. Chaux (CaO/CaCO3) pour corriger acidité (lessivage zones forestières Cameroun).
    \item \textbf{Dilution} : $C_i V_i = C_f V_f$. Respecter scrupuleusement les notices produits ménagers (eau de Javel, détartrants).
    \item \textbf{FDS} : Document légal de référence (16 sections). Lire : Identification $\rightarrow$ Dangers (H/P) $\rightarrow$ Premiers secours $\rightarrow$ Manipulation/Stockage $\rightarrow$ Élimination.
\end{itemize}
\end{aretenir}

\newpage

\coursec{Activité d'intégration}

\coursec{Solutions aqueuses : Activité d'intégration — Le cas du forage de Mokolo}

\courssub{Objectifs de l'activité}
Cette activité d'intégration vise à mobiliser l'ensemble des notions abordées dans la leçon « Solutions aqueuses » pour résoudre un problème concret de santé publique. À l'issue de ce travail, vous devez être capable de :
\begin{itemize}
    \item Identifier les constituants d'une solution (soluté, solvant) et écrire l'équation de dissolution.
    \item Calculer la concentration molaire d'une solution à partir de la masse de soluté et du volume de solution.
    \item Exploiter une courbe de solubilité pour déterminer l'état de saturation d'une solution.
    \item Interpréter une valeur de pH pour qualifier le caractère acido-basique d'une solution et vérifier sa conformité aux normes de potabilité.
    \item Calculer une concentration massique en ions et la comparer à une valeur guide organoleptique.
    \item Communiquer vos résultats à l'oral en structurant un raisonnement scientifique rigoureux et en rédigeant un avis technique argumenté.
\end{itemize}

\courssub{Situation : L'eau du forage du quartier « Mokolo » est-elle potable ?}
Dans le quartier Mokolo à Yaoundé, l'approvisionnement en eau potable par le réseau de la CAMWATER (Cameroon Water Utilities Corporation) étant parfois irrégulier, de nombreux habitants s'approvisionnent auprès d'un forage privé géré par un comité de quartier. Récemment, plusieurs familles ont signalé un goût inhabituel de l'eau et des irritations cutanées après la douche. Le comité de gestion mandate un laboratoire d'analyses pour contrôler la qualité de l'eau.

Le laboratoire remet le rapport d'analyse suivant pour un prélèvement de \SI{1,0}{L} d'eau brute du forage :
\begin{itemize}
    \item Masse de chlorure de sodium (\ce{NaCl}) dissous : \textbf{\SI{2,9}{g}}.
    \item pH mesuré à \SI{25}{\celsius} : \textbf{5,8}.
    \item Température de l'eau au moment du prélèvement : \SI{25}{\celsius}.
\end{itemize}

Le laboratoire joint également deux documents de référence :
\begin{enumerate}
    \item \textbf{Courbe de solubilité du chlorure de sodium (\ce{NaCl}) dans l'eau} en fonction de la température (document 1). On y lit qu'à \SI{25}{\celsius}, la solubilité maximale du \ce{NaCl} est de \SI{357}{g/L} (soit environ \SI{6,1}{mol/L}).
    \item \textbf{Normes de l'Organisation Mondiale de la Santé (OMS) pour l'eau de boisson} (document 2) : le pH doit être compris entre \textbf{6,5 et 8,5}. La concentration en chlorures (\ce{Cl-}) ne doit pas excéder \SI{250}{mg/L} (soit \SI{0,25}{g/L}) pour des raisons organoleptiques (goût), bien qu'aucune valeur limite sanitaire stricte ne soit fixée pour le \ce{NaCl} lui-même.
\end{enumerate}

\medskip
\textbf{Votre mission :} En tant qu'experts juniors en chimie, vous devez analyser ce dossier et rédiger un avis technique à l'intention du comité de quartier. Cet avis devra préciser la concentration molaire en \ce{NaCl}, l'état de saturation de la solution, son caractère acido-basique, sa conformité aux normes OMS, et proposer une conduite à tenir.

\courssub{Consignes}
Travaillez par groupes de 3 ou 4 élèves. Répondez aux questions suivantes en justifiant chaque étape par un calcul, une lecture de document ou une définition du cours. Préparez une restitution orale de 5 minutes.

\begin{enumerate}
    \item \textbf{Identification :} Quels sont le soluté et le solvant dans cette situation ? Écrire la dissolution du \ce{NaCl} dans l'eau sous forme d'équation chimique.
    \item \textbf{Concentration molaire :} Calculer la quantité de matière $n$ de \ce{NaCl} dissous dans \SI{1,0}{L} d'eau (masse molaire $M_{\ce{NaCl}} = \SI{58,5}{g/mol}$). En déduire la concentration molaire $C$ de la solution en \ce{NaCl} (en \SI{}{mol/L}).
    \item \textbf{Solubilité et saturation :} En vous référant au document 1 (courbe de solubilité), dire si la solution prélevée est \textbf{saturée} ou \textbf{non saturée} à \SI{25}{\celsius}. Justifier votre réponse par une comparaison numérique.
    \item \textbf{Caractère acido-basique :} En vous référant au document 2 (normes OMS) et à vos connaissances sur l'échelle de pH, qualifier la solution (acide, neutre, basique). Cette eau respecte-t-elle la norme de pH pour la potabilité ?
    \item \textbf{Concentration en chlorures :} Calculer la concentration massique en ions chlorure \ce{Cl-} (en \SI{}{g/L} et en \SI{}{mg/L}). Comparer à la valeur guide de l'OMS (\SI{250}{mg/L}). Commenter le goût probable de l'eau.
    \item \textbf{Synthèse et avis :} Rédiger la conclusion de votre avis technique : « L'eau de ce forage est-elle potable au regard des paramètres analysés ? » Proposer deux mesures concrètes à destination du comité de quartier (traitement, signalement, interdiction temporaire, etc.).
\end{enumerate}

\courssub{Ressources à mobiliser}

\begin{definition}[Solution aqueuse, soluté, solvant]
Une \cle{solution aqueuse} est un mélange homogène obtenu en dissolvant une substance (le \cle{soluté}) dans de l'eau (le \cle{solvant}). Le solvant est le composant majoritaire. La dissolution est un phénomène physique : le soluté se disperse à l'échelle microscopique (molécules ou ions) dans le solvant.
\end{definition}

\begin{propriete}[Concentration molaire d'une solution]
La concentration molaire $C$ (en \SI{}{mol/L} ou \SI{}{mol\cdot L^{-1}}) d'une solution est le quotient de la quantité de matière $n$ de soluté (en \SI{}{mol}) par le volume $V$ de la solution (en \SI{}{L}) :
\[ C = \frac{n}{V} \]
La quantité de matière $n$ (en \SI{}{mol}) se calcule à partir de la masse $m$ (en \SI{}{g}) et de la masse molaire $M$ (en \SI{}{g/mol}) : $n = \dfrac{m}{M}$.
\textbf{Au BEPC :} La formule $C = n/V$ et $n = m/M$ doivent être connues et appliquées. L'homogénéité des unités (L, g, mol) est obligatoire.
\end{propriete}

\begin{definition}[Solubilité et solution saturée]
La \cle{solubilité} $s$ d'un soluté dans un solvant à une température donnée est la concentration molaire (ou massique) maximale que peut atteindre la solution à cette température.
\begin{itemize}
    \item Solution \cle{saturée} : $C = s$ (présence éventuelle de solide en excès au fond).
    \item Solution \cle{non saturée} : $C < s$ (tout le soluté est dissous, on peut encore en dissoudre).
\end{itemize}
La solubilité dépend de la nature du soluté, du solvant et de la température.
\end{definition}

\begin{propriete}[Échelle de pH et normes OMS]
À \SI{25}{\celsius}, le pH permet de qualifier le caractère acido-basique d'une solution aqueuse :
\begin{itemize}
    \item pH $< 7$ : solution \cle{acide}.
    \item pH $= 7$ : solution \cle{neutre}.
    \item pH $> 7$ : solution \cle{basique}.
\end{itemize}
Pour l'eau de boisson, l'OMS recommande un pH entre \textbf{6,5 et 8,5}. Une eau de pH 5,8 est donc \textbf{non conforme} (trop acide), ce qui peut indiquer une pollution ou une corrosion des canalisations (relâchement de métaux).
\end{propriete}

\begin{popfigure}
\begin{tikzpicture}[scale=0.9, every node/.style={font=\small}]
% Axe pH
\draw[thick, popdark] (0,0) -- (14.5,0);
\foreach \x in {0,1,...,14} {
    \draw (\x,0.1) -- (\x,-0.1) node[below=2pt] {\x};
}
% Zones colorées
\fill[poppinkL] (0,0.2) rectangle (6.5,0.7) node[pos=0.5, above=4pt, text=popdark] {Domaine acide};
\fill[popgreenL] (6.5,0.2) rectangle (8.5,0.7) node[pos=0.5, above=4pt, text=popdark, align=center]{Norme OMS \\ $[$6,5 ; 8,5$]$};
\fill[popblueL] (8.5,0.2) rectangle (14,0.7) node[pos=0.5, above=4pt, text=popdark] {Domaine basique};
% Neutre
\draw[dashed, thick, popdark] (7,0) -- (7,0.7) node[above=4pt] {Neutre (7)};
% Échantillon
\draw[popred, very thick, ->] (5.8, -0.4) -- (5.8, -0.1) node[below=6pt, popred, font=\bfseries] {Eau forage : 5,8};
\end{tikzpicture}
\legende{Échelle de pH à 25 °C : position de l'eau du forage (pH 5,8) par rapport à la norme OMS de potabilité (zone verte).}
\end{popfigure}

\begin{popfigure}
\begin{tikzpicture}[scale=1.1, every node/.style={font=\small}]
% Réseau ionique NaCl (solide)
\foreach \x in {0,1,2} \foreach \y in {0,1,2} {
    \ifodd\numexpr\x+\y\relax
        \fill[poporange] (\x*0.5, \y*0.5) circle (0.18); % Na+
    \else
        \fill[popblue] (\x*0.5, \y*0.5) circle (0.18); % Cl-
    \fi
}
\node[align=center] at (0.75, -0.6) {Solide \ce{NaCl} \\ (réseau ionique)};
% Flèche dissolution
\draw[->, thick, popdark] (2.0, 0.5) -- (2.8, 0.5) node[midway, above] {\ce{H2O}};
% Ions solvatés (solution)
\begin{scope}[xshift=3.5cm]
\foreach \i in {1,...,12} {
    \pgfmathsetmacro{\rx}{0.5 + 1.5*rand}
    \pgfmathsetmacro{\ry}{0.5 + 1.0*rand}
    \ifodd\i
        \fill[poporange] (\rx, \ry) circle (0.15);
    \else
        \fill[popblue] (\rx, \ry) circle (0.15);
    \fi
}
% Molécules d'eau schématiques
\foreach \i in {1,...,8} {
    \pgfmathsetmacro{\rx}{0.5 + 1.5*rand}
    \pgfmathsetmacro{\ry}{0.5 + 1.0*rand}
    \draw[popblue!50, thin] (\rx, \ry) -- ++(0.1,0.1) (\rx, \ry) -- ++(-0.1,0.1);
}
\node[align=center] at (1.25, -0.6) {Solution aqueuse \\ \ce{Na+_{(aq)}} (orange), \ce{Cl-_{(aq)}} (bleu)};
\end{scope}
\end{tikzpicture}
\legende{Modélisation microscopique de la dissolution du chlorure de sodium : rupture du réseau ionique et solvatation des ions par les molécules d'eau.}
\end{popfigure}

\begin{attention}[Pièges fréquents à éviter]
\begin{itemize}
    \item Ne pas confondre \textbf{masse molaire} ($M$, en \SI{}{g/mol}) et \textbf{masse} ($m$, en \SI{}{g}).
    \item Le volume $V$ dans la formule $C = n/V$ est le volume de la \textbf{solution finale}, pas celui du solvant seul (ici, on considère que le volume de la solution est \SI{1,0}{L} car la quantité de sel est faible).
    \item La solubilité du \ce{NaCl} varie peu avec la température : ne pas extrapoler sans document.
    \item Le pH est une grandeur \textbf{sans unité}.
    \item Pour passer de \SI{}{g/L} à \SI{}{mg/L}, on multiplie par 1000 ($1\,\text{g} = 1000\,\text{mg}$).
    \item L'équation de dissolution montre la dissociation ionique : \ce{NaCl_{(s)} ->[H2O] Na+_{(aq)} + Cl-_{(aq)}}. Ne pas écrire \ce{NaCl_{(aq)}}.
\end{itemize}
\end{attention}

\courssub{Méthode : Calculer une concentration molaire}
\begin{methode}[Étapes pour passer de la masse à la concentration molaire]
\begin{enumerate}
    \item \textbf{Lister les données :} $m$ (en g), $V$ (en L), $M$ (en g/mol). Convertir si nécessaire.
    \item \textbf{Calculer la quantité de matière :} $n = \dfrac{m}{M}$ (résultat en mol).
    \item \textbf{Calculer la concentration molaire :} $C = \dfrac{n}{V}$ (résultat en mol/L).
    \item \textbf{Conclure avec une phrase :} « La concentration molaire de la solution en [soluté] est de [valeur] mol/L. »
\end{enumerate}
\end{methode}

\courssub{Démarche et corrigé commenté}
\begin{exoresolu}[Analyse complète de l'eau du forage Mokolo]
\textbf{Énoncé :} Voir la situation et les consignes ci-dessus.

\tcblower

\textbf{Solution détaillée :}

\medskip
\textbf{1. Identification du système}
\begin{itemize}
    \item \textbf{Solvant :} L'eau (\ce{H2O}), composant majeur du prélèvement.
    \item \textbf{Soluté :} Le chlorure de sodium (\ce{NaCl}), sel dissous.
    \item \textbf{Équation de dissolution (dissociation ionique) :}
    \[ \ce{NaCl_{(s)} ->[H2O] Na+_{(aq)} + Cl-_{(aq)}} \]
    Le réseau ionique solide se brise sous l'action des molécules d'eau : les ions \ce{Na+} et \ce{Cl-} se séparent et sont solvatés (entourés par des molécules d'eau). La solution obtenue est incolore.
\end{itemize}

\medskip
\textbf{2. Calcul de la concentration molaire en \ce{NaCl}}
\begin{itemize}
    \item \textit{Données :} $m_{\ce{NaCl}} = \SI{2,9}{g}$ ; $V_{\text{solution}} = \SI{1,0}{L}$ ; $M_{\ce{NaCl}} = \SI{58,5}{g/mol}$.
    \item \textit{Quantité de matière :}
    \[ n_{\ce{NaCl}} = \frac{m_{\ce{NaCl}}}{M_{\ce{NaCl}}} = \frac{2,9}{58,5} \approx \SI{0,0496}{mol} \]
    \item \textit{Concentration molaire :}
    \[ C_{\ce{NaCl}} = \frac{n_{\ce{NaCl}}}{V_{\text{solution}}} = \frac{0,0496}{1,0} \approx \SI{0,050}{mol/L} \]
    \textbf{Résultat :} La concentration molaire de la solution en \ce{NaCl} est $C_{\ce{NaCl}} \approx \SI{5,0 \times 10^{-2}}{mol/L}$ (arrondi à 2 chiffres significatifs cohérents avec les données).
\end{itemize}

\medskip
\textbf{3. État de saturation (Exploitation du document 1)}
\begin{itemize}
    \item \textit{Solubilité massique du \ce{NaCl} à \SI{25}{\celsius} (lu sur la courbe) :} $s_m = \SI{357}{g/L}$.
    \item \textit{Conversion en concentration molaire pour comparer :}
    \[ s_{\text{molaire}} = \frac{s_m}{M_{\ce{NaCl}}} = \frac{357}{58,5} \approx \SI{6,1}{mol/L} \]
    \item \textit{Comparaison :} $C_{\text{mesurée}} = \SI{0,050}{mol/L}$ \textbf{est très inférieure à} $s = \SI{6,1}{mol/L}$.
    \item \textbf{Conclusion :} La solution est \textbf{non saturée} (très loin de la saturation). Il est possible de dissoudre encore beaucoup de sel dans cette eau à cette température. Aucun précipité n'est visible au fond du prélèvement.
\end{itemize}

\medskip
\textbf{4. Caractère acido-basique et conformité pH (Exploitation du document 2)}
\begin{itemize}
    \item \textit{pH mesuré :} $\text{pH} = 5,8$.
    \item \textit{Interprétation (échelle de pH) :} $5,8 < 7$, donc la solution est \textbf{acide}.
    \item \textit{Norme OMS (Document 2) :} Intervalle autorisé $[6,5 ; 8,5]$.
    \item \textit{Comparaison :} $5,8 < 6,5$. La valeur est \textbf{hors norme} (trop acide).
    \item \textbf{Conséquence :} Une eau acide est agressive pour les canalisations (risque de dissolution de métaux comme le plomb, le cuivre ou le fer) et peut irriter les muqueuses et la peau. Elle n'est \textbf{pas potable} au regard de ce critère sanitaire.
\end{itemize}

\medskip
\textbf{5. Concentration en chlorures \ce{Cl-} et aspect organoleptique (goût)}
\begin{itemize}
    \item \textit{Stœchiométrie de la dissolution :} 1 mol \ce{NaCl} $\rightarrow$ 1 mol \ce{Na+} + 1 mol \ce{Cl-}. Donc $n_{\ce{Cl-}} = n_{\ce{NaCl}} = \SI{0,0496}{mol}$.
    \item \textit{Masse molaire de l'ion chlorure :} $M_{\ce{Cl}} = \SI{35,5}{g/mol}$.
    \item \textit{Concentration massique en \ce{Cl-} :}
    \[ \rho_{\ce{Cl-}} = \frac{n_{\ce{Cl-}} \times M_{\ce{Cl}}}{V} = \frac{0,0496 \times 35,5}{1,0} \approx \SI{1,76}{g/L} \]
    \item \textit{Conversion en mg/L :} $\SI{1,76}{g/L} = \SI{1,76}{\times} \SI{1000}{mg/L} = \SI{1760}{mg/L}$.
    \item \textit{Comparaison à la valeur guide OMS :} Valeur guide = \SI{250}{mg/L}.
    \item \textbf{Conclusion :} $\SI{1760}{mg/L} \gg \SI{250}{mg/L}$. La concentration en chlorures dépasse \textbf{7 fois} la valeur guide organoleptique. L'eau a un \textbf{goût salé très prononcé} (eau saumure), ce qui explique les réclamations des usagers. Bien que non toxique à cette dose pour un adulte sain, ce goût rend l'eau impropre à la consommation courante (rejet par la population, risque de déshydratation si consommation exclusive).
\end{itemize}

\medskip
\textbf{6. Synthèse et Avis Technique au Comité de Quartier Mokolo}

\begin{center}
\begin{tabularx}{\linewidth}{|l|X|}\hline
\rowcolor{popblueL}
\textbf{Paramètre analysé} & \textbf{Résultat \& Conformité} \\ \hline
Concentration en \ce{NaCl} & \SI{0,050}{mol/L} (\SI{2,9}{g/L}) -- Solution non saturée \\ \hline
pH & 5,8 -- \textbf{NON CONFORME} (Norme OMS : 6,5 -- 8,5) \\ \hline
Chlorures (\ce{Cl-}) & \SI{1760}{mg/L} -- \textbf{NON CONFORME} (Guide OMS : \SI{250}{mg/L}) \\ \hline
\end{tabularx}
\end{center}

\textbf{Avis :} \textbf{L'eau de ce forage n'est PAS POTABLE.}
\begin{itemize}
    \item \textbf{Critère sanitaire (pH) :} Le pH acide (5,8) indique une agressivité chimique. Cela peut provenir d'infiltrations de matières organiques en décomposition (humus, tourbe), de rejets acides industriels ou agricoles, ou de la nature du sous-sol (roches granitiques pauvres en carbonates qui ne tamponnent pas l'acidité). Ce pH favorise la corrosion des tuyaux (risque de métaux lourds : plomb, cuivre, fer).
    \item \textbf{Critère organoleptique (Goût/Chlorures) :} La teneur en sel (\SI{2,9}{g/L}) est 7 fois supérieure au seuil de goût acceptable. La population boit une eau « saumure » légère.
\end{itemize}

\textbf{Conduite à tenir (Recommandations urgentes) :}
\begin{enumerate}
    \item \textbf{Interdiction immédiate de la consommation} pour boisson et cuisson (affichage aux bornes-fontaines, sensibilisation par les chefs de rue, annonces radio locale).
    \item \textbf{Recherche de la cause du pH acide :} Faire réaliser une analyse complète (métaux lourds : fer, manganèse, plomb ; nitrates ; bactériologie) par un laboratoire accrédité (ex: Centre Pasteur du Cameroun ou Laboratoire National de Santé Publique).
    \item \textbf{Traitement correctif :} Si la source est conservée, installer un \textbf{filtre neutralisant} (lit de calcite / marbre concassé) pour remonter le pH vers 7-7,5 par dissolution de \ce{CaCO3} ($\ce{CaCO3 + H+ -> Ca^2+ + HCO3-}$). Pour les chlorures, seule une \textbf{osmose inverse} ou une \textbf{distillation} les élimine (coûteux). Une alternative durable : raccorder le quartier au réseau CAMWATER ou forer un nouveau puits dans une autre nappe.
    \item \textbf{Surveillance :} Mettre en place un suivi mensuel du pH et de la conductivité (indicateur de salinité) par le comité, avec formation d'un « brigadier eau » local.
\end{enumerate}
\end{exoresolu}

\courssub{Bilan de l'activité}
Cette activité a permis de mettre en évidence la puissance de la chimie pour diagnostiquer un problème de santé publique réel. Vous avez mobilisé :
\begin{itemize}
    \item La \textbf{modélisation microscopique} (dissociation ionique) pour comprendre la composition réelle de l'eau.
    \item Le \textbf{calcul de concentration molaire} ($C = n/V$) pour quantifier la pollution saline.
    \item L'\textbf{exploitation de documents techniques} (courbe de solubilité, normes OMS) pour situer les mesures par rapport à des seuils de référence.
    \item L'\textbf{argumentation scientifique} : une conclusion ne vaut que si elle s'appuie sur des chiffres comparés à des normes.
\end{itemize}

\medskip
\textbf{Leçon retenue :} Une eau claire et inodore n'est pas forcément potable. Le pH et la minéralisation (chlorures, sulfates, fluorures) sont des paramètres invisibles mais décisifs. Au Cameroun, comme ailleurs, l'accès à l'eau sûre (ODD 6) passe par une surveillance chimique rigoureuse des points de distribution, qu'ils soient publics ou privés. Vous savez désormais lire un bulletin d'analyse et alerter votre communauté en cas de non-conformité.

\courssub{À retenir}
\begin{aretenir}[Notions clés de la leçon « Solutions aqueuses »]
\begin{itemize}
    \item \textbf{Solution} = Solvant (majoritaire) + Soluté (dissous). Aqueuse $\Rightarrow$ solvant = eau.
    \item \textbf{Dissolution du sel :} \ce{NaCl_{(s)} ->[H2O] Na+_{(aq)} + Cl-_{(aq)}} (dissociation ionique).
    \item \textbf{Quantité de matière :} $n = \dfrac{m}{M}$ ($n$ en mol, $m$ en g, $M$ en g/mol).
    \item \textbf{Concentration molaire :} $C = \dfrac{n}{V}$ ($C$ en mol/L, $V$ en L). \textbf{Formule à connaître pour le BEPC.}
    \item \textbf{Solubilité $s$ :} Concentration maximale à une T° donnée. Saturée si $C = s$, non saturée si $C < s$.
    \item \textbf{pH :} Sans unité. Acide si $<7$, Neutre si $=7$, Basique si $>7$. Norme OMS eau de boisson : $[6,5 ; 8,5]$.
    \item \textbf{Concentration massique :} $\rho = \dfrac{m}{V}$ (g/L) ou $\rho = C \times M$. Passage g/L $\leftrightarrow$ mg/L : $\times 1000$ ou $\div 1000$.
\end{itemize}
\end{aretenir}

\courssub{Le saviez-vous ?}
\begin{savaistu}[L'eau au Cameroun : enjeux et réalités]
\begin{itemize}
    \item \textbf{CAMWATER} gère la production et la distribution d'eau potable en milieu urbain. La qualité de l'eau distribuée est contrôlée selon les normes OMS (pH, chlore résiduel, turbidité, bactériologie).
    \item \textbf{Forages et puits traditionnels :} En zone périurbaine ou rurale, de nombreuses populations dépendent d'ouvrages privés (forages, puits) non contrôlés. L'acidité naturelle des nappes (roches cristallines du socle granitique) ou les pollutions anthropiques (latrines, déchets, agriculture) dégradent la qualité.
    \item \textbf{Chlorures et goût :} Le seuil de détection du goût salé est autour de \SI{200}{--}{250}{mg/L} en \ce{Cl-}. Au-delà, l'eau est rejetée. Une eau à \SI{1760}{mg/L} a un goût de « soupe » ou de « sauce » très marqué.
    \item \textbf{Traitement domestique :} Bouillir l'eau tue les microbes mais \textbf{ne change pas le pH ni la teneur en sel}. Seuls des filtres spécifiques (charbon actif, osmose inverse, résines échangeuses d'ions) ou l'adduction d'un réseau traité garantissent une eau conforme.
    \item \textbf{ODD 6 (Objectif de Développement Durable) :} « Garantir l'accès de tous à l'eau et à l'assainissement et assurer une gestion durable des ressources en eau ». La chimie (analyse, traitement) est au cœur de cet objectif.
\end{itemize}
\end{savaistu}

\newpage

\coursec{Méthodes et exercices résolus}

\coursec{Leçon 06 : Solutions aqueuses}

\courssub{Dissolution, soluté, solvant, solution}

\begin{definition}[Notions fondamentales]
Une \cle{dissolution} est l'action de mélanger une substance (le \cle{soluté}) dans une autre (le \cle{solvant}) pour obtenir un mélange homogène appelé \cle{solution}.
\begin{itemize}
\item Le \textbf{soluté} est l'espèce chimique dissoute (celle qui « disparaît » à l'œil nu). Il peut être solide (sel, sucre), liquide (alcool, vinaigre) ou gazeux (dioxyde de carbone dans les boissons).
\item Le \textbf{solvant} est l'espèce qui dissout (généralement un liquide, présent en plus grande quantité).
\item La \textbf{solution} est le mélange homogène obtenu. On ne distingue plus les constituants à l'œil nu.
\item Si le solvant est l'eau, on parle de \cle{solution aqueuse}.
\end{itemize}
\end{definition}

\begin{methode}[Reconnaître le soluté, le solvant et la solution]
Pour analyser une situation de dissolution (sucre dans le thé, sel dans la sauce, gaz dans une boisson \emph{Top} ou \emph{Malta}) :
\begin{enumerate}
\item Repérer l'espèce qui est \cle{dissoute} : c'est le \cle{soluté}.
\item Repérer l'espèce qui \cle{dissout} (liquide en plus grande quantité) : c'est le \cle{solvant}.
\item Nommer le mélange homogène obtenu : c'est la \cle{solution}.
\item Vérifier l'homogénéité : si un dépôt reste au fond, la dissolution est incomplète ou la solution est \cle{saturée}.
\item Préciser la nature : si le solvant est l'eau $\rightarrow$ \cle{solution aqueuse}.
\end{enumerate}
\textbf{Réflexe de rédaction :} toujours écrire une phrase complète : « Le soluté est ..., le solvant est ..., la solution est ... ».
\end{methode}

\begin{exoresolu}[L'eau de cuisson de Maman Alice]
Maman Alice prépare une sauce. Elle verse une cuillère de sel dans de l'eau chaude et agite : le sel disparaît complètement. Elle ajoute ensuite beaucoup de sel : une partie reste au fond de la casserole.
\begin{enumerate}
\item Identifier le soluté, le solvant et la solution dans le premier cas.
\item La solution obtenue est-elle aqueuse ? Justifier.
\item Que peut-on dire de la solution dans le deuxième cas ?
\end{enumerate}
\tcblower
\textbf{1. Identification des constituants.}
\begin{itemize}
\item L'espèce dissoute est le sel (\ce{NaCl}) : c'est le \cle{soluté}.
\item L'espèce qui dissout est l'eau (\ce{H2O}) : c'est le \cle{solvant}.
\item Le mélange homogène obtenu (eau salée) est la \cle{solution}.
\end{itemize}

\textbf{2. Nature de la solution.} Le solvant est l'eau. Par définition, toute solution dont le solvant est l'eau est une \cle{solution aqueuse}. L'eau salée est donc une solution aqueuse de chlorure de sodium.

\textbf{3. Deuxième cas.} Une partie du sel reste solide au fond : l'eau ne peut plus dissoudre davantage de sel. On dit que la solution est \cle{saturée} en sel. Le sel non dissous forme un \cle{dépôt}.

\textbf{Conclusion :} le sel est le soluté, l'eau est le solvant, l'eau salée est une solution aqueuse ; lorsque le sel ne se dissout plus entièrement, la solution est saturée.
\end{exoresolu}

\courssub{Solubilité et saturation}

\begin{definition}[Solubilité]
La \cle{solubilité} $s$ d'un soluté dans un solvant (à une température donnée) est la \textbf{masse maximale} de ce soluté que l'on peut dissoudre dans \textbf{1 L de ce solvant} pour obtenir une solution saturée. Elle s'exprime en \SI{}{g/L} (ou en \SI{}{g/100mL}).
\end{definition}

\begin{propriete}[États de saturation]
Soit $m$ la masse de soluté introduite dans un volume $V$ de solvant, et $s$ la solubilité. La masse maximale dissoluble est $m_{\max} = s \times V$ (avec $V$ en L).
\begin{itemize}
\item Si $m < m_{\max}$ : la solution est \cle{non saturée} (tout se dissout).
\item Si $m = m_{\max}$ : la solution est \cle{juste saturée} (limite de dissolution).
\item Si $m > m_{\max}$ : la solution est \cle{saturée} et il reste un dépôt de masse $m_{\text{dépôt}} = m - m_{\max}$.
\end{itemize}
\textbf{Remarque :} la solubilité dépend de la température (généralement, elle augmente avec la température pour les solides).
\end{propriete}

\begin{popfigure}
\begin{tikzpicture}
\begin{axis}[
    width=\linewidth, height=7cm,
    xlabel={Température (\si{\celsius})}, ylabel={Solubilité $s$ (\si{g/L})},
    domain=0:100, samples=100,
    axis lines=middle, grid=major,
    xmin=0, xmax=100, ymin=0, ymax=400,
    tick label style={font=\small},
    label style={font=\small},
]
\addplot[popcurve, thick, domain=0:100] {300 + x}; % Exemple courbe croissante type sel
\addlegendentry{Solubilité du sel (exemple)}
\addplot[poporange, dashed, thick] coordinates {(20,0) (20,320) (0,320)};
\addplot[poporange, dashed, thick] coordinates {(100,0) (100,400) (0,400)};
\node[poporange, anchor=south west, font=\small] at (20,320) {À 20°C : $s \approx 360$ g/L};
\node[poporange, anchor=south west, font=\small] at (100,400) {À 100°C : $s \approx 400$ g/L};
\end{axis}
\end{tikzpicture}
\legende{Variation de la solubilité du chlorure de sodium (\ce{NaCl}) en fonction de la température. On observe une légère augmentation. Pour d'autres solutés (ex. \ce{KNO3}), l'augmentation est bien plus forte.}
\end{popfigure}

\begin{methode}[Exploiter la solubilité d'un soluté]
\begin{enumerate}
\item Relever la solubilité $s$ dans les données (attention : elle dépend de la température).
\item Calculer la masse maximale dissoluble : $m_{\max} = s \times V$ ($V$ en litres de solvant).
\item Comparer la masse $m$ introduite à $m_{\max}$ :
\begin{itemize}
\item $m < m_{\max}$ : solution non saturée ;
\item $m = m_{\max}$ : solution juste saturée ;
\item $m > m_{\max}$ : solution saturée, dépôt de masse $m - m_{\max}$.
\end{itemize}
\item Conclure par une phrase citant les valeurs numériques.
\end{enumerate}
\end{methode}

\begin{exoresolu}[Solubilité du sel dans l'eau]
À \SI{20}{\celsius}, la solubilité du chlorure de sodium (\ce{NaCl}) dans l'eau est $s = \SI{360}{g/L}$. Un élève verse \SI{100}{g} de sel dans \SI{0,25}{L} d'eau à \SI{20}{\celsius} et agite longtemps.
\begin{enumerate}
\item Calculer la masse maximale de sel que l'on peut dissoudre dans ce volume d'eau.
\item Tout le sel se dissout-il ? Si non, calculer la masse du dépôt.
\end{enumerate}
\tcblower
\textbf{1. Masse maximale dissoluble.}
On applique $m_{\max} = s \times V$ avec $s = \SI{360}{g/L}$ et $V = \SI{0,25}{L}$ :
\[ m_{\max} = 360 \times 0,25 = \SI{90}{g}. \]
On peut dissoudre au maximum \SI{90}{g} de sel dans \SI{0,25}{L} d'eau à \SI{20}{\celsius}.

\textbf{2. Comparaison.}
La masse introduite est $m = \SI{100}{g}$. Or $100 > 90$, donc $m > m_{\max}$ : la solution est \cle{saturée} et tout le sel ne se dissout pas.

Masse du dépôt restant au fond :
\[ m_{\text{dépôt}} = m - m_{\max} = 100 - 90 = \SI{10}{g}. \]

\textbf{Conclusion :} on peut dissoudre au maximum \SI{90}{g} de sel dans \SI{0,25}{L} d'eau à \SI{20}{\celsius} ; il reste donc \SI{10}{g} de sel non dissous au fond du récipient : la solution est saturée.
\end{exoresolu}

\courssub{Concentration molaire $C = n/V$}

\begin{definition}[Concentration molaire]
La \cle{concentration molaire} $C$ d'une solution est la quantité de matière $n$ de soluté dissous par litre de solution :
\[ C = \frac{n}{V} \qquad \text{avec } C \text{ en \SI{}{mol/L}},\ n \text{ en \SI{}{mol}},\ V \text{ en \SI{}{L}}. \]
\end{definition}

\begin{propriete}[Formules dérivées]
À partir de $C = \dfrac{n}{V}$, on déduit deux formules utiles :
\[ n = C \times V \qquad \text{et} \qquad V = \frac{n}{C}. \]
\end{propriete}

\begin{methode}[Calculer une concentration molaire]
\begin{enumerate}
\item Identifier $n$ (en mol) et $V$ (volume de la \textbf{solution}) dans l'énoncé.
\item \textbf{Convertir le volume en litres} : $1~\text{L} = 1000~\text{mL}$.
\item Appliquer $C = \dfrac{n}{V}$ en posant le calcul avec les unités.
\item Donner le résultat en \SI{}{mol/L} et conclure par une phrase.
\end{enumerate}
\end{methode}

\begin{exoresolu}[Solution de glucose pour perfusion]
À l'hôpital, on prépare une solution de glucose (\ce{C6H12O6}) en dissolvant $n = \SI{0,05}{mol}$ de glucose dans de l'eau distillée pour obtenir un volume final $V = \SI{250}{mL}$.
\begin{enumerate}
\item Convertir le volume en litres.
\item Calculer la concentration molaire $C$.
\item Quelle quantité de matière pour préparer \SI{1}{L} de même concentration ?
\end{enumerate}
\tcblower
\textbf{1. Conversion.}
\[ V = \frac{250}{1000} = \SI{0,25}{L}. \]

\textbf{2. Concentration.}
\[ C = \frac{n}{V} = \frac{0,05}{0,25} = \SI{0,2}{mol/L}. \]

\textbf{3. Quantité pour 1 L.}
On utilise $n = C \times V$ avec $C = \SI{0,2}{mol/L}$ et $V = \SI{1}{L}$ :
\[ n = 0,2 \times 1 = \SI{0,2}{mol}. \]

\textbf{Conclusion :} la concentration est $C = \SI{0,2}{mol/L}$ ; pour \SI{1}{L} il faut \SI{0,2}{mol} de glucose.
\end{exoresolu}

\begin{methode}[Préparer une solution de concentration donnée (par dissolution)]
Pour préparer un volume $V$ de solution de concentration $C$ à partir d'un solide :
\begin{enumerate}
\item Calculer $n = C \times V$ ($V$ en L).
\item Si la masse molaire $M$ est connue, calculer la masse à peser : $m = n \times M$.
\item \textbf{Protocole au laboratoire :}
\begin{itemize}
\item Peser la masse $m$ de soluté (balance, coupelle).
\item Introduire le soluté dans une \cle{fiole jaugée} de volume $V$ (entonnoir, rinçage coupelle à l'eau distillée).
\item Ajouter de l'eau distillée aux $\frac{2}{3}$, boucher, agiter pour dissoudre.
\item Compléter à l'eau distillée \textbf{jusqu'au trait de jauge}, boucher, homogénéiser (retourner plusieurs fois).
\end{itemize}
\item \textbf{Sécurité :} blouse, lunettes, ne jamais goûter, étiqueter la fiole.
\end{enumerate}
\end{methode}

\begin{exoresolu}[Préparation d'une solution de chlorure de sodium]
Préparer $V = \SI{500}{mL}$ d'une solution de \ce{NaCl} de concentration $C = \SI{0,1}{mol/L}$. $M(\ce{NaCl}) = \SI{58,5}{g/mol}$.
\begin{enumerate}
\item Quantité de matière $n$ nécessaire.
\item Masse $m$ de sel à peser.
\item Décrire le protocole.
\end{enumerate}
\tcblower
\textbf{1. Quantité de matière.}
$V = \SI{500}{mL} = \SI{0,5}{L}$.
\[ n = C \times V = 0,1 \times 0,5 = \SI{0,05}{mol}. \]

\textbf{2. Masse à peser.}
\[ m = n \times M = 0,05 \times 58,5 = \SI{2,925}{g}. \]
On pèse \SI{2,93}{g} (précision balance 0,01 g) ou \SI{2,9}{g} (précision 0,1 g).

\textbf{3. Protocole.}
\begin{enumerate}
\item Peser \SI{2,93}{g} de \ce{NaCl} dans une coupelle.
\item Verser dans une fiole jaugée de \SI{500}{mL} (entonnoir), rincer la coupelle à l'eau distillée dans la fiole.
\item Remplir aux $\frac{2}{3}$ d'eau distillée, boucher, agiter jusqu'à dissolution totale.
\item Compléter à l'eau distillée jusqu'au trait de jauge (bas du ménisque au niveau du trait), boucher, homogénéiser.
\end{enumerate}

\textbf{Conclusion :} on obtient \SI{500}{mL} de solution aqueuse de \ce{NaCl} à \SI{0,1}{mol/L}.
\end{exoresolu}

\courssub{Notion de pH : solutions acides, neutres, basiques}

\begin{definition}[pH d'une solution]
Le \cle{pH} est un nombre sans unité (échelle de 0 à 14) qui caractérise l'acidité ou la basicité d'une solution aqueuse. Il se mesure avec du \cle{papier pH} (bandelettes) ou un \cle{pH-mètre}.
\end{definition}

\begin{propriete}[Classification des solutions aqueuses]
\begin{center}
\begin{tabular}{|c|c|c|c|}
\hline
\rowcolor{popblueL} \textbf{Domaine de pH} & \textbf{Qualificatif} & \textbf{Exemple courant (Cameroun)} & \textbf{Couleur papier pH (indicatif)} \\
\hline
$\text{pH} < 7$ & \cle{Acide} & Jus de citron, vinaigre, boisson gazeuse, acide batterie & Rouge / Orange / Jaune \\
\hline
$\text{pH} = 7$ & \cle{Neutre} & Eau pure, eau minérale (source), sérum physiologique & Vert \\
\hline
$\text{pH} > 7$ & \cle{Basique} & Eau de Javel (eau de Javel \emph{Croix}), eau savonneuse, lessive, cendres & Bleu / Violet \\
\hline
\end{tabular}
\end{center}
\textbf{Règle mnémotechnique :} « Plus on s'éloigne de 7, plus l'acidité (vers 0) ou la basicité (vers 14) est forte. »
\end{propriete}

\begin{popfigure}
\begin{tikzpicture}
% Échelle pH horizontale
\draw[popline, thick] (0,0) -- (14,0);
\foreach \x in {0,1,...,14} {
    \draw (\x,0.1) -- (\x,-0.1) node[below, font=\tiny] {\x};
}
% Zones colorées
\fill[poppink!50] (0,-0.5) rectangle (7,0.5);
\fill[popgreen!30] (7,-0.5) rectangle (7.02,0.5); %Trait neutre
\fill[popblue!50] (7,-0.5) rectangle (14,0.5);
% Légendes zones
\node[popdark, font=\bfseries\small] at (3.5, 1) {ACIDE};
\node[popdark, font=\bfseries\small] at (7, 1) {NEUTRE};
\node[popdark, font=\bfseries\small] at (10.5, 1) {BASIQUE};
% Exemples
\node[poppink, font=\small, align=center] at (1, -1.2) {Acide batterie\\pH $\approx$ 1};
\node[poppink, font=\small, align=center] at (2.5, -1.2) {Jus citron\\pH $\approx$ 2,4};
\node[poppink, font=\small, align=center] at (4, -1.2) {Vinaigre\\pH $\approx$ 3};
\node[popgreen, font=\small, align=center] at (7, -1.2) {Eau pure\\pH = 7};
\node[popblue, font=\small, align=center] at (9.5, -1.2) {Eau savonneuse\\pH $\approx$ 9,5};
\node[popblue, font=\small, align=center] at (12, -1.2) {Eau de Javel\\pH $\approx$ 12};
% Flèches force
\draw[<->, thick, poppink] (0,-1.8) -- (7,-1.8) node[midway, below, font=\tiny] {Acidité croissante};
\draw[<->, thick, popblue] (7,-1.8) -- (14,-1.8) node[midway, below, font=\tiny] {Basicité croissante};
\end{tikzpicture}
\legende{Échelle du pH avec exemples de la vie courante au Cameroun.}
\end{popfigure}

\begin{methode}[Déterminer le caractère acide, neutre ou basique]
\begin{enumerate}
\item Lire la valeur du pH (papier pH ou pH-mètre).
\item Comparer à 7 :
\begin{itemize}
\item $\text{pH} < 7$ $\rightarrow$ \cle{acide} (plus c'est petit, plus c'est acide).
\item $\text{pH} = 7$ $\rightarrow$ \cle{neutre}.
\item $\text{pH} > 7$ $\rightarrow$ \cle{basique} (plus c'est grand, plus c'est basique).
\end{itemize}
\item Citer un exemple concret.
\item \textbf{Sécurité :} solutions très acides (pH $\le$ 2) ou très basiques (pH $\ge$ 12) = \cle{corrosives} $\rightarrow$ gants + lunettes obligatoires.
\end{enumerate}
\end{methode}

\begin{exoresolu}[Le pH des produits de la maison]
On mesure le pH de liquides usuels :

\begin{center}
\begin{tabularx}{\linewidth}{|l|X|X|X|X|}
\hline
\rowcolor{popblueL} Liquide & Jus de citron & Eau minérale & Eau savonneuse & Vinaigre \\
\hline
pH & 2,4 & 7,0 & 9,5 & 3,1 \\
\hline
\end{tabularx}
\end{center}

\begin{enumerate}
\item Classer ces liquides (acide, neutre, basique).
\item Quel est le plus acide ? Justifier.
\item Pourquoi porter des gants pour manipuler un détergent très basique ?
\end{enumerate}
\tcblower
\textbf{1. Classement.}
\begin{itemize}
\item Jus de citron : $\text{pH} = 2,4 < 7$ $\rightarrow$ \cle{acide}.
\item Eau minérale : $\text{pH} = 7,0$ $\rightarrow$ \cle{neutre}.
\item Eau savonneuse : $\text{pH} = 9,5 > 7$ $\rightarrow$ \cle{basique}.
\item Vinaigre : $\text{pH} = 3,1 < 7$ $\rightarrow$ \cle{acide}.
\end{itemize}

\textbf{2. Le plus acide.}
Plus le pH est \textbf{petit} (inférieur à 7), plus la solution est acide. $2,4 < 3,1$, donc le \cle{jus de citron} est le plus acide.

\textbf{3. Sécurité.}
Une solution très basique (pH très $> 7$, ex. soude concentrée, eau de Javel pure) est \cle{corrosive} : elle attaque la peau et les yeux (brûlures chimiques). Le port de gants et de lunettes est indispensable, comme pour les acides concentrés.

\textbf{Conclusion :} citron et vinaigre acides, eau neutre, eau savonneuse basique ; le citron est le plus acide (pH le plus faible).
\end{exoresolu}

\courssub{Sécurité et applications concrètes au Cameroun}

\begin{savaistu}[Le SRO : un mélange vital]
Au Cameroun, comme partout en zone tropicale, la diarrhée aiguë chez l'enfant est une cause majeure de mortalité. Le traitement de référence est le \textbf{SRO (Sérum de Réhydratation Orale)}. C'est une \textbf{solution aqueuse} précise dont la composition (concentrations molaires en ions \ce{Na+}, \ce{K+}, \ce{Cl-}, citrate, glucose) est normée par l'OMS.
\begin{itemize}
\item \textbf{Préparation :} on dissout un sachet (poudre = solutés) dans \textbf{1 L d'eau potable bouillie puis refroidie} (solvant). Respecter le volume est crucial : trop concentré $\rightarrow$ hypernatrémie (danger) ; trop dilué $\rightarrow$ inefficace.
\item \textbf{Lien au cours :} c'est une application directe de la notion de concentration molaire ($C = n/V$) et de préparation de solution (fiole jaugée $\leftrightarrow$ bouteille 1 L graduée).
\end{itemize}
\end{savaistu}

\begin{savaistu}[Eau de Javel et désinfection de l'eau]
L'\textbf{eau de Javel} (solution d'hypochlorite de sodium \ce{NaClO}) est une solution \textbf{basique} (pH $\approx$ 11-12). Elle est utilisée pour :
\begin{itemize}
\item Désinfecter l'eau de boisson (quelques gouttes par litre, attendre 30 min).
\item Blanchir le linge, nettoyer les sols.
\end{itemize}
\textbf{Attention :} ne jamais mélanger eau de Javel et produit acide (vinaigre, détartrant, WC) $\rightarrow$ dégage du \textbf{dichlore \ce{Cl2}} (gaz toxique, vert-jaune, suffocant). C'est une réaction chimique dangereuse : $\ce{ClO- + 2H+ + Cl- -> Cl2 ^ + H2O}$.
\end{savaistu}

\begin{savaistu}[Batterie moto/voiture : l'électrolyte]
La batterie au plomb (moto-taxi, taxi, groupe électrogène) contient de l'\textbf{acide sulfurique} \ce{H2SO4} dilué (solution aqueuse \textbf{très acide}, pH $\approx$ 0-1, concentration $\approx$ \SI{4}{mol/L} à \SI{5}{mol/L}).
\begin{itemize}
\item \textbf{Danger :} très corrosif (brûlures graves), dégage de l'hydrogène \ce{H2} (explosif) à la charge.
\item \textbf{Entretien :} vérifier le niveau d'électrolyte (ajouter de l'eau \textbf{distillée} seulement, jamais d'acide), porter gants et lunettes.
\end{itemize}
\end{savaistu}

\begin{experience}[TP : Préparer une solution et mesurer son pH]
\textbf{Objectif :} Préparer \SI{100}{mL} d'une solution de \ce{NaCl} à \SI{0,1}{mol/L} et mesurer le pH de solutions usuelles.
\textbf{Matériel :} Fiole jaugée \SI{100}{mL}, balance (précision 0,01 g), entonnoir, bécher, baguette de verre, eau distillée, \ce{NaCl} solide, papier pH (ou pH-mètre), tubes à essai, solutions test (vinaigre, eau savonneuse, jus citron, eau robinet).
\textbf{Protocole :}
\begin{enumerate}
\item Calculer la masse $m$ de \ce{NaCl} nécessaire ($M = \SI{58,5}{g/mol}$). Peser cette masse.
\item Réaliser la dissolution dans la fiole jaugée (méthode 4).
\item Verser \SI{5}{mL} de chaque solution test dans des tubes étiquetés.
\item Tremper une bandelette pH dans chaque tube, comparer à l'échelle colorée, noter le pH.
\item Classer les solutions (acide/neutre/basique).
\end{enumerate}
\textbf{Interprétation attendue :} La solution de \ce{NaCl} est neutre (pH $\approx$ 7). Le vinaigre et le citron sont acides. L'eau savonneuse est basique. L'eau du robinet (ENEO) est souvent légèrement basique (traitement chlore).
\textbf{Sécurité :} Lunettes obligatoires. Ne pas goûter. Eliminer les solutions dans l'évier avec beaucoup d'eau.
\end{experience}

\begin{attention}[Les 5 erreurs qui coûtent des points au BEPC]
\begin{enumerate}
\item \textbf{Confondre soluté et solvant.} Le solvant dissout (souvent l'eau, grande quantité) ; le soluté est dissous. Écrire « l'eau est le soluté » = 0 point.
\item \textbf{Oublier de convertir le volume en litres.} Dans $C = \dfrac{n}{V}$, $V$ \textbf{doit} être en \SI{}{L}. Utiliser \SI{250}{mL} directement donne un résultat 1000 fois trop grand. Écrire : $V = \SI{250}{mL} = \SI{0,25}{L}$.
\item \textbf{Résultat sans unité ou mauvaise unité.} Concentration molaire = \SI{}{mol/L}. Solubilité = \SI{}{g/L}. Un nombre seul n'est pas une réponse.
\item \textbf{Inverser l'échelle du pH.} $\text{pH} < 7$ = acide ; $\text{pH} > 7$ = basique. Plus le pH est \textbf{petit}, plus c'est \textbf{acide} (pas l'inverse).
\item \textbf{Oublier la phrase de conclusion et les justifications.} Au BEPC, il faut : \textbf{Formule} $\rightarrow$ \textbf{Remplacement numérique avec unités} $\rightarrow$ \textbf{Résultat + unité} $\rightarrow$ \textbf{Phrase de conclusion}. Pour la saturation : comparer $m$ et $m_{\max}$ explicitement.
\end{enumerate}
\end{attention}

\begin{exercice}[Entraînement BEPC - Solutions aqueuses]
\begin{enumerate}
\item On dissout \SI{5,85}{g} de chlorure de sodium (\ce{NaCl}, $M = \SI{58,5}{g/mol}$) dans de l'eau pour obtenir \SI{250}{mL} de solution. Calculer la concentration molaire $C$ de cette solution.
\item À \SI{20}{\celsius}, la solubilité du sulfate de cuivre(II) (\ce{CuSO4}) est $s = \SI{230}{g/L}$. On introduit \SI{30}{g} de \ce{CuSO4} dans \SI{100}{mL} d'eau. La solution est-elle saturée ? Justifier par un calcul. Si oui, quelle est la masse du dépôt ?
\item On mesure le pH de trois solutions : A : pH = 13 ; B : pH = 6,5 ; C : pH = 7,0. Classer ces solutions par ordre d'acidité croissante. Laquelle est la plus corrosive pour la peau ?
\item \textbf{Pratique :} Tu dois préparer \SI{500}{mL} d'une solution de glucose (\ce{C6H12O6}, $M = \SI{180}{g/mol}$) à \SI{0,2}{mol/L}. Indique la masse à peser et décris les 4 étapes clés de la préparation en fiole jaugée.
\end{enumerate}
\end{exercice}

\begin{corrige}[Exercice Entraînement BEPC]
\begin{enumerate}
\item \textbf{Données :} $m = \SI{5,85}{g}$, $M = \SI{58,5}{g/mol}$, $V = \SI{250}{mL} = \SI{0,25}{L}$.
$n = \frac{m}{M} = \frac{5,85}{58,5} = \SI{0,1}{mol}$.
$C = \frac{n}{V} = \frac{0,1}{0,25} = \SI{0,4}{mol/L}$.
\textbf{Conclusion :} $C = \SI{0,4}{mol/L}$.

\item \textbf{Données :} $s = \SI{230}{g/L}$, $m = \SI{30}{g}$, $V_{\text{eau}} = \SI{100}{mL} = \SI{0,1}{L}$.
$m_{\max} = s \times V = 230 \times 0,1 = \SI{23}{g}$.
Comparaison : $m = \SI{30}{g} > m_{\max} = \SI{23}{g}$.
La solution est \textbf{saturée}.
Masse dépôt $= m - m_{\max} = 30 - 23 = \SI{7}{g}$.
\textbf{Conclusion :} solution saturée, dépôt de \SI{7}{g}.

\item \textbf{Classement acidité croissante (du plus acide au moins acide / plus basique) :}
Plus le pH est grand, moins c'est acide.
B (pH 6,5, légèrement acide) $<$ C (pH 7, neutre) $<$ A (pH 13, très basique).
\textbf{Ordre :} B $<$ C $<$ A.
\textbf{Plus corrosive :} Solution A (pH 13, très basique, comme la soude ou l'eau de Javel concentrée). Les solutions très basiques sont aussi dangereuses que les très acides.

\item \textbf{Données :} $C = \SI{0,2}{mol/L}$, $V = \SI{500}{mL} = \SI{0,5}{L}$, $M = \SI{180}{g/mol}$.
$n = C \times V = 0,2 \times 0,5 = \SI{0,1}{mol}$.
$m = n \times M = 0,1 \times 180 = \SI{18}{g}$.
\textbf{Masse à peser :} \SI{18,0}{g} (balance 0,1 g) ou \SI{18,00}{g} (balance 0,01 g).
\textbf{Étapes préparation :}
1. Peser \SI{18}{g} de glucose dans une coupelle.
2. Verser dans fiole jaugée \SI{500}{mL} (entonnoir + rinçage coupelle).
3. Ajouter eau distillée aux 2/3, boucher, agiter dissolution.
4. Compléter eau distillée jusqu'au trait de jauge, boucher, homogénéiser.
\end{enumerate}
\end{corrige}

\newpage

\coursec{Exercices}

\courssub{Je vérifie mes connaissances}

\begin{exercice}[QCM : soluté et solvant]
Pour chaque solution de la vie courante, identifie le \cle{soluté} et le \cle{solvant}. Entoure la bonne réponse.

\begin{enumerate}
\item Dans un verre d'eau sucrée, le sucre est :
\begin{enumerate}
\item le solvant \item le soluté \item la solution
\end{enumerate}
\item Dans une tisane (feuilles infusées dans l'eau chaude), l'eau joue le rôle de :
\begin{enumerate}
\item soluté \item solvant \item précipité
\end{enumerate}
\item Dans l'eau de Javel diluée utilisée pour laver les sols, l'eau de Javel concentrée ajoutée à l'eau est :
\begin{enumerate}
\item le solvant \item le soluté \item le mélange hétérogène
\end{enumerate}
\item Une solution dans laquelle on ne peut plus dissoudre de soluté est dite :
\begin{enumerate}
\item diluée \item saturée \item neutre
\end{enumerate}
\end{enumerate}
\end{exercice}

\begin{exercice}[Vrai ou faux ? Justifie]
Réponds par Vrai ou Faux, puis justifie ta réponse en une phrase.

\begin{enumerate}
\item Le sel se dissout dans l'eau : le sel est le solvant.
\item La solubilité d'un sel dans l'eau augmente généralement quand la température augmente.
\item Une solution de concentration $C = \SI{0,5}{mol/L}$ contient $0,5$ mol de soluté dans chaque litre de solution.
\item Une solution de pH égal à 9 est acide.
\item L'eau pure a un pH égal à 7 : elle est neutre.
\end{enumerate}
\end{exercice}

\begin{exercice}[Définitions à compléter]
Recopie et complète les phrases suivantes avec les mots : \emph{solvant}, \emph{soluté}, \emph{solution}, \emph{solubilité}, \emph{concentration molaire}.

\begin{enumerate}
\item Le \ldots\ldots\ldots\ldots est la substance que l'on dissout.
\item Le \ldots\ldots\ldots\ldots est le liquide dans lequel on dissout une substance.
\item Le mélange homogène obtenu après dissolution s'appelle une \ldots\ldots\ldots\ldots
\item La \ldots\ldots\ldots\ldots est la masse maximale de soluté que l'on peut dissoudre dans un litre de solvant à une température donnée.
\item La \ldots\ldots\ldots\ldots d'une solution est le quotient de la quantité de matière de soluté par le volume de la solution : $C = \dfrac{n}{V}$.
\end{enumerate}
\end{exercice}

\begin{exercice}[Calculs directs]
\begin{enumerate}
\item On dissout $n = \SI{0,2}{mol}$ de glucose dans de l'eau de façon à obtenir $V = \SI{500}{mL}$ de solution. Calcule la concentration molaire $C$ de cette solution.
\item Une solution de sel de cuisine a une concentration $C = \SI{1}{mol/L}$. Quelle quantité de matière de sel contient un volume $V = \SI{250}{mL}$ de cette solution ?
\item On donne les pH de trois boissons : jus d'orange (pH = 4), eau minérale (pH = 7), eau savonneuse (pH = 9). Indique pour chacune si elle est acide, neutre ou basique.
\end{enumerate}
\end{exercice}

\courssub{Je m'entraîne}

\begin{exercice}[Le sérum physiologique]
Sur l'étiquette d'une poche de sérum physiologique utilisée à l'hôpital, on lit : « chlorure de sodium, $C = \SI{0,154}{mol/L}$ ». On donne la masse molaire du chlorure de sodium : $M(\ce{NaCl}) = \SI{58,5}{g/mol}$.

\begin{enumerate}
\item Quelle quantité de matière de chlorure de sodium contient une poche de volume $V = \SI{500}{mL}$ ?
\item En déduire la masse de sel contenue dans cette poche.
\item Vérifie que cette masse est proche de \SI{4,5}{g}. Pourquoi appelle-t-on parfois ce sérum « solution à 9 pour mille » ?
\end{enumerate}
\end{exercice}

\begin{exercice}[Courbe de solubilité du nitrate de potassium]
La solubilité du nitrate de potassium \ce{KNO3} dans l'eau vaut environ \SI{64}{g} par litre d'eau à \SI{40}{\celsius} et \SI{32}{g} par litre d'eau à \SI{20}{\celsius}.

\begin{enumerate}
\item Quelle masse maximale de \ce{KNO3} peut-on dissoudre dans \SI{300}{mL} d'eau à \SI{40}{\celsius} ?
\item On laisse refroidir cette solution saturée jusqu'à \SI{20}{\celsius}. Quelle masse de \ce{KNO3} reste dissoute ?
\item Que devient le reste du sel ? Comment s'appelle ce phénomène ?
\end{enumerate}
\end{exercice}

\begin{exercice}[Classement selon le pH]
On mesure le pH de quatre solutions de la vie courante : jus de citron (pH = 2), eau pure (pH = 7), eau savonneuse (pH = 10), vinaigre (pH = 3).

\begin{enumerate}
\item Classe ces quatre solutions par acidité croissante.
\item Indique celles qui sont acides, celle qui est neutre et celle qui est basique.
\item Lequel de ces produits faut-il manipuler avec le plus de précautions ? Justifie.
\end{enumerate}
\end{exercice}

\begin{exercice}[Préparation d'une solution pour la maman]
Pour préparer une solution de réhydratation maison, une maman dissout $n = \SI{0,05}{mol}$ de sel dans $V = \SI{1}{L}$ d'eau bouillie refroidie.

\begin{enumerate}
\item Calcule la concentration molaire de la solution obtenue.
\item Elle prélève ensuite \SI{200}{mL} de cette solution pour son enfant. Quelle est la concentration du prélèvement ? Justifie ta réponse.
\item Quelle quantité de matière de sel l'enfant reçoit-il avec ces \SI{200}{mL} ?
\end{enumerate}
\end{exercice}

\courssub{Je me prépare au Bac}

\begin{exercice}[Type BEPC : la saumure du poisson fumé]
Pour conserver le poisson fumé, une vendeuse du marché dissout $m = \SI{5,85}{g}$ de chlorure de sodium \ce{NaCl} dans de l'eau et obtient $V = \SI{200}{mL}$ de solution. On donne : $M(\ce{NaCl}) = \SI{58,5}{g/mol}$.

\begin{enumerate}
\item Définis les termes soluté et solvant, puis identifie-les dans cette situation.
\item Calcule la quantité de matière $n$ de chlorure de sodium dissous.
\item Convertis le volume $V$ en litres, puis calcule la concentration molaire $C$ de la solution.
\item La vendeuse ajoute \SI{300}{mL} d'eau pure à sa solution. La quantité de matière de sel change-t-elle ? La concentration change-t-elle ? Calcule la nouvelle concentration.
\item On dit que la solution obtenue à la question 3 est « décimolaire » si $C = \SI{0,1}{mol/L}$. La solution de la vendeuse est-elle décimolaire ?
\end{enumerate}
\end{exercice}

\begin{exercice}[Type BEPC : expérience et mesures de pH]
Au laboratoire du collège, un groupe d'élèves de 3ème réalise les expériences suivantes :

\textbf{Expérience 1 :} ils dissolvent progressivement du sel dans \SI{100}{mL} d'eau à température ambiante en agitant. À partir d'une certaine masse ajoutée, le sel ne se dissout plus et se dépose au fond du bécher.

\textbf{Expérience 2 :} ils mesurent le pH de trois liquides avec du papier pH : vinaigre (pH = 3), eau de robinet (pH = 7), solution de lessive (pH = 11).

\begin{enumerate}
\item Comment appelle-t-on l'opération réalisée dans l'expérience 1 ?
\item Comment qualifie-t-on la solution obtenue lorsque le sel ne se dissout plus ? Définis ce terme.
\item Propose une méthode permettant de dissoudre davantage de sel dans cette solution. Justifie.
\item Pour chaque liquide de l'expérience 2, précise s'il est acide, neutre ou basique. Justifie à l'aide de l'échelle de pH.
\item La lessive a un pH très élevé. Cite deux précautions à prendre pour la manipuler.
\end{enumerate}
\end{exercice}

\begin{exercice}[Type BEPC : danger des mélanges à la maison]
Dans une maison à Yaoundé, un élève veut nettoyer les toilettes. Il verse de l'eau de Javel (solution basique contenant des ions hypochlorite \ce{ClO-}) puis un détartrant acide (contenant de l'acide chlorhydrique). Le mélange dégage un gaz toxique, le dichlore \ce{Cl2}, selon une transformation chimique qui peut provoquer de graves irritations des yeux et des poumons.

\begin{enumerate}
\item Le détartrant a un pH voisin de 2 et l'eau de Javel un pH voisin de 11. Indique la nature (acide ou basique) de chaque produit.
\item Explique pourquoi il est dangereux de mélanger ces deux produits.
\item Rédige une consigne de sécurité claire à afficher dans la maison concernant l'utilisation des produits d'entretien.
\item Un flacon de détartrant porte la mention « $C = \SI{1}{mol/L}$ d'acide ». Calcule la quantité de matière d'acide contenue dans un volume $V = \SI{250}{mL}$ de détartrant.
\item On dilue ce prélèvement en ajoutant de l'eau jusqu'à obtenir \SI{1}{L} de solution. Calcule la nouvelle concentration molaire. Le pH de la solution diluée est-il plus proche de 2 ou de 7 ? Justifie.
\end{enumerate}
\end{exercice}

\newpage

\coursec{Corrigés des exercices}

\begin{corrige}[Exercice 1 — QCM : soluté et solvant]
\begin{enumerate}
\item \textbf{Réponse b : le soluté.} Le sucre est la substance dissoute dans l'eau ; c'est donc le \cle{soluté}. L'eau est le \cle{solvant} et le mélange obtenu est la \cle{solution}.
\item \textbf{Réponse b : solvant.} L'eau chaude est le liquide qui dissout les substances contenues dans les feuilles de thé : elle joue le rôle de \cle{solvant}.
\item \textbf{Réponse b : le soluté.} L'eau de Javel concentrée est ajoutée en petite quantité dans une grande quantité d'eau : c'est la substance dissoute, donc le \cle{soluté}. L'eau est le \cle{solvant}.
\item \textbf{Réponse b : saturée.} Par définition, une solution est \cle{saturée} lorsqu'on ne peut plus y dissoudre de soluté à la température considérée.
\end{enumerate}
\end{corrige}

\begin{corrige}[Exercice 2 — Vrai ou faux ? Justifie]
\begin{enumerate}
\item \textbf{Faux.} Le sel est la substance dissoute : c'est le \cle{soluté}. Le solvant est l'eau.
\item \textbf{Vrai.} Pour la plupart des solides, la \cle{solubilité} dans l'eau augmente lorsque la température augmente : c'est pourquoi on chauffe pour dissoudre davantage de sel ou de sucre.
\item \textbf{Vrai.} Par définition, $C = \dfrac{n}{V}$ : une concentration de \SI{0,5}{mol/L} signifie qu'il y a \SI{0,5}{mol} de soluté par litre de solution.
\item \textbf{Faux.} Un pH supérieur à 7 caractérise une solution \cle{basique}. Une solution acide a un pH inférieur à 7.
\item \textbf{Vrai.} À pH = 7, la solution est neutre : c'est le cas de l'eau pure.
\end{enumerate}
\end{corrige}

\begin{corrige}[Exercice 3 — Définitions à compléter]
\begin{enumerate}
\item Le \textbf{soluté} est la substance que l'on dissout.
\item Le \textbf{solvant} est le liquide dans lequel on dissout une substance.
\item Le mélange homogène obtenu après dissolution s'appelle une \textbf{solution}.
\item La \textbf{solubilité} est la masse maximale de soluté que l'on peut dissoudre dans un litre de solvant à une température donnée.
\item La \textbf{concentration molaire} d'une solution est le quotient de la quantité de matière de soluté par le volume de la solution : $C = \dfrac{n}{V}$.
\end{enumerate}
\end{corrige}

\begin{corrige}[Exercice 4 — Calculs directs]
\begin{enumerate}
\item On convertit le volume : $V = \SI{500}{mL} = \SI{0,5}{L}$. On applique la formule :
\[ C = \frac{n}{V} = \frac{\SI{0,2}{mol}}{\SI{0,5}{L}} = \SI{0,4}{mol/L}. \]
La concentration molaire de la solution de glucose est $\boxed{C = \SI{0,4}{mol/L}}$.
\item On convertit : $V = \SI{250}{mL} = \SI{0,25}{L}$. De $C = \dfrac{n}{V}$ on tire $n = C \times V$ :
\[ n = \SI{1}{mol/L} \times \SI{0,25}{L} = \SI{0,25}{mol}. \]
Le prélèvement contient $\boxed{n = \SI{0,25}{mol}}$ de sel.
\item Jus d'orange (pH = 4 $< 7$) : \textbf{acide}. Eau minérale (pH = 7) : \textbf{neutre}. Eau savonneuse (pH = 9 $> 7$) : \textbf{basique}.
\end{enumerate}
\end{corrige}

\begin{corrige}[Exercice 5 — Le sérum physiologique]
\begin{enumerate}
\item On convertit le volume : $V = \SI{500}{mL} = \SI{0,5}{L}$. Alors :
\[ n = C \times V = \SI{0,154}{mol/L} \times \SI{0,5}{L} = \SI{0,077}{mol}. \]
La poche contient $\boxed{n = \SI{0,077}{mol}}$ de chlorure de sodium.
\item On utilise $m = n \times M$ (avec $M_{\ce{NaCl}} = \SI{58,5}{g/mol}$) :
\[ m = \SI{0,077}{mol} \times \SI{58,5}{g/mol} \approx \SI{4,5}{g}. \]
La masse de sel dans la poche est $\boxed{m \approx \SI{4,5}{g}}$.
\item On a bien $m \approx \SI{4,5}{g}$ pour \SI{500}{mL}, soit \SI{9}{g} pour un litre. Or \SI{9}{g} de sel pour \SI{1000}{g} (environ 1 L) de solution représente 9 grammes pour mille : d'où le nom de « solution à 9 pour mille » (9 \textperthousand).
\end{enumerate}
\end{corrige}

\begin{corrige}[Exercice 6 — Courbe de solubilité du nitrate de potassium]
\begin{enumerate}
\item D'après la courbe fournie, à \SI{40}{\celsius}, la solubilité est de \SI{64}{g} de \ce{KNO3} par litre d'eau. Pour un volume d'eau $V = \SI{300}{mL} = \SI{0,3}{L}$ :
\[ m_{\max} = 64 \times 0,3 = \SI{19,2}{g}. \]
On peut dissoudre au maximum $\boxed{\SI{19,2}{g}}$ de \ce{KNO3}.
\item À \SI{20}{\celsius}, la solubilité n'est plus que \SI{32}{g/L} (d'après la courbe). La masse qui reste dissoute dans \SI{0,3}{L} d'eau est :
\[ m = 32 \times 0,3 = \SI{9,6}{g}. \]
Il reste $\boxed{\SI{9,6}{g}}$ de \ce{KNO3} dissous.
\item Le reste du sel, soit $19,2 - 9,6 = \SI{9,6}{g}$, ne peut plus rester dissous : il redevient solide et se dépose au fond du récipient. Ce phénomène s'appelle la \cle{cristallisation} (ou précipitation) du sel par refroidissement.
\end{enumerate}
\end{corrige}

\begin{corrige}[Exercice 7 — Classement selon le pH]
\begin{enumerate}
\item Plus le pH est petit, plus la solution est acide. L'ordre d'acidité croissante (du moins acide au plus acide) est donc :
\[ \text{eau savonneuse (10)} < \text{eau pure (7)} < \text{vinaigre (3)} < \text{jus de citron (2)}. \]
\item Solutions \textbf{acides} (pH $< 7$) : jus de citron et vinaigre. Solution \textbf{neutre} (pH = 7) : eau pure. Solution \textbf{basique} (pH $> 7$) : eau savonneuse.
\item Le jus de citron est le plus acide (pH = 2), mais c'est un produit alimentaire. En pratique, c'est l'\textbf{eau savonneuse} (pH = 10) qui demande le plus de précautions parmi ces produits d'usage ménager, car les bases fortes sont corrosives pour la peau et les yeux : il faut éviter le contact avec la peau et se laver les mains après usage.
\end{enumerate}
\end{corrige}

\begin{corrige}[Exercice 8 — Préparation d'une solution pour la maman]
\begin{enumerate}
\item On applique la formule :
\[ C = \frac{n}{V} = \frac{\SI{0,05}{mol}}{\SI{1}{L}} = \SI{0,05}{mol/L}. \]
La concentration de la solution est $\boxed{C = \SI{0,05}{mol/L}}$.
\item La concentration du prélèvement est \textbf{la même} : $C = \SI{0,05}{mol/L}$. En effet, une solution est un mélange \cle{homogène} : prélever une partie de la solution ne change pas sa concentration, seulement la quantité de matière prélevée.
\item On convertit : $V = \SI{200}{mL} = \SI{0,2}{L}$. Alors :
\[ n = C \times V = \SI{0,05}{mol/L} \times \SI{0,2}{L} = \SI{0,01}{mol}. \]
L'enfant reçoit $\boxed{n = \SI{0,01}{mol}}$ de sel avec ces \SI{200}{mL}.
\end{enumerate}
\end{corrige}

\begin{corrige}[Exercice 9 — Type BEPC : la saumure du poisson fumé]
\begin{enumerate}
\item Le \cle{soluté} est la substance que l'on dissout ; le \cle{solvant} est le liquide dans lequel on dissout cette substance. Ici, le soluté est le chlorure de sodium \ce{NaCl} (le sel) et le solvant est l'eau.
\item On applique $n = \dfrac{m}{M}$ (avec $M_{\ce{NaCl}} = \SI{58,5}{g/mol}$) :
\[ n = \frac{\SI{5,85}{g}}{\SI{58,5}{g/mol}} = \SI{0,1}{mol}. \]
La quantité de matière de sel dissous est $\boxed{n = \SI{0,1}{mol}}$.
\item Conversion : $V = \SI{200}{mL} = \SI{0,2}{L}$. Alors :
\[ C = \frac{n}{V} = \frac{\SI{0,1}{mol}}{\SI{0,2}{L}} = \SI{0,5}{mol/L}. \]
La concentration de la saumure est $\boxed{C = \SI{0,5}{mol/L}}$.
\item En ajoutant de l'eau pure, on n'ajoute pas de sel : la \textbf{quantité de matière de sel ne change pas} ($n = \SI{0,1}{mol}$). En revanche, le volume augmente, donc la \textbf{concentration diminue}. Le nouveau volume est :
\[ V' = 200 + 300 = \SI{500}{mL} = \SI{0,5}{L}. \]
La nouvelle concentration est :
\[ C' = \frac{n}{V'} = \frac{\SI{0,1}{mol}}{\SI{0,5}{L}} = \SI{0,2}{mol/L}. \]
\item Non : la solution de la question 3 a une concentration $C = \SI{0,5}{mol/L}$, alors qu'une solution décimolaire a $C = \SI{0,1}{mol/L}$. La saumure est cinq fois plus concentrée qu'une solution décimolaire.
\end{enumerate}
\textbf{Barème indicatif (10 points) :} définitions 2 pts ; calcul de $n$ 2 pts ; conversion et calcul de $C$ 2 pts ; dilution (quantité conservée, nouveau $C'$) 3 pts ; comparaison décimolaire 1 pt.

\textbf{Points de vigilance :} toujours convertir les volumes en litres avant d'appliquer $C = \dfrac{n}{V}$ ; lors d'une dilution, citer explicitement que la quantité de matière de soluté est conservée ; ne pas confondre ajout d'eau (dilution) et ajout de sel.
\end{corrige}

\begin{corrige}[Exercice 10 — Type BEPC : expérience et mesures de pH]
\begin{enumerate}
\item L'opération réalisée dans l'expérience 1 est une \cle{dissolution} du sel dans l'eau.
\item Lorsque le sel ne se dissout plus, la solution est dite \cle{saturée} : une solution saturée est une solution dans laquelle on ne peut plus dissoudre de soluté à la température considérée ; l'excès de soluté se dépose au fond.
\item On peut \textbf{chauffer la solution} : la solubilité de la plupart des solides (comme le nitrate de potassium) dans l'eau augmente avec la température, donc une partie du sel déposé se dissout à chaud. (On peut aussi ajouter de l'eau, c'est-à-dire augmenter la quantité de solvant.)
\item Vinaigre (pH = 3 $< 7$) : \textbf{acide}. Eau de robinet (pH = 7) : \textbf{neutre}. Solution de lessive (pH = 11 $> 7$) : \textbf{basique}. Sur l'échelle de pH, une solution est acide si pH $< 7$, neutre si pH = 7 et basique si pH $> 7$.
\item Pour manipuler la lessive (produit très basique et corrosif), il faut par exemple : porter des \textbf{gants} de protection ; éviter tout contact avec les \textbf{yeux} (porter des lunettes de protection) ; ne jamais la mélanger avec un autre produit ; la conserver hors de portée des enfants. (Deux réponses suffisent.)
\end{enumerate}
\textbf{Barème indicatif (10 points) :} question 1 : 1 pt ; question 2 : 2 pts ; question 3 : 2 pts ; question 4 : 3 pts ; question 5 : 2 pts.

\textbf{Points de vigilance :} la justification de la question 4 doit citer la comparaison du pH à 7 ; pour la question 3, la justification attendue est l'augmentation de la solubilité avec la température ; ne pas écrire « le sel fond » : il se \emph{dissout}.
\end{corrige}

\begin{corrige}[Exercice 11 — Type BEPC : danger des mélanges à la maison]
\begin{enumerate}
\item Le détartrant (pH $\approx 2 < 7$) est \textbf{acide} ; l'eau de Javel (pH $\approx 11 > 7$) est \textbf{basique}.
\item Le mélange de ces deux produits provoque une \cle{transformation chimique} qui dégage du dichlore \ce{Cl2}, un gaz \textbf{toxique} qui irrite gravement les yeux et les poumons. C'est pourquoi il ne faut jamais mélanger un produit acide avec de l'eau de Javel.
\item Consigne de sécurité : « \textbf{Ne jamais mélanger l'eau de Javel avec un détartrant ou tout autre produit acide. Utiliser un seul produit à la fois, porter des gants et aérer la pièce.} »
\item On convertit : $V = \SI{250}{mL} = \SI{0,25}{L}$. Alors :
\[ n = C \times V = \SI{1}{mol/L} \times \SI{0,25}{L} = \SI{0,25}{mol}. \]
Le prélèvement contient $\boxed{n = \SI{0,25}{mol}}$ d'acide.
\item Lors de la dilution, la quantité de matière d'acide est conservée : $n = \SI{0,25}{mol}$ dans le nouveau volume $V' = \SI{1}{L}$. La nouvelle concentration est :
\[ C' = \frac{n}{V'} = \frac{\SI{0,25}{mol}}{\SI{1}{L}} = \SI{0,25}{mol/L}. \]
La solution diluée est \textbf{moins acide} que le détartrant initial : son pH augmente et se rapproche donc de \textbf{7} (il reste toutefois inférieur à 7, car la solution demeure acide).
\end{enumerate}
\textbf{Barème indicatif (10 points) :} question 1 : 2 pts ; question 2 : 2 pts ; question 3 : 2 pts ; question 4 : 2 pts ; question 5 : 2 pts.

\textbf{Points de vigilance :} pour la question 5, citer la conservation de la quantité de matière lors de la dilution ; le pH d'une solution acide diluée se rapproche de 7 mais ne le dépasse jamais ; la réponse à la question 2 doit mentionner le dégagement d'un gaz toxique (le dichlore).
\end{corrige}

\newpage

\coursec{Fiche bilan}

\begin{popfigure}
\begin{tikzpicture}[
  mindmap,
  concept color=popblue,
  level 1 concept/.append style={level distance=4.5cm, sibling angle=60},
  level 2 concept/.append style={level distance=3cm, sibling angle=45},
  every node/.style={font=\small},
  text=white,
  popblue/.style={concept color=popblue},
  poporange/.style={concept color=poporange},
  popgreen/.style={concept color=popgreen},
  poppurple/.style={concept color=poppurple},
  poppink/.style={concept color=poppink},
  popgold/.style={concept color=popgold}
]
\node[concept, align=center]{Solutions\\ aqueuses}
  child[popblue] { node[concept] {Dissolution}
    child { node[concept] {Soluté} }
    child { node[concept] {Solvant} }
    child { node[concept] {Solution} }
  }
  child[poporange] { node[concept] {Solubilité}
    child { node[concept] {Saturée} }
    child { node[concept] {Non saturée} }
    child { node[concept] {Sur-saturée} }
  }
  child[popgreen] { node[concept, align=center]{Concentration\\ molaire}
    child { node[concept] {$C = n/V$} }
    child { node[concept] {Unité : mol/L} }
    child { node[concept] {Préparation} }
  }
  child[poppurple] { node[concept] {pH}
    child { node[concept] {Acide $<7$} }
    child { node[concept] {Neutre $=7$} }
    child { node[concept] {Basique $>7$} }
  }
  child[poppink] { node[concept] {Sécurité}
    child { node[concept] {Étiquettes} }
    child { node[concept] {EPI} }
    child { node[concept] {Déchets} }
  }
  child[popgold] { node[concept] {Applications}
    child { node[concept] {Cuisine} }
    child { node[concept] {Médecine} }
    child { node[concept] {Industrie} }
  };
\end{tikzpicture}
\legende{Carte mentale de la leçon : notions clés, formules, sécurité et applications au Cameroun (eau potable, sérum, batterie, lessive).}
\end{popfigure}

\begin{bilan}
\courssub{Définitions essentielles}
\begin{itemize}
  \item \textbf{Dissolution} : mélange homogène obtenu en introduisant un \cle{soluté} (solide, liquide ou gazeux) dans un \cle{solvant} (souvent l'eau). La solution aqueuse est le résultat.
  \item \textbf{Solubilité} ($s$) : masse maximale de soluté qui peut se dissoudre dans un volume donné de solvant (souvent \SI{1}{L} ou \SI{100}{mL}) à une température fixée. Unité : \SI{}{g/L} ou \SI{}{g/100mL}.
  \item \textbf{Solution saturée} : solution qui contient la quantité maximale de soluté dissous ; tout excédent reste non dissous (dépôt au fond).
  \item \textbf{Concentration molaire} ($C$) : quantité de matière de soluté par litre de solution. $C = \dfrac{n}{V}$.
  \item \textbf{pH} : nombre sans unité qui caractérise l'acidité ou la basicité d'une solution aqueuse. Échelle de 0 à 14.
\end{itemize}

\courssub{Formules à connaître par cœur}
\begin{center}
\begin{tabularx}{\linewidth}{|l|X|X|}\hline
\rowcolor{popblueL}
\textbf{Grandeur} & \textbf{Formule} & \textbf{Unités} \\ \hline
Concentration molaire & $C = \dfrac{n}{V}$ & $C$ en \SI{}{mol/L} ; $n$ en \SI{}{mol} ; $V$ en \SI{}{L} \\ \hline
Quantité de matière & $n = \dfrac{m}{M}$ & $m$ en \SI{}{g} ; $M$ en \SI{}{g/mol} \\ \hline
Masse de soluté & $m = C \times V \times M$ & $m$ en \SI{}{g} \\ \hline
Dilution (conservation $n$) & $C_1 V_1 = C_2 V_2$ & $C$ en \SI{}{mol/L} ; $V$ en \SI{}{L} \\ \hline
pH (notion) & $\text{pH} = -\log[\ce{H3O+}]$ & pH sans unité \\ \hline
\end{tabularx}
\end{center}

\courssub{Méthodes express}
\begin{enumerate}
  \item \textbf{Calculer une concentration} : identifier $m$ (ou $n$) et $V$ ; convertir $V$ en L ; appliquer $C = n/V$.
  \item \textbf{Préparer une solution mère} : peser $m = C \times V \times M$ ; dissoudre dans un bécher ; transvaser dans une fiche jaugée ; compléter au trait de jauge ; homogénéiser.
  \item \textbf{Préparer une solution fille par dilution} : calculer $V_1 = \dfrac{C_2 V_2}{C_1}$ ; prélever $V_1$ à la pipette jaugée ; verser dans fiche jaugée $V_2$ ; compléter au trait de jauge.
  \item \textbf{Caractériser une solution par le pH} : utiliser papier pH ou pH-mètre ; comparer à 7 : $\text{pH}<7$ acide, $\text{pH}=7$ neutre, $\text{pH}>7$ basique.
\end{enumerate}
\end{bilan}

\begin{aretenir}[Checklist avant le BEPC]
\begin{itemize}
  \item $\square$ Je définis correctement soluté, solvant, solution, dissolution.
  \item $\square$ Je distingue dissolution (phénomène physique) et réaction chimique.
  \item $\square$ Je définis la solubilité et je donne son unité (\SI{}{g/L} ou \SI{}{g/100mL}).
  \item $\square$ Je reconnais une solution saturée (dépôt) d'une solution non saturée.
  \item $\square$ J'écris la formule $C = n/V$ avec les bonnes unités ($C$ en \SI{}{mol/L}, $V$ en \SI{}{L}).
  \item $\square$ Je sais calculer $n$, $C$, $V$ ou la masse $m$ à peser ($m = C V M$).
  \item $\square$ Je décris les étapes de préparation d'une solution mère (fiche jaugée) et d'une dilution (pipette + fiche jaugée).
  \item $\square$ Je définis le pH et je classe une solution : acide ($\text{pH}<7$), neutre ($\text{pH}=7$), basique ($\text{pH}>7$).
  \item $\square$ J'interprète une étiquette de danger (pictogrammes corrosif, irritant, danger pour l'environnement).
  \item $\square$ Je cite 2 applications concrètes au Cameroun : eau de javel (désinfection), sérum physiologique (soins), boisson gazeuse ($\ce{CO2}$ dissous), eau de robinet (traitement), batterie (électrolyte).
  \item $\square$ Je respecte la sécurité : porter lunettes/blouse, verser l'acide dans l'eau (jamais l'inverse), étiqueter les fioles, trier les déchets chimiques.
\end{itemize}
\end{aretenir}

\findecours

\end{document}
